% ============================================================
% PAPER 11 of 11 -- forecasting_baselines
% Assembled 2026-09-29 by content-and-merits partition.
% Title: Forecasting Baselines and the Null
% Source: arena agent 1/agent workspace/fam/e1_v60.tex
% Source: arena agent 1/paper rewrites/latex/paperE3_edwards_forecast_ladder_v17.tex
% Source: arena agent 1/agent workspace/fam/ws_v17.tex
% Partition note: MERGED. E1 (cod forecast baselines) + E3 (the Edwards ladder, whose null result leads) + ws (welfare/support).
%% STATUS: structurally assembled as a NEW version. Editorial integration
%%         (de-duplication of shared framework, sibling cross-citations)
%%         is NOT yet done. See PAPERS_MANIFEST.md.
%% ============================================================



% ===== PART: E1_v60  (source: e1_v60.tex) =====
\clearpage

% Pandoc body identical to the wave-9/11 output (asserted at build time), with the
% declaration subsections relocated to a trailing Declarations section.
% Front matter: Amin Abaee, Independent Researcher, clickable ORCID
%(https://orcid.org/0000-0002-0019-1842) and email (amin_abaee@ut.ac.ir)
% beneath the name; date September 6, 2026.
% Back matter: the paper's own declarations, each a titled subsection, plus
% the AI declaration (GLM (Z.ai), Qwen (Alibaba Cloud) and DeepSeek AI
% assisted with drafting and iterative review).
% References carry the Zenodo DOIs.
% Edit the markdown source and re-run the build script to regenerate.
% Compiles error-free with tectonic (also compatible with pdflatex/xelatex).
\documentclass[11pt]{article}
\usepackage[a4paper,margin=1in]{geometry}
\usepackage{amsmath,amssymb}
\usepackage{graphicx}
\usepackage{booktabs,longtable,array,calc}
\usepackage{xcolor}
\usepackage[colorlinks=true,allcolors=blue!45!black]{hyperref}
\usepackage[font=small,labelfont=bf]{caption}
\usepackage{etoolbox}
\AtBeginEnvironment{longtable}{\small}  % dense data tables
\setlength{\tabcolsep}{4pt}
\providecommand{\tightlist}{\setlength{\itemsep}{0pt}\setlength{\parskip}{0pt}}
\setcounter{secnumdepth}{-1}
\emergencystretch=3em
\graphicspath{{../}}
\begin{document}
\title{Does a surplus-production ladder improve forecasts of Northern cod? A scored test on NAFO 2J3KL}
\author{Amin Abaee\\[0.35em]
{\small Independent Researcher}\\[0.55em]
{\small\href{https://orcid.org/0000-0002-0019-1842}{ORCID: 0000-0002-0019-1842}}\\[0.3em]
{\small\href{mailto:amin\_abaee@ut.ac.ir}{amin\_abaee@ut.ac.ir}}}
\date{September 24, 2026}
\maketitle

{\centering\small \textbf{Prepared in the format of Fisheries Research}\par}

\subsection{Highlights}\label{highlights}

\begin{itemize}
\tightlist
\item
  A five-module structural ladder is scored against two naive baselines
\item
  No module attains lower pooled rolling-origin RMSE than persistence
\item
  From a pre-collapse origin, every model misses the 1991--1995 collapse
\item
  Simulation: the rule is specific, and powerful only where signal is strong
\item
  Non-retention is strong evidence for the autonomous module, weak for others
\item
  Outcome years 2016--2021, threshold definition-dependence (29 of 71
  years), and the 2022--2023 survey suspension: the verdict layer, not
  forecast structure, moved
\end{itemize}

\begin{abstract}

Five surplus-production modules and two naive baselines are scored on
Northern cod (\emph{Gadus morhua}), NAFO 2J3KL, under a rule whose
scoring core was fixed beforehand: a module is kept only if it lowers
RMSE against the next-simpler module and persistence. The predictand is
NCAM (Northern Cod Assessment Model) M-shift SSB (DFO, 2016, Table A2),
1983--2015, LRP (limit reference point) 884.6 kt.

No module is retained on these two unpooled series, on rolling RMSE
pooled over origins within each specification and horizon: under the
coarse pass one-year RMSE is 98 kt for persistence against 115--196 kt,
five-year 265 against 289--488 kt. By period the ordering is not uniform. From the fixed pre-collapse origin every model misses the 1991--1995
collapse (694--819 kt against 670 and 688 kt for baselines): constant
productivity with a 1992 catch drop cannot reproduce it. Added structure
is not uniformly harmful, yet none closes the gap. A second, unpooled specification (xteNCAM, 1954--2024, LRP 276 kt) gives
the same outcome (matched persistence 84 kt against M1's 120).

Simulation under known ground truth retains a true autonomous module in
97\% of replicates at collapse-window parameters but has power below
15\% for the stock-flow and depensation alternatives; against truth
outside the tested model class the rule over-retains, so its specificity
is conditional on that class. No Specification A margin against
persistence has an interval excluding zero, though several on
Specification B do. The predictand is retrospectively reconstructed rather than the vintages
at each origin, and catch is supplied along the horizon: a conditional
hindcast, not an operational forecast.

Published outcome data now cover 2016--2021: the outcome window is an
increase regime (340--440 kt on the updated vintage, read within 1--26 kt
by the independent xteNCAM vintage), persistence errs low throughout, and
the window contains no collapse event, so the frozen verdict is neither
confirmed nor reversed by it. On the published files the two reference
points disagree on 29 of 71 years, the 1980s-mean LRP itself moves from
884.6 to 790.0 kt across vintages, and the real 2022--2023 survey
suspension --- two assessment cycles with the leading indicator
uncomputable and removals frozen --- supplies the review-timing boundary
case. The decision-relevant uncertainty of the period accumulated in the
observation and reference-point layer, not in forecast structure.

\end{abstract}

\textbf{Keywords:} northern cod; biomass forecasting; surplus
production; forecast evaluation; stock assessment; prediction skill;
model identifiability

\subsection{1. Introduction}\label{introduction}

Fisheries assessment builds structure by default. State-space
age-structured assessments, surplus-production models, and
ecosystem-linked extensions each add parameters, state variables, or
couplings in the hope of better advice. In the surplus-production implementations evaluated here, that structure
does not pay for itself: across a five-module ladder scored on two
unpooled assessment specifications, no tested implementation attains
lower pooled rolling-origin RMSE than last-value persistence at either
evaluated horizon. Split by period the pooled result is not uniform ---
persistence is not the lowest-error rule during the collapse band ---
and Section 3.5 reports that split. The result is scoped to the estimators and
series tested, and the exercise is a conditional hindcast rather than
an operational forecast test. Whether that structure improves
out-of-sample forecasts --- as opposed to in-sample fit --- is a
separate, testable question, and on it the burden of proof sits with the
added structure.

That burden is now formalised in assessment practice. Comparison against
a naive baseline has become a standard acceptance diagnostic: hindcast
cross-validation scores a model's out-of-sample predictions with the
mean absolute scaled error (MASE) of Hyndman and Koehler (2006), which
divides prediction error by the error of a random-walk forecast, so that
a value above one means the model predicts less accurately than assuming
no change (Kell et al., 2016; Kell et al., 2021). The diagnostic
cookbook of Carvalho et al. (2021) lists prediction skill so measured as
one of four criteria for accepting an integrated assessment, alongside
convergence, fit, and retrospective consistency, and the good-practice
guidance for surplus-production models adopts the same test (Kokkalis
et al., 2024). The reference forecast in that literature is therefore
already the one used here.

Two features of that practice leave the present question open. First,
MASE is almost always computed on the \emph{abundance index} a model
fits, not on the estimated stock state that advice is written about, so
a model can pass on the index it was tuned to while its biomass
trajectory remains untested. Second, the diagnostic is applied to
certify a single accepted assessment rather than to rank a graded
sequence of structural additions against each other and against the
naive rule at once. What is missing is a direct, scored comparison in
which the predictand is the assessment's own biomass output, each
increment of structure has a declared comparator, and the naive rule is
a competitor rather than a scaling constant. That is the comparison
reported here.

Northern cod (\emph{Gadus morhua}) in NAFO divisions 2J3KL is the
canonical demanding case --- and the running instance of this paper. The late-1980s and early-1990s collapse --- a
spawning stock falling from a high, weakly trending phase to a few
percent of its 1980s mean --- occurred under an assessment and
management regime whose failures were dissected in the canonical
post-mortems (Hutchings and Myers, 1994; Walters and Maguire, 1996). The
non-recovery that followed the July 1992 moratorium has kept the
depensation question open for three decades (Shelton and Healey, 1999).
A partial rebuilding was documented in the 2010s (Rose and Rowe, 2015),
but surplus production has since stalled, with some years being negative
(Rose, 2026).

The measured quantity throughout is the condition of the productive
stock --- the spawning-stock biomass whose persistence underwrites
future yield. This is distinct from the catch taken from it. A fishery
can return a high extracted yield while the stock that produces it
declines, and it is the latter that the forecast here must track. Any
model family proposed for this stock must therefore first clear a
minimal bar: it must at least beat the forecast that nothing changes.
One-dimensional autonomous surplus-production maps carry a structural
constraint under sustained constant catch, discussed with its fitted
values and its scope in Section 4. It is a property of a map held at a
fixed removal level, not a reason the forecasts below fail, and the
scored comparison does not isolate it.

This evaluation follows a stated retention rule whose scoring core was
coded before the first scoring pass and applied unchanged (Section 4 records its protocol
status): additional model structure is retained only when it reduces
primary RMSE relative to the next-simpler model and relative to
persistence (\S{}2.3). Early-warning and intervention-selection criteria
are secondary objectives; neither is operationalised here (the
Brier score is reported as a secondary diagnostic and never changes a
retention verdict, and the moratorium is deliberately not evaluated).
The instrument is a forward-ordered ladder of surplus-production models
(not a strict nesting for M2 and M4) evaluated against two naive
baselines. It runs output-only, stock-and-flow, residual, then lagged initialisation;
prey- and index-informed observation variants are evaluated separately
in Section 3.4 and are not part of the ladder. A module is kept only if
it improves the stated score.

The paper asks whether stock-flow, residual, delay, or prey-informed
modules reduce SSB forecast error relative to last-value persistence and
to the autonomous surplus-production model. The scope is deliberately
narrow: this paper does not estimate an Allee threshold, identify the
cause of the 1990s collapse, or evaluate whether the 1992 moratorium was
adequate.

The stock was selected because of its data availability (a full
state-space assessment series), its policy relevance, and the sharpness
of its collapse-and-stall trajectory. The same scored design is applied to a groundwater system, the Edwards
Aquifer, Texas, in Abaee (2026, Does a one-pool water-balance model improve forecasts of Edwards Aquifer head). The two systems' series are not
pooled, and no retention outcome is transferred between them.

The rule that decides retention is itself evaluated. Because a negative
result is only as informative as the instrument that produced it,
Section 3.7 reports a pre-registered simulation of the rule's operating
characteristics under known ground truth: how often it retains a module
that is genuinely present, how often it retains structure when
persistence is the truth, and how it compares with selection by an
information criterion or by scaled error alone. That evaluation is what
licenses the reading of the empirical result module by module rather
than as a single verdict.

The remainder of the article is organized as follows. Section 2 states
the data, the two assessment specifications, the model ladder, and the
evaluation design. Section 3 reports the results: the primary
specification (3.1), annual landings and the survey-start variant (3.2),
the alternative assessment specification (3.3), prey-informed
productivity (3.4), the Diebold--Mariano and block-bootstrap
the Diebold--Mariano and block-bootstrap uncertainty analysis (3.5),
the fitted parameters (3.6), and the simulated operating characteristics
of the retention rule (3.7).
Section 4 discusses the findings, the identification limits, the
reconstruction-level corroboration, and the protocol status.
Section 5 concludes. The contributions are: the scored forecast-ladder
comparison on the full record --- the surplus-production ladder against
naive and assessment-based competitors under the declared evaluation
design (Section 3.1), with the landings, survey-start, alternative
assessment, and prey-informed variants (Sections 3.2--3.4); the
uncertainty analysis of the score differentials --- the
Diebold--Mariano tests and the block-bootstrap, with the fitted
parameters reported (Sections 3.5--3.6); and the retention rule's
simulated operating characteristics with the delayed-revelation reading
of the documented one-year lead (Section 3.7).

\subsection{2. Material and methods}\label{material-and-methods}

\subsubsection{2.1 Data and
specifications}\label{data-and-specifications}

The specification tables type each field as follows: D, data; E,
empirical construct; M, model; N, normative threshold.

\textbf{Table 1.} Primary specification --- Specification A (the
1983--2015 NCAM M-shift SSB series of DFO, 2016, Table A2; the
state-space assessment model of Cadigan, 2016).

\begin{longtable}[]{@{}
  >{\raggedright\arraybackslash}p{(\columnwidth - 4\tabcolsep) * \real{0.3333}}
  >{\raggedright\arraybackslash}p{(\columnwidth - 4\tabcolsep) * \real{0.3333}}
  >{\raggedright\arraybackslash}p{(\columnwidth - 4\tabcolsep) * \real{0.3333}}@{}}
\toprule\noalign{}
Field & Contents & Type \\
\midrule\noalign{}
\endhead
\bottomrule\noalign{}
\endlastfoot
System & Northern cod, NAFO 2J3KL, as represented by NCAM M-shift SSB &
D \\
Interest & Continuity of a spawning stock on the 2016
precautionary-approach cautious or healthy side of the 2016 LRP & N \\
Domain & Stock area in DFO (2016) Figure 1; calendar years 1983--2015 &
D \\
Safe set & \(S_t \ge \mathrm{LRP} = 884.6\) kt (1983--1989 mean of Table
A2) & N \\
Disturbance & Unspecified productivity shocks; not a fitted \(M(t)\) &
M \\
Theoretical catch & Any \(C_t \ge 0\) & M \\
Implementable catch & Pre-1992 directed fishery; post-2 July 1992
moratorium and low inshore removals & E \\
Horizon & Hindcast 1983--2015; two fixed test windows and rolling origin
& D \\
Food web & Capelin excluded from the primary pass; tested as a variant
in Section 3.4 & M \\
Norm & 2010/2016 DFO precautionary-approach LRP (1980s mean SSB), not
the 2023 40\% \(B_{\mathrm{MSY}}\) LRP & N \\
\end{longtable}

Table 1 note: the safe-set and LRP fields type the reference frame of
the evaluation; the primary RMSE score never uses them (only the
secondary Brier indicator score does).

Table A2 also reports fishing mortality \(F\) and natural mortality
\(M\). They are joint outputs with SSB and are not used as exogenous
inputs.

Regular et al.~(2025) extend the assessment to 1954, revise the LRP to
276 kt (95\% interval 180--423 kt; 40\% of \(B_{\mathrm{MSY}}\)), and
estimate 2024 SSB at 342 kt. That is a second specification ---
Specification B (the 1954--2024 extended xteNCAM series; Section 3.3).
The two objects differ in five specification fields: (1) the state
equation (xteNCAM versus NCAM M-shift), (2) the start year (1954 versus
1983), (3) the reference threshold (276 kt versus 884.6 kt), (4) the
treatment of catch, and (5) the horizon.
The two SSB columns are not mixed. The reconstructions are scored separately because they differ in
estimated SSB trajectory, assessment formulation, temporal coverage and
catch specification; the reference threshold enters only the secondary
classification diagnostic and not the primary RMSE comparison. No
retention outcome is carried from one specification to the other.

\subsubsection{2.2 Forecast models}\label{forecast-models}

Northern cod surplus production is represented as a discrete-time map,
with \(S\) in kt and \(C\) in kt yr\textsuperscript{-1}.

\textbf{Definition 2.1 (Surplus-production map).} The autonomous
surplus-production map with multiplicative depensation is

\[
\begin{gathered}
S_{t+1}=\bigl[S_t+g(S_t)-C_t+\varepsilon_t\bigr]_+,
\qquad
g(S_t)=rS_t\bigl(1-S_t/K\bigr)\,a(S_t),
\\
a(S_t)=
\begin{cases}
1 & \text{(Schaefer; no Allee term declared)},\\[2pt]
\dfrac{S_t-\mathfrak s}{K-\mathfrak s} & \text{(depensation term active, $\mathfrak s>0$)}.
\end{cases}
\end{gathered}
\]

Two conventions fix how this map is used. The disturbance
\(\varepsilon_t\) is a regression error at estimation only: parameters
are fitted to the one-step increments with \(\varepsilon_t\) as the
residual, and all forecasts reported here set \(\varepsilon_t = 0\), so
every forecast trajectory is a deterministic plug-in path rather than a
multi-step conditional mean, except in M3 and M4, which carry the
fitted residual forward as a state. In those two modules the training residuals are
\(e_u = \Delta S_u - [g(S_u) - C_u]\) over the training transitions,
formed before any clipping. The coefficient \(\hat\phi\) is the
no-intercept lag-one least-squares estimate
\(\sum_u e_u e_{u-1}/\sum_u e_{u-1}^2\); it is set to zero when the
denominator is not positive or the window supplies fewer than four
residuals, and is clipped to \([-0.95, 0.95]\). The residual carried
into the forecast is the last training residual \(e_{\mathrm{last}}\),
and the projected residual at step \(k \ge 1\) is
\(\hat\phi^{\,k} e_{\mathrm{last}}\), added inside the update before
the state is clipped to \([10^{-3}, 10^{6}]\) kt.
Catch is dated so that \(C_t\) is the catch taken during the interval
from \(t\) to \(t+1\); a forecast issued at origin \(t\) therefore
uses the recorded catch for the year it is forecasting. Under the
coarse regime this places the 1992 drop in the step from 1992 to 1993,
so on the fixed collapse test window (1991--1995) the drop does not
enter the first two forecast steps; the models are already rising before
it applies.
For multi-step forecasts the catch path is supplied for the
whole horizon, so the exercise is a conditional hindcast with respect to
catch, not an operational forecast; Section 4 returns to this.

The factor \(a(S_t)\) is a multiplicative depensation term, not a
parameter switch. Setting \(\mathfrak s=0\) gives \(a(S_t)=S_t/K\) (a
cubic modification of the logistic surplus; Pella and
Tomlinson, 1969), not the Schaefer (1954) law. The
Schaefer model is the separate member of the family in which the factor
is replaced by \(1\); Abaee (2026, Periodic review as sampled governance) makes the same distinction.

\medskip
\noindent\textbf{Shared symbols across the companion papers.} The
companion papers share this vocabulary with paper-local meanings; the
letters below carry meanings here that differ there:
\begin{center}\small
\begin{tabular}{@{}p{0.40\linewidth}p{0.52\linewidth}@{}}
\toprule
\textbf{Here} & \textbf{In the companions} \\
\midrule
\(K\): the carrying capacity &
the safe set and kernel \(K\) (obstruction calculus); the
full-information kernel (certification); the window depth and contact
set \(\mathcal{K}(B)\) (minimax duals) \\
\(r\): the intrinsic growth rate &
the catch share \(r = C/B_t\) (applied regime viability) \\
\(\sigma\): the process/noise scale &
the kernel time variable \(\sigma\) (certification; minimax duals);
the floored exit count \(\sigma^{*}\) (worked systems) \\
\bottomrule
\end{tabular}
\end{center}

Section SI-2 gives the statement and proof.

The model ladder is given in Table 2.

The ladder is the forward-ordered set of seven entries \(\{\)persist,
mean, M1, M1b, M2, M3, M4\(\}\) of Table 2, ordered by structural
complexity (autonomous surplus map, Allee extension, stock-flow,
residual, delay): two naive baselines (persist, mean) and five
structural modules. Where this article refers to a five-module ladder
evaluated against two naive baselines, it means this set. It is not a
strict nesting for M2 and M4.

\textbf{Table 2.} Model ladder.

\begin{longtable}[]{@{}
  >{\raggedright\arraybackslash}p{(\columnwidth - 6\tabcolsep) * \real{0.2500}}
  >{\raggedright\arraybackslash}p{(\columnwidth - 6\tabcolsep) * \real{0.2500}}
  >{\raggedright\arraybackslash}p{(\columnwidth - 6\tabcolsep) * \real{0.2500}}
  >{\raggedright\arraybackslash}p{(\columnwidth - 6\tabcolsep) * \real{0.2500}}@{}}
\toprule\noalign{}
ID & Class & Free on the training window & Frozen into the test window \\
\midrule\noalign{}
\endhead
\bottomrule\noalign{}
\endlastfoot
persist & baseline & --- & \(\hat S_{t+h}=S_t\) \\
mean & baseline & training mean &
\(\hat S_{t+h}=\bar S_{\mathrm{train}}\) \\
M1 & autonomous & \(r,K\); \(C\) = training-mean catch (plugged, not
estimated) & same \(C\) \\
M1b & autonomous with Allee & \(r,K,\mathfrak s\); \(C\) =
training-mean catch (plugged, not estimated) & same \(C\) \\
M2 & stock-flow & \(r,K\); \(C_t\) prescribed & prescribed \(C_t\) \\
M3 & residual & M2 + AR(1) residual & \(\phi\) persisted \\
M4 & lagged initialisation & \(r,K\) and \(\phi\) refitted exactly as for M3 & forecast
starts from \(S_{t-1}\) \\
\end{longtable}

The coarse catch regime, taken from DFO (2016) prose (the
precautionary-approach framework of DFO, 2009), is \(C_t=240\) kt for
\(t\le 1991\), \(C_t=120\) kt for \(t=1992\), and \(C_t=5\) kt for
\(t\ge 1993\). Year-by-year landings (Schijns et al., 2021, Table 1)
replace that regime in Section 3.2. Regular et al.~(2025) Table 1
matches Schijns exactly on 1983--1993 (11 years, maximum absolute
difference 0 t); STATLANT matches Schijns on 1983--1993, so the
STATLANT-versus-Schijns sensitivity is closed on 1983--1993 (they are
the same column there). The Specification A collapse test window runs to
1995, two years beyond the verified match; the residual exposure is
immaterial, since Schijns records 1.31 kt in 1994 and 0.41 kt in 1995
against a collapse-window \(\Delta S\) of hundreds of kilotonnes. A
1956 discrepancy between those sources (236,210 t versus 263,210 t) lies
outside the 1983--2015 Specification A range; Specification B begins in
1954 and so contains 1956, but it draws its catch column from the
xteNCAM reconstruction rather than from either of these two sources, so
the discrepancy is unused on both specifications. The coarseness of the regime series is a limitation.

Parameters are estimated by one-step least squares on the training
window only. Rolling-origin training windows expand rather than slide:
each origin refits on all history up to that origin, subject to the
minimum length. A window of \(m\) annual states supplies
\(m-1\) one-step transitions, and it is the transition count that sets
the effective regression sample. Bounds: \(r\in(0.001,2]\), and \(K\) constrained above the largest state
the one-step regression conditions on. The estimation code optimises
\(K\) on \([\max_{\mathrm{train}} S + 10, 5000]\) kt, where
\(\max_{\mathrm{train}} S\) is taken over the predictor states \(S_t\)
of the training transitions and so excludes the terminal state of the
window. Here \(500\) kt is the multi-start initialiser rather than the
lower bound, consistent with the frozen specification's constraint of
\(K\) above the training maximum.
The remaining estimation bounds are
\(\mathfrak s \in [0, \max_{\mathrm{train}} S_t]\) with the further
feasibility restriction \(0 < \mathfrak s < 0.8K\) enforced in the
objective, \(\phi\) clipped to \([-0.95, 0.95]\), and the index
exponent \(b\) fitted jointly with \(r\) and \(K\) without bounds.
The profile of Section 3.2 restricts \(\mathfrak s\) further, to
\([0, \min_{\mathrm{train}} S_t]\), for the reason given there. The
reported fits attain the upper endpoint (\(K = 5000\) kt) where the
objective is flat out to the imposed bound; M1b's recovery-window
\(K = 105.8\) kt is a valid interior fit, not a bound violation. Where a
reported fit sits at \(500.0\) kt this is the multi-start initialiser,
which the optimiser has not left because the objective is flat there,
and not an active constraint. M3's autoregressive coefficient is
estimated in a second stage: \(r\) and \(K\) are fitted first by
one-step least squares, and \(\phi\) is then obtained as the
lag-one regression coefficient of the resulting training residuals on
their own lag, clipped to \([-0.95, 0.95]\); it is not jointly
optimised with \(r\) and \(K\). Every
fitted-parameter value this article prints, with the window it belongs
to and the bound it attains, is collected in Table 10; the per-origin
rolling fits and the parameters the article does not print are archive
items, marked there as such.

A survey-start variant replaces the surplus initial condition by
\(\hat q\,I_t\), where \(I_t\) is the autumn research-vessel abundance
index and \(\hat q\) is the training-window median of
\(\mathrm{SSB}/I\). The index is an NCAM input, not an independent
stock.

Prey-informed variants (Section 3.4) scale surplus production by a 1991
capelin regime or by the tabulated 3L spring acoustic index (DFO,
2024b). Pre-1991 acoustic values are not carried across 1991. Missing
survey years are not interpolated; within the post-1991 regime the last
eligible observation is carried forward, and training transitions with
no eligible index value are dropped from the fit.

\subsubsection{2.3 Evaluation design}\label{evaluation-design}

\textbf{Table 2b.} Information set at forecast origin \(t\). ``Available''
means the quantity is dated at or before \(t\) and is used by the module
named; ``supplied'' means the quantity is dated after \(t\) and is
nevertheless provided to the forecast, which is what makes the exercise a
conditional hindcast rather than an operational test.

\begin{longtable}[]{@{}
  >{\raggedright\arraybackslash}p{(\columnwidth - 4\tabcolsep) * \real{0.34}}
  >{\raggedright\arraybackslash}p{(\columnwidth - 4\tabcolsep) * \real{0.30}}
  >{\raggedright\arraybackslash}p{(\columnwidth - 4\tabcolsep) * \real{0.36}}@{}}
\toprule\noalign{}
Quantity & Status at origin \(t\) & Used by \\
\midrule\noalign{}
\endhead
\bottomrule\noalign{}
\endlastfoot
Target series \(S_u,\ u\le t\) & available; single retrospective vintage, not the vintage current at \(t\) & all modules and baselines \\
Origin state \(S_t\) & available & all except M4 \\
Lagged state \(S_{t-1}\) & available & M4 (forecast start) \\
Training transitions \(u\le t\) & available; window expands with \(t\) & parameter estimation, all structural modules \\
Fitted residual \(e_{\mathrm{last}}\) & available; last training transition & M3, M4 \\
Catch \(C_u,\ u\le t\) & available & M1, M1b (training mean); M2--M4 \\
Catch \(C_u,\ u>t\) & \textbf{supplied} along the horizon & M2, M3, M4 \\
Prey index \(I_u,\ u\le t\) & available; last observation carried forward & M\_cap\_index \\
Prey index \(I_u,\ u>t\) & not used & --- \\
Reference point (LRP) & fixed by specification & secondary Brier score only \\
\end{longtable}

The distinction the table draws --- between quantities dated at or
before the origin and quantities dated after it but supplied anyway ---
is not specific to this study, and reporting it explicitly is proposed
here as a template for hindcast evaluations generally: an evaluation
that does not state which of its inputs were unavailable in real time
cannot be read as evidence about operational skill. The single row that
makes this a conditional hindcast is the supplied future catch: M2, M3 and M4 receive the realised catch path over the
forecast horizon, which no operational forecast would have. That asymmetry
strengthens rather than weakens the comparison: the three modules given
information no operational forecast could possess still do not attain
lower pooled rolling RMSE than a rule that uses only the current state.
Persistence receives no catch information at all. M1 and M1b
do not, using instead the training-mean catch, which is why their scores
are nearly invariant to the catch treatment. The target series is a
single retrospective reconstruction throughout, so no module sees the
assessment vintage that would have been current at its origin.

\textbf{Table 2c.} Rolling-origin sets, as constructed by the scored
code and archived per experiment. Minimum training length is eight years
for the naive baselines throughout, eight for the Specification A
ladder, and twelve for the Specification B ladder; the baseline and
ladder origin sets on Specification B therefore differ, and every
comparison reported against a baseline is recomputed on the module's own
origins.

\begin{longtable}[]{@{}llrl@{}}
\toprule\noalign{}
Experiment & \(h\) & \(n\) & Origin range \\
\midrule\noalign{}
\endhead
\bottomrule\noalign{}
\endlastfoot
Specification A ladder & 1 & 25 & 1990--2014 \\
Specification A ladder & 5 & 21 & 1990--2010 \\
Specification B ladder & 1 & 59 & 1965--2023 \\
Specification B ladder & 5 & 55 & 1965--2019 \\
Specification B baselines & 1 & 63 & 1961--2023 \\
Specification B baselines & 5 & 59 & 1961--2019 \\
M\_cap\_index, Specification A & 1 & 24 & 1991--2014 \\
M\_cap\_index, Specification A & 5 & 20 & 1991--2010 \\
M\_cap\_index, Specification B & 1 & 36 & 1988--2023 \\
M\_cap\_index, Specification B & 5 & 32 & 1988--2019 \\
M\_cap\_regime, Specification A & 1 & 25 & 1990--2014 \\
M\_cap\_regime, Specification B & 1 & 59 & 1965--2023 \\
\end{longtable}

The index experiments start later than the ladder because an origin is
scored only where a carried-forward index value exists. On Specification
B that admits three pre-break origins (1988--1990), which draw on
pre-1991 surveys; the carry is reset at 1991 so that no pre-break value
is carried into the post-break regime.


What the design does and does not test should be stated before the
scores. The predictand at every origin is a single retrospectively
reconstructed assessment series, not the assessment vintage that would
have been available at that origin; assessment vintages were not available for this
evaluation. Earlier SSB values therefore embed information produced
after the nominal forecast date, and the smoothing inherent in an
assessment reconstruction is present in both the target and the
persistence baseline. Catch is supplied along the forecast horizon from
the catch file rather than forecast. The exercise is accordingly
a retrospective prediction experiment and, with respect to catch and the
prey index, a conditional hindcast; it is not a test of operational
forecasting skill, and the comparison against persistence inherits the
same status. Section 4 returns to what this does and does not license.

Fixed windows on Specification A: collapse, train 1983--1990, test
1991--1995; recovery, train 1995--2007, test 2008--2015. The year 1995
appears in the collapse test window and again as the first recovery
training year. The two fixed-window exercises are separate fits scored
separately and neither informs the other, so this is not leakage within
an exercise; it is noted because the two windows are reported side by
side. Rolling origin:
minimum eight training years; horizons \(h=1\) and \(h=5\). Estimation
is by one-step least squares on the training window, while the five-year
score evaluates iterated trajectories of parameters fitted for one-step
fit --- a training/testing cost-function mismatch;
trajectory-matched estimation over the multi-year horizon is an
untested alternative, not scored here.

The primary score is RMSE of SSB (kt), taken over all scored
origin--target pairs at the stated horizon; the \(h=5\) score is the
root mean square of the five-year-ahead errors themselves, not an
average along the intervening trajectory. Log-RMSE uses natural
logarithms of states floored at \(10^{-3}\) kt. The sign-hit rate
compares \(\mathrm{sign}(\Delta S)\) between forecast and observation
over adjacent years within a window, giving one fewer comparison than
there are scored years, and is reported as NA where a rule supplies no
directional prediction rather than as a zero hit rate. Secondary scores are mean
absolute error, RMSE on \(\log S\), the threshold misclassification
rate for \(\mathbf{1}\{\hat S<\mathrm{LRP}\}\) --- equivalently the
Brier score for a deterministic binary forecast, which is what a hard
\(0/1\) prediction makes it, so no probability forecast is implied ---
and the sign-hit rate of \(\Delta S\) on fixed windows. On Specification A the Brier indicator
rests on a degenerate event --- the \(884.6\)-kt LRP sits above the
entire origin range, so the indicator is \(1\) at every origin and
persistence attains \(0.00\) by construction. The score does separate
models in this sample (persistence and the training mean at \(0.00\)
against \(0.80\)--\(1.00\) for the structural modules, Table 3), but it
separates them on an outcome that never varies, so it carries no
early-warning information and is reported as a secondary diagnostic only
(Section 3.3 gives the record). The
sign-hit rate is conventionally \(0.00\) for persistence, whose forecast
\(\Delta S = 0\) has no sign; persistence is excluded from that ranking
by declaration.

A descriptive scale comparison is available because the extended
assessment publishes an uncertainty band on its own predictand:
\texttt{xtencam\_table17\_ssb.csv} carries 95\% limits for every year,
with a median relative width of \(0.44\). Writing
\(D = |\hat S_{t+h|t} - S_{t+h}| / [(\mathrm{hi}_{t+h} -
\mathrm{lo}_{t+h})/2]\), the proportion of hindcast errors smaller than
half the published interval width is, on Specification B, \(56\%\) for
persistence, \(36\%\) for M1 and \(59\%\) for M3 at \(h=1\)
(\(n=59\)), and \(18\%\), \(5\%\) and \(7\%\) at \(h=5\) (\(n=55\)).

These are \textbf{not} coverage probabilities and are not comparable to a
nominal \(95\%\). The published interval expresses uncertainty about
true biomass given the assessment's data and model; re-centring its
width on a hindcast produces an interval with no established coverage,
and the bounds are in any case right-skewed on the biomass scale
(\(67\) of \(71\) years asymmetric by more than \(10\%\), close to
symmetric only in logarithms), so a symmetric normal reading of the
width would be the wrong transform. Nor does the statistic establish
that model \emph{rankings} are robust to assessment uncertainty: it
compares each model's errors with a target-state interval, not the
differences between models. It indicates only that hindcast
discrepancies are mostly of the same order as, or larger than, reported
target-state uncertainty. On this statistic M3 (\(59\%\)) reads
marginally above persistence (\(56\%\)) at \(h=1\), which the
rolling-origin RMSE ordering does not; the two need not agree, and no
retention verdict rests on either this statistic or that reading.

The retention rule is stated as explicit hypotheses.

\textbf{Retention rule.} \emph{A module \(M\) on the ladder of Table 2
is retained only if all of the following hold on the rolling-origin
primary RMSE score:}

\begin{itemize}
\tightlist
\item
  \emph{(H1) \(M\) reduces primary RMSE relative to its declared
  comparator --- the next-simpler rung by complexity ordering (M1b against
  M1, M3 against M2). Only M3 against M2 is a strict nesting: M1b is a
  separate branch of Definition 2.1 rather than a parameter restriction
  of M1, since \(\mathfrak s \to 0\) yields the cubic factor
  \(a(S)=S/K\) and not the Schaefer law (Section SI-2), so its
  comparator is declared by ordering rather than by nesting. For the two
  rungs that are not nestings at all, M1
  for M2 (the autonomous constant-catch map whose catch treatment M2
  changes) and M3 for M4 (the module the lagged initialisation acts on). M1b is reported
  as the alternative comparator for M2 and never decides retention. M1
  is the first structural rung and has no simpler structural comparator,
  so H1 is vacuous for M1 and its retention turns on H2 and H3;}
\item
  \emph{(H2) \(M\) reduces primary RMSE relative to last-value
  persistence;}
\item
  \emph{(H3) each required reduction holds at both horizons \(h=1\) and
  \(h=5\) --- the frozen specification's primary score is the
  rolling-origin RMSE pair --- and exceeds a tie band of 5\% of the
  comparator's score; improvements inside the band are ties and do not
  retain.}
\end{itemize}

\begin{verbatim}
Retention rule (as scored; band = 0.05, horizons = {1, 5})

  for each module M on the ladder:
      retain <- TRUE
      for h in horizons:
          if RMSE(M, h) >= (1 - band) * RMSE(persistence, h):
              retain <- FALSE                    # H2: beat persistence
          if comparator(M) exists and
             RMSE(M, h) >= (1 - band) * RMSE(comparator(M), h):
              retain <- FALSE                    # H1: beat next-simpler
      report retain                              # H3: both horizons, band
\end{verbatim}

Scores are origin-matched: every comparison uses the module's own origin
set, and the baseline is recomputed on that set rather than taken from a
longer one. The rule is deliberately conservative in two places, and
Section 3.7 quantifies what each costs: the band suppresses retention on
margins inside sampling noise, and the comparator requirement forbids a
module from being retained merely because it beats persistence when a
simpler member of the same ladder does so too.

\emph{Retention is decided separately on each specification. A module
failing any of (H1)--(H3) is not retained. Two disclosures complete the
rule. First, it is a point-rule ranking under a pre-set score, not a
test of equal predictive ability; the uncertainty analysis of
Section 3.5 attaches Diebold--Mariano and moving-block-bootstrap
intervals to its margins and changes no verdict. Second, the tie band
and the comparator declarations are completions recorded after the
scores of Section 3 were computed: no recorded verdict depends on them
--- the smallest structural deficit against persistence (M1b's on
Specification A at \(h=1\); \((114.80-98.05)/98.05 = 17.1\%\)) is far
outside the \(5\%\) band, and no (H1) comparison in Tables
3--8 reverses under either comparator reading.}

A module is not retained when the one-step least-squares implementation
evaluated here fails at least one of (H1)--(H3) on the rolling-origin
primary RMSE score on the series tested. Non-retention is a statement
about these implementations on these series, not about every possible
implementation of the underlying model class, and it is weaker than a
statistical null result: it records that a module did not clear a stated
bar, not that no difference exists.

The scoring is deterministic: the scripts are listed in the Data
availability statement with pinned output checksums and an archived
environment specification, so the reported arithmetic regenerates byte
for byte within that environment. Cross-environment numerical
identity is not claimed, and one M1b row is environment-sensitive at
the \(\pm 17\)-kt level, as recorded there.
Reproducibility of the arithmetic is distinct from the analytical result
of Section 4, which is proved rather than computed.

On Specification B the collapse window is train 1954--1989, test
1990--1995, and the recovery-stall window is train 1995--2012, test
2013--2024. As on Specification A, 1995 appears both in the collapse
test window and at the start of the later training window; the two
windows are scored independently and no fitted quantity crosses between
them, so the shared year is not leakage. Catch is Regular et al. (2025) Table 1 landings (2024 persisted from
2023). No row
is taken from NCAM 2016. The ladder's rolling origin on Specification B
uses a minimum of twelve training years (\(n=59\) at \(h=1\), \(n=55\)
at \(h=5\)); the naive baselines retain the eight-year minimum (\(n=63\)
and \(n=59\)). The two origin sets therefore overlap but are not identical, so a
baseline computed over the longer naive origin sets would mix four extra
early origins into the comparison. Table 6 accordingly reports the
origin-matched baselines, computed on the ladder's own origins: on the
identical twelve-year origins persistence reads 84 kt (\(n=59\)) at
\(h=1\) and 300 kt (\(n=55\)) at \(h=5\), and the training mean 458 kt
and 522 kt, against M1's 120 kt and 432 kt. Persistence remains the
lowest-RMSE forecast. The mixed-origin reading (88 kt, quoted in the
Table 6 caption) exceeds the controlled reading of 84.4 kt by 3.6 kt ---
an origin-mix effect of no retention consequence. On the identical post-break origins of the Table 7 control
(\(n=33\), \(h=1\)) the recomputed baseline is 75 kt, matching the
post-break value reported with the two-regime control of Section 3.4. On
the index modules' origin sets (Table 8), persistence reads 97 kt
(\(n=24\), \(h=1\)) and 193 kt (\(n=20\), \(h=5\)) on Specification A,
and 79 kt (\(n=36\)) and 288 kt (\(n=32\)) on Specification B. The index
module loses to the origin-matched baseline in every cell, and the
verdict --- the module is not retained --- is unchanged.

\subsection{3. Results}\label{results}

\subsubsection{3.1 Primary specification}\label{primary-specification}

\begin{figure}[htbp]
\centering
\includegraphics[width=\linewidth]{figs_e1/fig1_series.png}
\caption{NCAM M-shift SSB (DFO, 2016, Table A2). Dashed line: 2016 LRP = 884.6 kt. 2015 SSB is 33.8\% of that LRP, matching the advisory report statement of 34\%.}
\end{figure}

\begin{figure}[htbp]
\centering
\includegraphics[width=\linewidth]{figs_e1/fig2_windows.png}
\caption{Multi-step forecasts from the end of each training window. Recovery is under-predicted when an AR residual fitted on a slow early-recovery train is persisted.}
\end{figure}

\textbf{Table 3.} Fixed-window scores (RMSE in kt), with both baselines
reported on every window, computed on the same fixed origins set
in Section 2.3. Table 3 is not to be read as a retention table ---
retention is decided on rolling-origin primary RMSE only.

\begin{longtable}[]{@{}llrrrrr@{}}
\toprule\noalign{}
Window & Model & RMSE & MAE & log-RMSE & Brier & Direction \\
\midrule\noalign{}
\endhead
\bottomrule\noalign{}
\endlastfoot
Collapse & persist & 670 & 610 & 2.71 & 0.00 & NA \\
& M1 & 694 & 638 & 2.73 & 1.00 & 0.50 \\
& M1b & 694 & 636 & 2.73 & 0.80 & 0.00 \\
& M2 & 819 & 751 & 2.85 & 1.00 & 0.25 \\
& M3 & 819 & 751 & 2.85 & 1.00 & 0.25 \\
& M4 & 819 & 750 & 2.85 & 1.00 & 0.25 \\
& mean & 688 & 630 & 2.72 & 0.00 & NA \\
Recovery & persist & 104 & 72 & 0.69 & 0.00 & NA \\
& M1 = M2 & 120 & 105 & 0.61 & 0.00 & 0.57 \\
& M1b & \textbf{90} & 55 & 0.52 & 0.00 & 0.57 \\
& M3 & 220 & 200 & 0.92 & 0.00 & 0.57 \\
& M4 & 214 & 195 & 0.91 & 0.00 & 0.57 \\
& mean & 144 & 124 & 1.60 & 0.00 & NA \\
\end{longtable}

On collapse, fitted \(r\) reaches \(1.935\), close to but not at the
imposed upper bound of \(2\).
The 1983--1990 window is a high, weakly trending stock. The fitted pre-collapse model does not reproduce the subsequent
decline, and dropping \(C\) from 240 to 5 raises forecast SSB. In the model the 1992 drop is supplied as an
exogenous catch path; historically it was largely a policy response to
depletion. Because production is stationary, prescribing the lower catch
in a constant-\(r\) stock-flow equation raises the predicted increment
and forces a mechanical rebound instead of reproducing the crash. Within this
accounting, had the crash been driven by the catch regime one would
expect supplying the catch path to help; M2 does not improve on M1. Supplying the prescribed catch path did not improve this
implementation's collapse-window hindcast. That is not proof of one
alternative: a misspecified and weakly identified module need not
improve even when the variable it adds is genuinely relevant, so the
result bounds what this ladder can attribute rather than identifying the
mechanism.

Under the coarse catch regime, M1 and M2 coincide on the recovery
window; under annual landings they do not (264 against 303 kt on the fixed recovery window, Section 3.2). The reason is structural:

On the recovery window (train 1995--2007, test 2008--2015) of
Specification A, M1 and M2 produce identical forecasts: M1 plugs the
training-mean catch into the autonomous map and M2 prescribes \(C_t\)
year by year, but \(C_t\equiv 5\) kt on both train and test windows, so
the two prescriptions coincide and the least-squares \((r,K)\) optimum is
identical.

M1b reports a lower RMSE (90 kt against M1's 120 kt), but
\(\mathfrak s\) descends towards zero and \(K\) settles just above the
training predictor range at \(105.8\) kt: an
unidentified Allee parameter, not a biological threshold. The estimate
approaches the lower boundary to numerical tolerance, for a reason visible in the objective.

\textbf{Observation 3.2 (Boundary Allee estimate on the recovery
window).} Any \(\mathfrak s\) strictly above the training minimum makes
the depensation factor
\(a(S_t)=(S_t-\mathfrak s)/(K-\mathfrak s)\) negative for observations
below \(\mathfrak s\), so wherever the observed \(\Delta S\) is positive
at such a state the fit incurs a large squared residual. On the
1995--2007 recovery window seven of the twelve annual increments are
positive --- the five negative ones fall in a single consecutive run,
1999--2004 --- and on balance this penalty pushes the estimate downward: the fitted \(\mathfrak s\)
descends towards zero (\(2.1\times10^{-23}\) under the coarse regime and
\(9.4\times10^{-6}\) under annual landings) while \(r\) reaches its
upper bound of \(2\) and \(K\) settles just above the training range at
\(105.8\) kt. Zero is excluded by the objective's feasibility
restriction \(0 < \mathfrak s < 0.8K\), so these values are a numerical
approach to an unattained infimum rather than an attained boundary
estimate. At such values the fitted object is effectively the
zero-threshold cubic branch \(a(S)=S/K\) of Definition 2.1, not a
positive-threshold Allee model, so M1b's lower error here is not
evidence for depensation: \(r\) and \(K\) absorb what \(\mathfrak s\)
does not.

This is an empirical finding about these fits, not a theorem. Two qualifications matter. The pressure just
described is a tendency and not a proof that the global optimum must lie
below the training minimum: the objective trades errors across all
observations, and \(r\), \(K\) and \(\mathfrak s\) can compensate (for M1b the catch is
plugged, not estimated, and the profile re-optimises only \(r\) and
\(K\))
for one another. The window is also not monotone --- SSB declines in
five consecutive annual decreases, from \(34.59\) kt in 1999 to
\(20.07\) kt in 2004, before rising to \(81.10\) kt in 2007 --- so no argument that
assumes a monotone recovery applies to it. Nor is a boundary estimate by
itself proof of non-identifiability, and a biologically meaningful
threshold may lie outside the sampled range.

A profile of the least-squares objective shows how weakly the data
constrain \(\mathfrak s\). The grid for \(\mathfrak s\) spans the profile range
\([0, \min_{\mathrm{train}} S_t] = [0, 9.68]\) kt. Two different ranges are in play and should not be conflated. Both are
built from the \emph{predictor} states of the training window --- the
\(S_t\) that the one-step regression conditions on, which exclude the
terminal response --- so on the recovery window (states 1995--2007) the
relevant maximum is \(S^{\mathrm{pred}}_{\max} = 40.83\) kt over the
twelve transitions beginning 1995--2006, not the \(81.10\) kt terminal
state of 2007.

\begin{longtable}[]{@{}lll@{}}
\toprule\noalign{}
Quantity & Range & Recovery window \\
\midrule\noalign{}
\endhead
\bottomrule\noalign{}
\endlastfoot
\(K\) & \([S^{\mathrm{pred}}_{\max}+10,\ 5000]\) kt & \([50.83,\ 5000]\) kt \\
\(\mathfrak s\), scored fit & \((0,\ S^{\mathrm{pred}}_{\max}]\), open at zero & \((0,\ 40.83]\) kt \\
\(\mathfrak s\), profile & \([0,\ \min_{\mathrm{train}} S_t]\) & \([0,\ 9.68]\) kt \\
\end{longtable}

The profile is computed on the narrower range so that the production
multiplier is non-negative at every predictor state. That is an
interpretive restriction adopted for the profile, not a feasibility
condition of the model and not a statistical ground for excluding higher
thresholds: a negative multiplier can be compatible with negative
observed increments, and the objective trades errors across
observations. Both ranges therefore exclude
a threshold placed above the collapsed range: at \(\mathfrak s = 50\) kt
--- outside the scored bound --- the factor \(a(S)\) would be negative
at all twelve predictor states, so the model would predict negative
production across the whole window. The estimate at the lower edge is
consequently a statement about a restricted feasible set, and this
design neither supports nor excludes depensation hypotheses --- a
predator pit, for instance --- that place the stock below a high
threshold throughout the recovery. Testing one requires lifting the
range bound, which is outside the frozen specification scored here. Fixing
\(\mathfrak s\) on that grid and re-optimising \(r\) and \(K\) at each
value gives a profile objective that rises only from \(126.50\) to
\(135.04\) kt\textsuperscript{2} under the coarse regime, and from
\(122.90\) to \(126.50\) kt\textsuperscript{2} under annual landings. Expressed as
\(n_{\mathrm{tr}}\log\{\mathrm{SSE}(\mathfrak s)/
\mathrm{SSE}(\hat{\mathfrak s})\}\) with \(n_{\mathrm{tr}}=12\), the
largest deviation across the grid is \(0.78\) under the coarse regime
and \(0.35\) under annual landings. These are descriptive measures of curvature, not calibrated
likelihood-ratio tests:
the least-squares fit carries no error model, the estimate
\(\hat{\mathfrak s}\) lies on the boundary of its range, \(r\) rests
at its own upper bound, the window supplies twelve transitions, and the
states are assessment reconstructions rather than direct observations.
Any one of these would compromise a \(\chi^2_1\) reference. The
substantive reading does not depend on the reference distribution: the
objective is nearly flat across the whole range, so the data do not
locate \(\mathfrak s\) within it.

The mechanism is parameter compensation. In
\(a(S_t)=(S_t-\mathfrak s)/(K-\mathfrak s)\), raising \(\mathfrak s\)
reduces the numerator while lowering \(K\) raises \(1/(K-\mathfrak s)\),
and along the profile \(K\) slides from \(106.0\) to \(82.7\) kt
(coarse) or \(129.8\) to \(95.6\) kt (annual landings) while \(r\)
remains at its upper bound of \(2\) throughout. The product
\(rS_t(1-S_t/K)\,a(S_t)\) is nearly invariant along this path, so the
data cannot separate \(\mathfrak s\) from \(K\). The training
observations therefore do not constrain behaviour near a putative
threshold, and the boundary estimate is not evidence about a biological
Allee threshold in either direction.

M3's \(\phi=0.95\) persists a negative residual, which raises its error
relative to persistence on this window; against its own declared
comparator M2 it still reads lower.

\begin{figure}[htbp]
\centering
\includegraphics[width=\linewidth]{figs_e1/fig3_rmse.png}
\caption{Rolling-origin RMSE. Persistence is the best one-year and five-year point forecast.}
\end{figure}

\textbf{Table 4.} Rolling-origin summary, Specification A, coarse catch
regime.

\begin{longtable}[]{@{}lrrrr@{}}
\toprule\noalign{}
Model & \(h=1\) RMSE & \(h=1\) MAE & \(h=1\) log-RMSE & \(h=5\) RMSE \\
\midrule\noalign{}
\endhead
\bottomrule\noalign{}
\endlastfoot
persist & \textbf{98} & \textbf{48} & \textbf{0.52} & \textbf{265} \\
M1 & 121 & 80 & 8.02 & 289 \\
M1b & 115 & 80 & 8.70 & 289 \\
M2 & 144 & 61 & 0.59 & 398 \\
M3 & 135 & 53 & 3.39 & 366 \\
M4 & 196 & 82 & 0.76 & 488 \\
mean & 424 & 375 & 2.35 & 507 \\
\end{longtable}

Log-RMSE for M1 and M1b is large because some origins produce
trajectories that hit the numerical floor. M2 keeps the state off zero
and stabilizes the log score, but loses on raw RMSE. M4, the only model started from a year-old state, is the worst
structural model on raw RMSE; its decomposition against a stale-start
persistence control is reported in Section 4.

No structural model --- M1 through M4 --- is retained on Specification
A. Persistence is the lowest-RMSE forecast.

\subsubsection{3.2 Annual landings and survey
start}\label{annual-landings-and-survey-start}

Year-by-year landings for 2015 are 4.436 kt, matching DFO reported
landings. Pre-collapse catches are 172--269 kt, not a flat 240. The 1992
drop is to 41 kt, then 11 kt, then 0.4--1.9 kt.

\textbf{Table 5.} Rolling-origin RMSE (kt) with annual landings.

\begin{longtable}[]{@{}lrr@{}}
\toprule\noalign{}
Model & \(h=1\) & \(h=5\) \\
\midrule\noalign{}
\endhead
\bottomrule\noalign{}
\endlastfoot
persist & \textbf{98} & \textbf{265} \\
M1 & 121 & 289 \\
M1b & 115 & 289 \\
M2 & 160 & 394 \\
M3 & 154 & 352 \\
M4 & 206 & 486 \\
M2, survey start & 128 & 331 \\
\end{longtable}

Annual landings make M2 worse than the coarse regime on one-year RMSE
(160 versus 144 kt). Collapse-window RMSE under annual landings remains
about 821 kt for M2, M3 and M4 (819 kt under the coarse regime; M1 and
M1b are at 694 kt on both treatments, Table 3). On the annual-catch recovery window (train 1995--2007, test 2008--2015)
the ladder scores M1 264, M1b 78, M2 303, M3 609, and M4 586 kt, against
a frozen-persistence error of 104 kt (forecast held at the 2007 level of
81 kt). Annual landings therefore do not repair the recovery-window fit,
and the residual-persistence models deteriorate severely (609 and 586 kt
against 220 and 214 kt under the coarse regime). M1b's 78 kt carries the
same unidentified-Allee caveat as its regime-window fit.

Apparent net production \(S_{t+1}-S_t+C_t\) is strongly negative in
1991--93 even after subtracting reconstructed catch. This quantity is a diagnostic construct, not an exact stock balance: \(S\) here is spawning-stock biomass while \(C_t\) is total
landings, and changes in SSB also reflect maturation into the spawning
stock, changing maturity and weight-at-age schedules, and the age
composition of catch and of natural deaths. Within the surplus-production
accounting the ladder itself assumes --- in which \(S\) is the modelled
biomass state and \(C_t\) is subtracted from it --- matching the observed
\(\Delta S\) with no production at all would require removals of order
\(-\Delta S\), and the observed decline is far larger than \(C_t\). On
that reading the unexplained term dwarfs the difference between the two
catch treatments.
Regular et al.'s \(M \approx 2.5\) peak is an age-structured assessment's
estimate informed by additional data and assumptions; it is not the same
object as this scalar residual, though the two point in the same
direction. Within the model class tested, substituting annual recorded landings
for the coarse regime does not produce the crash in a constant-\(r\)
surplus model.

Reading the same quantity as an implied productivity makes the
non-stationarity explicit. Inverting the Schaefer form at a fixed \(K\)
gives \(\hat r_t = P_t / [S_t(1-S_t/K)]\), a reconstruction of
productivity rather than a forecast model. Periods below are labelled by the year in which a transition starts, so
the last usable transition begins in 2014. At \(K = 1032.7\) kt the
annual-landings series averages \(+1.82\) over the eight transitions
beginning 1983--1990, turns negative through 1991--1994 (mean
\(-0.79\), all four negative), and reads \(+0.37\) over 1995--2007
(thirteen) and \(+0.25\) over 2008--2014 (seven). At \(K = 5000\) kt the
sign pattern is the same --- \(+0.32\), \(-0.55\), \(+0.36\),
\(+0.22\) --- but the magnitudes are not: the post-collapse mean is
\(0.20\) of the pre-collapse mean at the fitted \(K\) and \(1.12\) of
it at \(K = 5000\) kt. The robust feature is therefore the sign
reversal, not the size of the pre-to-post contrast, and no constant
\(r\) reproduces a sequence that changes sign.
The ratio is undefined only at \(S_t = 0\) and \(S_t = K\); for
\(S_t > K\) it remains defined but its denominator is negative, so the
reading as a positive productivity fails. The coarse-regime recovery-window fit, whose \(K\) is 500 kt, cannot be
read this way at all over the pre-collapse states, which exceed its
\(K\). The reconstruction above therefore uses the collapse-window
\(K = 1032.7\) kt and the annual-landings \(K = 5000\) kt, both of
which exceed every observed state.

\begin{figure}
\centering
\includegraphics[width=\linewidth]{figs_e1/fig4_production.png}
\caption{Apparent net production \(P_t = S_{t+1}-S_t+C_t\) against \(S_t\)
under both catch treatments, with the fitted collapse-window and the two
recovery-window Schaefer curves overlaid. \(P_t\) is a diagnostic
construct, not a closed biomass budget: it inherits any error in the
catch reconstruction and any assessment revision in the SSB series, and
it is not a mass balance. The 1991--1993 points lie below every fitted
curve under both treatments, and the two recovery-window
parameterisations are almost indistinguishable over the observed range
of \(S_t\) while diverging widely when extrapolated beyond it.}
\label{fig:production}
\end{figure}

One numerical disclosure travels with the window: the annual-landings M1
(264 kt) and M1b (78 kt) differ from the coarse-regime recovery rows (M1
= 120, M1b = 90) although both treatments sit on the same SSB column.
The reconciliation is that M1's constant catch is not estimated on that
column. It is the training mean of the catch file: 5.0 kt for the coarse
regime, and 3.19 kt (the 1995--2007 mean of the annual landings) for the
annual treatment. The two least-squares objectives are nearly flat at
their minima (MSE 128.35 versus 127.84 kt\textsuperscript{2}), so the
\((r, K)\) minimizer slides along the valley as the constant changes:
\(r = 0.458\) with \(K\) resting at the multi-start initialiser of
500.0 kt --- not a bound, the lower bound on this window being 50.8 kt
--- against \(r = 0.370\) with \(K\) at its upper bound, 5000.0 kt. The optimiser leaves \(K\) at its starting value in this fit; the
profile of Section 3.2, not the optimiser's behaviour, is what
measures the flatness. The
same geometry explains the M1b pair (\(r\) pinned at its
upper bound in both treatments; \(K = 105.8\) against 129.8 kt): the
dependence of the M1b fit on the catch file is the flat-objective
identification of a two-parameter autonomous fit on a 13-year window,
not an inconsistency between fitting objects. The same flat valley caps what the ranking can resolve. Refitting the
recovery-window Schaefer map with \(K\) fixed anywhere on
\([60, 5000]\) kt and \(r\) re-optimised moves the one-step training
objective only from mean squared error \(127.4\) to \(149.9\) kt\(^2\)
(training RMSE \(11.29\)--\(12.24\) kt, a spread under \(0.95\) kt),
while \(r\) compensates from \(0.435\) to \(0.773\). On this window the
\((r, K)\) pair is not identified, and the ordering of valley variants
is not a robust ranking.

The collapse window is weakly identified in the complementary direction,
though less severely than the recovery window. There the stock is high
and weakly trending under large catch, so the data approximately impose
the single condition \(rS(1-S/K)\approx C\) on two parameters. Refitting
with the estimator of Section 2.3 at fixed \(K\) --- the same objective
on increments, the same bounds and multi-start --- raises the training
root mean squared error from \(31.3\) kt at the fitted
\(K = 1032.7\) kt to \(56.3\) kt at \(K = 1500\) kt and \(67.0\) kt at
\(K = 5000\) kt, while \(r\) falls from \(1.935\) to \(0.674\) and
\(0.331\). The objective therefore does discriminate here: a more than
twofold rise in training error is not a flat ridge. What the window does
not do is pin the pair tightly enough for the implied curve maxima to
agree, and the two windows fail in complementary directions --- a stock
near \(K\) constrains the product more than \(r\), and low-biomass
observations constrain \(K\) only weakly, because the density-dependent
term \(\partial g/\partial K = rS^{2}/K^{2}\) is small relative to the
approximately linear component rather than absent.

The implied maxima of the fitted curves make the consequence concrete.
The formal maximum \(rK/4\) and the reference state \(K/2\) of each
printed parameterisation are \(500\) kt and \(516\) kt for the collapse
window, \(57\) kt and \(250\) kt for the coarse-regime recovery fit, and
\(463\) kt and \(2500\) kt for the annual-landings recovery fit. These
are maxima of the fitted scalar curves, not sustainable-yield estimates
for this stock: the accounting caveat above --- SSB net of catch is not
a closed biomass budget --- applies to them unchanged. A spread of under
\(1\) kt in recovery-window training RMSE nevertheless separates curves
whose formal maxima differ by an order of magnitude. The fitted \(r\) is
correspondingly hard to read as a demographic rate: in this discrete
annual map the low-biomass multiplier is \(1+r\), so \(r = 1.935\)
corresponds to an annual factor of \(2.94\) and a continuous-equivalent
\(\log(1+r) = 1.08\,\mathrm{yr}^{-1}\), which is high for a
late-maturing gadoid but is not the same quantity as a continuous-time
intrinsic rate of increase, and SSB can also change through maturation
of existing cohorts. These are identification limits with biological
content, not numerical curiosities. One mechanism
unifies both faces of the catch-dependence: a constant catch shifts the
flat valley's minimiser without moving the objective, so the rolling
M1 and M1b scores (refit at every origin) are insensitive to the catch
treatment --- both are identical after the displayed rounding across the two
treatments (the archived values differ by at most 0.04 kt), because
future catch is not supplied to these modules --- though the
training-mean catch does change with treatment and can move their fitted
parameters, leaving the rolling scores at 121 and 115 kt at
\(h=1\) --- while the single
fixed-window fit is not. The property is specific to the two
constant-catch modules: M2, M3 and M4 take \(C_t\) into the state
update at every step, and their one-year rolling scores do move with the
treatment (144 to 160, 135 to 154, and 196 to 206 kt respectively).

Starting from \(\hat q I_t\) instead of SSB gives one-year RMSE 128 kt
(still worse than persist 98); one-year log-RMSE 0.49 versus persist
0.52; five-year RMSE 331 versus persist 265. On the primary score the
variant is not retained. The log-score difference is recorded and not
used for selection. The survey index contributes to NCAM and therefore
does not provide an independent validation target.

\subsubsection{3.3 Alternative assessment
specification}\label{alternative-assessment-specification}

\begin{figure}[htbp]
\centering
\includegraphics[width=\linewidth]{figs_e1/fig5_xtencam.png}
\caption{The two specifications. Overlap 1983--2015 RMSE = 126 kt (2015: NCAM 299 kt, xteNCAM 273 kt). Different safe sets. Not pooled.}
\end{figure}

Checkpoints from Regular et al.~(2025) Table 17: 2005 SSB = 26 kt (95\%
interval 22--31), 2017 = 451 kt (381--534), 2021 = 423 kt (NCAM and
xteNCAM said to agree), 2024 = 342 kt (246--475), 2024/LRP = 1.24. The
2005 values (26 kt versus NCAM 25.18 kt) can agree on a low year without
the two series sharing a reference point; agreement on a low year does
not warrant splicing the two columns. The 2015 value is
33.8\% of the old LRP, and the series is above the new one after 2016.

\textbf{Table 6.} Rolling-origin RMSE on Specification B (kt). The
persistence and training-mean rows are the origin-matched baselines,
computed on the ladder's own origin sets (\(n=59\) at \(h=1\),
\(n=55\) at \(h=5\)); the mixed-origin baselines over the longer naive
origin sets (\(n=63\) and \(n=59\)) are 88 and 318 kt for persistence
and 449 and 506 kt for the training mean.

\begin{longtable}[]{@{}lrr@{}}
\toprule\noalign{}
Model & \(h=1\) & \(h=5\) \\
\midrule\noalign{}
\endhead
\bottomrule\noalign{}
\endlastfoot
persist & \textbf{84} & \textbf{300} \\
M1 & 120 & 432 \\
M1b & 152 & 446 \\
M2 & 166 & 1059 \\
M3 & 127 & 930 \\
M4 & 206 & 1031 \\
mean & 458 & 522 \\
\end{longtable}

Collapse window (train 1954--89, test 1990--95): M1 RMSE 817 kt; M2 1898
kt. Official landings worsen the crash forecast. No module is retained.
Persistence remains the lowest-RMSE forecast. This is a second non-retention outcome under the same rule on a
different specification, and it is weaker than a statistical null
result. The longer series and the revised LRP
do not justify retaining the additional modules.

Recovery-stall window (train 1995--2012, test 2013--2024; \(n=12\)): M1
254 kt, M1b 178 kt, M2 269 kt, M3 268 kt, M4 269 kt. The 2013--2024 test
path rises from 167 kt to a 2022 peak of 462 kt before easing to 342 kt
in 2024, all above the 112 kt training-end value,
so a persistence forecast frozen at that value errs by 262 kt --- more
than M1 and M1b on this single origin. This fixed-origin comparison does
not enter the retention rule, which is decided on rolling-origin primary
RMSE (Table 6, where persistence is not beaten at either horizon). M1b's
window fit carries the identification fragility familiar from its other
window fits: a small fitted Allee parameter
(\(\mathfrak s = 5.3\times10^{-3}\)) with a carrying capacity (\(K=500\)
kt) extrapolated far beyond the training range (maximum 117 kt).

On the rolling Brier score for \(\mathbf{1}\{\hat S<\mathrm{LRP}\}\) on
Specification B (the secondary score, LRP = 276 kt):
persistence, on its origin-matched sets, \(4/59 = 0.07\) at \(h{=}1\)
and \(16/55 = 0.29\) at \(h{=}5\) (the same counts over the longer naive
origin sets give \(4/63 = 0.06\) and \(16/59 = 0.27\)); M1 0.05 and 0.45;
M1b 0.15 and 0.47; M2 0.07 and 0.36; M3 0.02 and 0.31; M4 0.07 and 0.38.
M1 and M3 improve the one-year Brier score over persistence (0.05 and
0.02 versus 0.07); no model improves it at \(h{=}5\), and no model
improves the primary RMSE score at either horizon (Table 6), so the
retention verdict is unchanged. For reference, the same secondary score on Specification A is 0.00 at
both horizons for persistence. The persistence indicator forecasts
\(\mathbf{1}\{S_t<884.6\}\) from the origin state, and its score is zero
because every origin state is already below the 884.6-kt LRP. In
particular the first origin, \(S_{1990}\), lies below the 1983--89 mean
bound, so the collapse train starts on the downslope of its own
reference period. The target years 1991--2015 are below the LRP as well,
so indicator and outcome agree at every origin. (The bare fact that targets are below the LRP
would not by itself make the persistence indicator correct; the origin
states being below it does.) Persistence supplies no directional prediction, so its sign-hit rate
is reported as NA rather than as a zero hit rate. The structural
scores are M1 0.04/0.05, M1b
0.00/0.05, M2 0.08/0.14, M3 0.08/0.10, M4 0.08/0.14 --- none below
persistence, which scores zero there by construction, and reported only to keep the two specifications' records
separate and complete.

Regular et al.~assign the 1992--94 disappearance primarily to \(M\)
(peak \(\approx 2.5\)), informed by tagging and a capelin/cod predictor,
and note that some of that \(M\) could still be unreported \(F\). The
relation between that attribution and the scalar residual of this ladder
is stated in Section 3.2.

\subsubsection{3.4 Prey-informed
productivity}\label{prey-informed-productivity}

Murphy et al.~(2025) report a 1991 acoustic collapse (1985--90 median
3704 kt versus 1991--2022 median 174 kt). A two-regime \(r\) with break
1991 uses the high regime only for forecasts issued before 1991. Capelin
enters as a candidate disturbance class, not as a fitted parameter.

\textbf{Table 7.} Rolling RMSE, two-regime \(r\). Baselines are origin-matched to the module's own origin sets (Specification A \(n=25/21\); Specification B \(n=59/55\)); the mixed-origin Specification B values over the longer naive sets are 88 and 318 kt.

\begin{longtable}[]{@{}llrr@{}}
\toprule\noalign{}
Specification & Model & \(h=1\) & \(h=5\) \\
\midrule\noalign{}
\endhead
\bottomrule\noalign{}
\endlastfoot
A & persist & \textbf{98} & \textbf{265} \\
A & M\_cap & 154 & 334 \\
B & persist & \textbf{84} & \textbf{300} \\
B & M\_cap & 147 & 894 \\
\end{longtable}

On post-break origins only (Specification B, \(h=1\); \(n=33\)): M\_cap
107 versus persistence 75 on the identical origin set (the all-origins value
of 88 is quoted in the Table 6 caption). Not retained. On Specification A post-break
origins (\(n=24\)) the corresponding M\_cap value is 149 kt.

The year-by-year 3L spring acoustic biomass is tabulated in the
Northwest Atlantic Fisheries Centre's Zenodo record (2025),
10.5281/zenodo.17515115, with 2023 = 331.3 kt from Murphy et al.~(2025).
Because \((I_{\mathrm{known}}/I_{\mathrm{ref}})^{b}>0\) for any
positive index value and any sign of \(b\), the scaled surplus keeps the
sign of \(g(S_t)\): the module can shrink production but cannot make it
negative, so food limitation severe enough to drive a net loss lies
outside what this specification can represent. Surplus is scaled by
\((I_{\mathrm{known}}/I_{\mathrm{ref}})^{b}\),
where \(I_{\mathrm{known}}(t)\) is the last observation at or before
\(t\), subject to the regime rule of Section 2. The index is carried forward from the last observation available at the
origin, and the carry is reset at 1991 only when the last observation
predates the break: an origin at or after 1991 with no post-1991
observation yet is not evaluated. Origins before 1991 draw on pre-break
surveys and are evaluated, so on Specification B the module is scored
from 1988 (\(n=36\) at \(h=1\), \(n=32\) at \(h=5\), of which three
origins precede the break); on Specification A the first eligible origin
is 1991 (\(n=24\) and \(n=20\)). Training transitions with no eligible
index value are dropped, and pre-1991 index observations are admissible
in training for pre-break origins on the same rule. Along a multi-step trajectory \(t\) is the forecast
year being stepped to, not the origin, so a five-year path may use index
values published after its issue date; this is a conditional hindcast
with respect to the index and is not an operational forecast. \(I_{\mathrm{ref}}\) is the training-window median of the
observed index, and \(b\) is a third free parameter fitted jointly with
\(r\) and \(K\) by one-step least squares on the training window. The
module additionally faces degrees-of-freedom starvation: three
parameters (\(r, K, b\)) are fitted on as few as eight training years at
the shortest origins, and after dropping years with no carried index
value the shortest fit uses only six one-step transitions (origin 1991
on Specification A, 1988 on Specification B; the estimator refuses to
fit on fewer) which invites in-sample overfitting and penalizes
the one-year out-of-sample comparison. Survey years are unobserved in
part of the record, so the index-forecast origins number \(n=24\)
(\(h=1\)) and \(n=20\) (\(h=5\)) on Specification A and \(n=36\) and
\(n=32\) on Specification B, against the SSB-based rolling-origin counts
(\(n=25\) at \(h=1\) and \(n=21\) at \(h=5\) on Specification A; Section
4).

\textbf{Table 8.} Rolling RMSE, observed acoustic index. Bold marks the
origin-matched baseline, which is the valid comparison for the module;
the mixed-origin rows over the longer naive origin sets are retained for
reference and are not the comparator. At \(h=5\) on Specification A
the module (262 kt) reads below the mixed-origin baseline (265 kt); that
apparent 3-kt advantage is an origin-mix artefact and is not marked,
because the origin-matched baseline on the module's own origins is
193 kt. The comparison that matters is the matched one.

\begin{longtable}[]{@{}llrr@{}}
\toprule\noalign{}
Specification & Model & \(h=1\) & \(h=5\) \\
\midrule\noalign{}
\endhead
\bottomrule\noalign{}
\endlastfoot
A & persist (mixed origins) & 98 & 265 \\
A & persist (origin-matched) & \textbf{97} & \textbf{193} \\
A & M\_cap\_index & 150 & 262 \\
B & persist (mixed origins) & 88 & 318 \\
B & persist (origin-matched) & \textbf{79} & \textbf{288} \\
B & M\_cap\_index & 132 & 492 \\
\end{longtable}

On the origin-matched reading the five-year near-tie on Specification A
dissolves --- the baseline on the module's own origins reads 193 kt
against the module's 262 kt --- and one-year RMSE is worse under both
readings (150 versus 98 kt on the SSB origins and 97 kt on the
module's); on Specification B the module's 132/492 kt stand against
origin-matched baselines of 79/288 kt. The module is not retained under
either baseline reading.

The scope of that result is set by the transmission the module encodes.
The index scales contemporaneous surplus, so what is tested is whether
prey information available at the origin improves the next steps through
a proportional change in aggregate production. Published maturity
schedules for this stock place the age at 50\% maturity between ages
five and six (Cadigan, 2016), so the pathway from a prey-driven change in
recruitment to a change in spawning-stock biomass is delayed by roughly
half a decade, and a contemporaneous multiplier cannot carry it. The same
timing objection applies to M4, whose one-year stale start is an
information delay rather than a recruitment delay. Non-retention here is
therefore evidence about the tested transmission and its timing, not
evidence that prey availability is unimportant to this stock.

\subsubsection{3.5 Uncertainty on the retention
margins}\label{uncertainty-on-the-retention-margins}

The uncertainty analysis attaches descriptive Diebold--Mariano loss-differential statistics
(Diebold and Mariano, 1995; unweighted HAC truncation at lag \(h-1\)) and
moving-block bootstrap intervals (K\"unsch, 1989; block length
\(\max(h,3)\), 20,000 replications, seeded) to the retention margins of
the retention rule. It is computed from the archived per-origin forecast
files: on Specification B the xteNCAM rolling file; on Specification A
the layer is computed on the primary coarse-regime pass. Its per-origin
rows were regenerated from the registered estimator and archived; the
regeneration reproduces the archived annual-landings per-origin file
and the archived coarse-regime summary exactly, so no score changes.
The layer is unchanged when recomputed on the annual-landings pass.
The persistence baseline is not archived per-origin; it is recomputed on the
identical origin sets from the registered series, and the recomputation
reproduces the recorded origin-matched baselines (98 and 265 kt on
Specification A; 84 and 300 kt on the matched Specification B origins)
--- the script fails loudly otherwise. The moving-block intervals quantify sampling variation in the archived
origin-level loss sequence conditional on the fitted forecast paths; they
are not predictive intervals and do not propagate parameter-estimation,
assessment-revision, catch-reconstruction or covariate uncertainty. The
layer attaches uncertainty to margins the point rule has already
ranked; it changes no frozen verdict,
score, or table value.

The two columns of Table 9 summarise the same comparison through
different approximations. The DM statistic \(z\) tests the mean of the
per-origin squared-loss differential \(d_i = L_{A,i} - L_{B,i}\) against
zero, scaled by a HAC standard error. The interval and \(p\) come from
the moving-block bootstrap of the difference of RMSEs,
\(\sqrt{\overline{L_A}} - \sqrt{\overline{L_B}}\). The two effect
measures are monotonically related, since
\(\sqrt{\overline{L_A}}-\sqrt{\overline{L_B}} =
(\overline{L_A}-\overline{L_B})/(\sqrt{\overline{L_A}}+
\sqrt{\overline{L_B}})\), so they share a sign and a zero null; the
square root rescales the gap but does not by itself make one summary
more stable than the other. What differs is the approximation used to
gauge sampling variability --- a HAC-scaled asymptotic normal reference
against a resampling distribution over blocks --- and in samples this
small and this dependent the two need not agree. Four of the
twenty-eight rows have an interval excluding zero while \(|z| < 1.96\)
(Specification A \(h=1\) M4 against M3; Specification B \(h=1\) M3
against persistence; Specification B \(h=5\) M1b against persistence and
M4 against M3; see Section SI-3). We therefore report the two
summaries side by side
and do not treat the bootstrap interval as a second calibrated test. The interval and the \(p\) are mutually
consistent in every row. The quantity reported as \(p\) is the percentile-tail fraction
\(2\min\{\#(\Delta^{*}_b\le 0),\ \#(\Delta^{*}_b\ge 0)\}/B\) over the
\(B\) block-resampled differences, obtained by inverting the percentile
interval rather than from a null-centred resampling scheme; it is an
interval-based summary and not a calibrated test of a sharp null, and a
reported value of zero means no resampled difference fell on the
opposite side of zero; where the two columns disagree both are
reported and neither carries a verdict.

Because the origins overlap at \(h=5\), the estimation window expands,
the predictand is a smoothed reconstruction and the sample straddles a
structural break, the automatic \(h-1\) HAC truncation and the
\(\max(h,3)\) block length should not be taken as neutral choices. The
M1-against-persistence margin was therefore recomputed over a grid of
both (Section SI-3). The DM statistic crosses the conventional
threshold within the range of defensible bandwidths, upward on
Specification A and downward on Specification B, so nothing rests on it;
the bootstrap is markedly more stable across block lengths \(3\) to
\(9\). The qualitative reading --- Specification A within noise,
Specification B separating --- is therefore a bootstrap result rather
than a DM result, established for that comparison and not asserted for
every module. Read
``within noise'' as ``this analysis does not resolve the difference'',
not as evidence that the models are equivalent.

\textbf{Table 9.} Uncertainty on the retention margins. Both
specifications are scored on their primary pass: Specification A rows
are the coarse catch regime, so the RMSE values reproduce Table 4, and
Specification B rows are the xteNCAM rolling pass.
(Diebold--Mariano with HAC lag \(h-1\); K\"unsch moving-block
bootstrap, block length \(\max(h,3)\), 20,000 replications, seeded).
Gaps are module minus comparator, in kt; positive gaps mean the module
is worse.

{\footnotesize\setlength{\tabcolsep}{3pt}
\begin{longtable}[]{@{}
  >{\raggedright\arraybackslash}p{(\columnwidth - 22\tabcolsep) * \real{0.0550}}
  >{\raggedleft\arraybackslash}p{(\columnwidth - 22\tabcolsep) * \real{0.0500}}
  >{\raggedright\arraybackslash}p{(\columnwidth - 22\tabcolsep) * \real{0.0850}}
  >{\raggedright\arraybackslash}p{(\columnwidth - 22\tabcolsep) * \real{0.1400}}
  >{\raggedleft\arraybackslash}p{(\columnwidth - 22\tabcolsep) * \real{0.0950}}
  >{\raggedleft\arraybackslash}p{(\columnwidth - 22\tabcolsep) * \real{0.0950}}
  >{\raggedleft\arraybackslash}p{(\columnwidth - 22\tabcolsep) * \real{0.1000}}
  >{\raggedleft\arraybackslash}p{(\columnwidth - 22\tabcolsep) * \real{0.1000}}
  >{\raggedleft\arraybackslash}p{(\columnwidth - 22\tabcolsep) * \real{0.1000}}
  >{\raggedright\arraybackslash}p{(\columnwidth - 22\tabcolsep) * \real{0.1050}}
  >{\raggedleft\arraybackslash}p{(\columnwidth - 22\tabcolsep) * \real{0.0600}}@{}}
\toprule\noalign{}
Spec & \(h\) & Module & Comparator & \(n\) & RMSE module & RMSE comp. & Gap (kt) & DM \(z\) & 95\% CI (kt) & \(p\) \\
\midrule\noalign{}
\endhead
\bottomrule\noalign{}
\endlastfoot
A & 1 & M1 & persist & 25 & 120.5 & 98.0 & +22.5 & 1.13 & {[}\ensuremath{-}13.4, +52.2{]} & 0.192 \\
A & 1 & M1b & persist & 25 & 114.8 & 98.0 & +16.8 & 0.98 & {[}\ensuremath{-}19.4, +59.6{]} & 0.384 \\
A & 1 & M2 & persist & 25 & 144.1 & 98.0 & +46.0 & 1.29 & {[}\ensuremath{-}12.3, +83.0{]} & 0.452 \\
A & 1 & M3 & persist & 25 & 134.6 & 98.0 & +36.6 & 0.99 & {[}\ensuremath{-}16.8, +64.2{]} & 0.932 \\
A & 1 & M4 & persist & 25 & 195.6 & 98.0 & +97.5 & 1.53 & {[}\ensuremath{-}9.1, +172.7{]} & 0.278 \\
A & 1 & M2 & M1 & 25 & 144.1 & 120.5 & +23.5 & 0.91 & {[}\ensuremath{-}60.7, +66.9{]} & 0.684 \\
A & 1 & M4 & M3 & 25 & 195.6 & 134.6 & +60.9 & 1.21 & {[}+4.4, +134.4{]} & 0.000 \\
A & 5 & M1 & persist & 21 & 288.7 & 264.7 & +24.0 & 1.84 & {[}\ensuremath{-}43.1, +56.4{]} & 0.279 \\
A & 5 & M1b & persist & 21 & 288.6 & 264.7 & +23.9 & 1.84 & {[}\ensuremath{-}43.1, +56.4{]} & 0.284 \\
A & 5 & M2 & persist & 21 & 398.2 & 264.7 & +133.5 & 1.10 & {[}\ensuremath{-}21.0, +217.5{]} & 0.421 \\
A & 5 & M3 & persist & 21 & 366.4 & 264.7 & +101.6 & 1.12 & {[}\ensuremath{-}0.8, +151.4{]} & 0.078 \\
A & 5 & M4 & persist & 21 & 488.2 & 264.7 & +223.5 & 1.14 & {[}\ensuremath{-}20.8, +409.4{]} & 0.405 \\
A & 5 & M2 & M1 & 21 & 398.2 & 288.7 & +109.5 & 0.98 & {[}\ensuremath{-}77.1, +232.8{]} & 0.805 \\
A & 5 & M4 & M3 & 21 & 488.2 & 366.4 & +121.9 & 1.15 & {[}\ensuremath{-}27.4, +263.2{]} & 0.707 \\
B & 1 & M1 & persist & 59 & 119.5 & 84.4 & +35.0 & 2.81 & {[}+2.7,
+70.8{]} & 0.032 \\
B & 1 & M1b & persist & 59 & 151.6 & 84.4 & +67.2 & 3.60 & {[}+18.3,
+117.7{]} & 0.009 \\
B & 1 & M2 & persist & 59 & 166.0 & 84.4 & +81.6 & 3.22 & {[}+36.4,
+122.8{]} & 0.000 \\
B & 1 & M3 & persist & 59 & 126.7 & 84.4 & +42.3 & 1.85 & {[}+1.0,
+92.5{]} & 0.042 \\
B & 1 & M4 & persist & 59 & 205.7 & 84.4 & +121.3 & 3.20 & {[}+53.7,
+193.0{]} & 0.000 \\
B & 1 & M2 & M1 & 59 & 166.0 & 119.5 & +46.6 & 1.78 & {[}\ensuremath{-}20.5,
+104.5{]} & 0.170 \\
B & 1 & M4 & M3 & 59 & 205.7 & 126.7 & +79.0 & 3.07 & {[}+35.6,
+116.4{]} & 0.000 \\
B & 5 & M1 & persist & 55 & 431.9 & 300.0 & +131.9 & 2.06 & {[}+33.2,
+218.6{]} & 0.008 \\
B & 5 & M1b & persist & 55 & 445.5 & 300.0 & +145.5 & 1.80 & {[}+18.2,
+250.9{]} & 0.023 \\
B & 5 & M2 & persist & 55 & 1058.9 & 300.0 & +758.9 & 2.15 & {[}+363.6,
+1038.5{]} & 0.000 \\
B & 5 & M3 & persist & 55 & 930.1 & 300.0 & +630.2 & 2.39 & {[}+352.8,
+845.6{]} & 0.000 \\
B & 5 & M4 & persist & 55 & 1030.7 & 300.0 & +730.8 & 2.35 & {[}+407.3,
+978.3{]} & 0.000 \\
B & 5 & M2 & M1 & 55 & 1058.9 & 431.9 & +627.0 & 1.97 & {[}+195.7,
+929.7{]} & 0.004 \\
B & 5 & M4 & M3 & 55 & 1030.7 & 930.1 & +100.6 & 1.88 & {[}+20.2,
+177.4{]} & 0.007 \\
\end{longtable}}

Splitting the same origins by period rather than by horizon shows what
the pooled score aggregates. On Specification B at \(h=1\), computed on
the archived per-origin records with no refit, persistence reads
\(94.7\) kt before 1991 (\(n=26\)) against \(142.4\) kt for M1;
\(167.0\) kt across 1991--1995 (\(n=5\)) against \(126.8\) kt for M1
and \(66.6\) kt for M1b; \(20.7\) kt across 1996--2012 (\(n=17\))
against \(19.7\) kt for M3; and \(61.0\) kt from 2013 (\(n=11\))
against \(96.4\) kt for M3. These rolling scores answer a different question from the fixed-window
collapse test of Section 3.1, where models fitted on 1983--1990 are
asked whether a pre-collapse fit anticipates the crash, and none does.
Here each forecast is issued from an origin already inside the decline,
with training extended to that origin, so the question is whether the
next year is closer to a refitted model or to the last observed state; a
one-step forecast issued in 1993 from an already-collapsed state is not
evidence that a model predicted the collapse. On that second question
persistence is the lowest-error forecast in the two quiescent bands but
not in the collapse band, where the structural modules read below it on
five origins; leave-one-out on
that band leaves M1 below persistence in all five cases, so the ordering
is stable, though five origins remain too few to score. The
1996--2012 margin is \(0.93\) kt on a \(20.67\)-kt baseline, inside the
\(5\%\) tie band of \(1.03\) kt and with a bootstrap interval of
\([-14.2, +14.9]\) kt that includes zero: a tie under the stated rule,
not a win. No retention verdict changes, because the rule scores pooled
rolling RMSE at both horizons and the one sub-band advantage falls
inside the band the rule declares non-deciding. The reading that changes
is interpretive: the pooled result combines regimes in which persistence
wins by a wide margin with a collapse band in which it does not, so the
non-retention statement is a statement about the pooled score rather
than about every period of this series.

Readings. On Specification A no non-retention margin against persistence separates
from zero. At \(h=1\) the deficits against persistence (16.8--97.5 kt
on \(n=25\)) carry DM \(z\) from 0.98 to 1.53 with bootstrap \(p\) from
0.19 to 0.93; at \(h=5\) (23.9--223.5 kt on \(n=21\)) \(z\) runs from
1.10 to 1.84 with \(p\) from 0.08 to 0.42. The heavy-tailed
collapse-window losses dominate the bootstrap variance at this sample
size, so the Specification A outcome is a point-rule ranking whose
margins lie within noise: the scores suffice to order the models and do
not resolve small skill differences. On Specification B the non-retention
separates: every module's deficit against persistence has a bootstrap interval excluding zero at
\(h=1\) (\(p\) from 0.042 to below 0.001 on \(n=59\)) and at \(h=5\)
(\(p\) at most 0.023 on \(n=55\)). The declared (H1) comparators behave
asymmetrically: M2's deficit against M1 is within noise at \(h=1\) on
both specifications and at \(h=5\) on Specification A, with the interval excluding zero only
on Specification B at \(h=5\) (\(p=0.004\)), while M4's delay cost
against M3 separates at \(h=1\) on both specifications (\(p\) below
0.001) and at \(h=5\) on Specification B (\(p=0.007\)) but not on
Specification A (\(p=0.71\)). The alternative comparator rows are held in the archive rather than
printed in Table 9. On the primary passes they read: Specification A
\(h=1\), gap \(+29.3\) kt, \(z=0.91\), interval
\([-67.6, +82.8]\) kt; \(h=5\), \(+109.7\) kt, \(z=0.98\),
\([-77.1, +232.8]\) kt; Specification B \(h=1\), \(+14.4\) kt,
\(z=0.51\), \([-67.5, +89.1]\) kt; \(h=5\), \(+613.4\) kt,
\(z=1.93\), \([+170.9, +945.1]\) kt. Accordingly, on them the reading (M2
against M1b) is within noise at \(h=1\) on both specifications and at
\(h=5\) on Specification A, separating on Specification B at \(h=5\)
(\(p=0.005\)): the comparator declaration of the retention rule is
immaterial to every recorded verdict, as the protocol disclosure
records. The layer is produced by a registered script, deterministic under seed 0,
with its output archived alongside it; Section SI-3 names both.

\subsubsection{3.6 Fitted parameters}\label{fitted-parameters}

This section collects in one place every fitted-parameter value the
article prints, the window or treatment it belongs to, and the bound it
attains. The table collects values reported in the preceding
sections; nothing is recomputed. Quantities not reported individually
(the per-origin rolling fits of every module, the index module's
per-window exponent \(b\), and the delay module's structural setting)
are held in the archive rather than
invented. The collection exists so the identification record --- bounds
attained, interior fits, the flat valley --- can be read at a glance.

\textbf{Table 10.} Fitted parameters, with the section in which each is
reported.

\begin{longtable}[]{@{}
  >{\raggedright\arraybackslash}p{(\columnwidth - 6\tabcolsep) * \real{0.2500}}
  >{\raggedright\arraybackslash}p{(\columnwidth - 6\tabcolsep) * \real{0.2500}}
  >{\raggedright\arraybackslash}p{(\columnwidth - 6\tabcolsep) * \real{0.2500}}
  >{\raggedright\arraybackslash}p{(\columnwidth - 6\tabcolsep) * \real{0.2500}}@{}}
\toprule\noalign{}
Module & Window / treatment & Fitted values & Bound / identification status \\
\midrule\noalign{}
\endhead
\bottomrule\noalign{}
\endlastfoot
M1 & Collapse, Specification A (train 1983--1990) (\S{}3.1, \S{}4) &
\(r = 1.935\), \(K = 1032.7\) kt, constant catch \(240\) kt; repeller
\(144\) kt, attractor \(889\) kt, monotone below \(783\) kt,
\(F'(S^*) \approx -0.39\) & \S{}3.1 records the collapse-window fitted
\(r\) as constrained near, but not attaining, the upper bound of \(2\) \\
M1 & Recovery, coarse catch regime (\S{}3.2) & \(r = 0.458\), constant
catch \(5.0\) kt (the training mean); training MSE \(128.35\) kt\textsuperscript{2}
& \(K\) at \(500.0\) kt, the multi-start initialiser, not a bound (the
lower bound on this window is \(\max_{\mathrm{train}}S + 10 =
50.8\) kt) \\
M1 & Recovery, annual landings (\S{}3.2) & \(r = 0.370\), constant catch
\(3.19\) kt (the 1995--2007 landings mean); training MSE \(127.84\)
kt\textsuperscript{2} & \(K\) at its upper bound, \(5000.0\) kt \\
M1b & Recovery, coarse catch regime (\S{}2.2, \S{}3.2) &
\(\mathfrak{s} = 2.1\times10^{-23}\), \(r = 2.000\),
\(K = 105.8\) kt & \(\mathfrak{s}\) approaches lower bound to numerical tolerance; \(r\)
pinned at its upper bound in both treatments; \(K\) a valid interior
fit \\
M1b & Recovery, annual landings (\S{}3.2) &
\(\mathfrak{s} = 9.4\times10^{-6}\), \(r = 2.000\),
\(K = 129.8\) kt & \(\mathfrak{s}\) at the lower boundary to
numerical tolerance \\
M1b & Recovery-stall, Specification B (train 1995--2012) (\S{}3.3) &
\(\mathfrak{s} = 5.3\times10^{-3}\), \(K = 500\) kt, training-window
maximum \(117\) kt & identification fragility: \(\mathfrak{s}\)
small, \(K\) extrapolated far beyond the training range \\
M3 & Rolling, Specification A (\S{}3.1) & \(\phi = 0.95\) (AR(1) residual
coefficient) & printed value; the per-window rolling \(\phi\) values are
archive items \\
\texttt{M\_cap\_index} & Index module (\S{}3.4) & \(r\), \(K\), \(b\) fitted jointly by
one-step least squares & fitted values not printed (archive);
degrees-of-freedom note at \(n = 8\) training years \\
M1 & Flat-valley sweep, recovery window (\S{}3.2) & \(K\) fixed
anywhere on \([60, 5000]\) kt moves the one-step training objective only
from MSE \(127.4\) to \(149.9\) kt\textsuperscript{2} (training RMSE \(11.29\)--\(12.24\)
kt) while \(r\) compensates over \([0.435, 0.773]\) & the \((r, K)\)
pair is not identified on this window; the ordering of valley variants
is not a robust ranking \\
M4 & Rolling, all specifications & \(r\), \(K\), \(\phi\) refitted by
the same procedure as M3 at every origin; the module adds no parameter
of its own, the one-year delay being structural & archive: the refitted
per-origin values coincide with M3's and are not printed separately \\
All modules & Per-origin rolling fits & archive items (not printed);
M1b's Specification B rolling Allee optimum is environment-sensitive at
the \ensuremath{\pm}17 kt level (\S{}Data availability) & --- \\
\end{longtable}

Bounds: \(r \in (0.001, 2]\), and \(K\) optimised on
\([\max_{\mathrm{train}} S + 10, 5000]\) kt, where the maximum is taken
over the states the one-step regression conditions on. On the
1995--2007 recovery window those are the years 1995--2006, with maximum
\(40.8\) kt, giving a lower bound of \(50.8\) kt. Optimisation starts
from \(K = 500\) kt. The recovery-window M1 fit reports \(K = 500.0\)
kt: the optimiser does not move from its starting value, which lies
\(449\) kt above the bound, because the objective is flat there. This is
the same flat-valley geometry quantified in Section 3.2, and a
stationary starting value is a sharper symptom of it than an active
bound would be.

\subsubsection{3.7 Operating characteristics of the retention
rule}\label{operating-characteristics}

The result above is negative, so its interpretation turns on whether the
rule could have detected added structure had it been present. That
question cannot be answered from the observed series alone. It is
addressed by simulation, under a design --- data-generating processes,
replicate count, and the thresholds separating an adequate from an
inadequate rule --- fixed in advance of generating any synthetic
series.

Five processes were simulated, each a member of the ladder's own class
and parameterised from the fitted values of Section 3.6, using the
estimation and trajectory code of Section 2.3 unmodified: an autonomous
Schaefer map at the collapse-window fit (D1) and at the recovery-window
fit (D2), a stock-flow map with the coarse regime path (D3), a
depensatory map with an identifiable threshold \(\mathfrak s = 15\) kt
(D4), and a persistence-true null with no surplus term (D5). Each was
run at \(\sigma = 11.8\) and \(33.8\) kt, the archived recovery- and
collapse-window residual standard deviations, on 33-year series matching
Specification A, with 200 seeded replicates per cell. The full retention
rule was applied unchanged.

\begin{longtable}[]{@{}llrr@{}}
\toprule\noalign{}
Process & Generating module & \(\sigma=11.8\) & \(\sigma=33.8\) \\
\midrule\noalign{}
\endhead
\bottomrule\noalign{}
\endlastfoot
D1 autonomous, collapse fit & M1 & 0.965 & 0.985 \\
D2 autonomous, recovery fit & M1 & 0.710 & 0.130 \\
D3 stock-flow & M2 & 0.090 & 0.110 \\
D4 depensation, \(\mathfrak s=15\) kt & M1b & 0.005 & 0.015 \\
D5 persistence-true & none retained & 0.985 & 0.970 \\
\end{longtable}

The first four rows give the proportion of replicates in which the
generating module was retained; the last gives the proportion in which
no structural module was retained when persistence was true. The
false-retention rate across all structural processes was \(0.044\).

\begin{figure}
\centering
\includegraphics[width=0.86\linewidth]{figs_e1/fig6_power.png}
\caption{Operating characteristics of the retention rule. Cells give
the proportion of replicates in which the generating module was retained
(D1--D4) or in which no structural module was retained under a
persistence-true process (D5), against the thresholds fixed before
execution: a tick marks power at or above \(0.80\) and specificity at or
above \(0.90\), a cross marks power below \(0.30\). Power is adequate
only where the generating signal is strong (D1); the rule is specific
throughout.}
\label{fig:power}
\end{figure}

Two properties follow. The rule is \emph{specific against in-class
alternatives}: under a persistence-true process drawn from the ladder's
own family it declines to retain structure in \(97\)--\(99\%\) of
replicates, so the empirical non-retention is not an artefact of an
instrument that rejects everything. That specificity does not extend
beyond the class. Two further processes were simulated whose truth the
ladder cannot represent --- productivity declining smoothly with time,
and a deterministic trajectory observed with error --- and the rule
retained a structural module in \(68\%\) to \(98\%\) of replicates,
almost always the autonomous map. Against a truth outside its class the
rule over-retains, and the specificity figure above should be read as
conditional on the class rather than as a property of the rule. It is also \emph{conditionally
powerful}: where the generating signal is strong, as in the
high-biomass, large-catch conditions of D1, it recovers the true module
in \(97\)--\(99\%\) of replicates.

Power is nevertheless low for three of the four structural processes,
and the binding constraint is identification rather than the rule's
gates. Requiring a module only to beat persistence, ignoring the
comparator and the band, still retains the true module in just
\(33\%\) of D3 and \(18\%\) of D4 replicates; conditional on clearing
that bar the comparator requirement then removes a further \(69\%\) and
\(94\%\). With five modules, chance alone would place the generating
module first \(20\%\) of the time, and on D3 it holds that position in
only \(3.8\%\) of replicates --- less often than a random draw, with a
mean one-year error above that of two rival modules. Section SI-5 gives
the full decomposition.

The consequence for the empirical result is that its strength is not
uniform across modules. Had the collapse window been generated by an
autonomous map at the fitted parameters, the rule would have retained
M1 in \(97\%\) of replicates; it does not retain M1 on the observed
series, so those data are inconsistent with that generating process and
the non-retention is evidence about the stock rather than about the
instrument. Non-retention of the stock-flow and depensation modules
under recovery-window conditions is by contrast weak evidence against
those structures, because the rule would have retained them in at most
\(11\%\) and \(1.5\%\) of replicates even had they been true. The
residual and lagged-initialisation modules were not simulated as
generating processes, so no power estimate is available for them; their
power is plausibly bounded by the stock-flow figure, since both extend
the same base map and must clear an additional comparator, but that is
an inference rather than a measurement.

Two limits are inherent to the design. The processes are members of the ladder's own class, so this measures power
against correctly specified alternatives and the D1 figure is an upper
bound; and the synthetic predictand carries no assessment smoothing, so
persistence is a weaker baseline here than on the reconstructed series.

The same ladder has been scored on a groundwater system, the Edwards
Aquifer, in Abaee (2026, Does a one-pool water-balance model improve forecasts of Edwards Aquifer head), and the comparison isolates a property that
this series cannot show. That analysis applied the retention rule
without a tie band and horizon by horizon; the figures below restate its
verdict under the rule used here, so that both systems are judged by one
criterion. The restatement changes no outcome: the retained set is empty
on both systems under either form of the rule. Six Edwards margins fall
inside the \(5\%\) band, and the module they bear on was already
excluded on other grounds. There the autonomous module reads below
persistence on the point ranking at one year (12.84 against 13.23 ft)
and is still not retained, because the margin of \(2.96\%\) falls
inside the band and the ordering reverses at five years (21.25 against
21.11 ft). On the present
series no module comes within the band at either horizon: the closest
approach is M1b at \(9.0\%\) adrift at five years and \(17.1\%\) at one
year, against a band of \(5\%\). The two cases therefore reach the same verdict by
different routes: here the ranking alone decides, whereas on the
groundwater series the tie band and the two-horizon requirement are what
withhold retention from a module the point rule would have kept. Neither
series is pooled with the other and no verdict is transferred; what
recurs is the decision rule, not the data or the outcome.

Because the simulation archives every replicate's scores, alternative
decision rules can be evaluated on the same synthetic data without
refitting. Averaging over the four structural processes, and over the
persistence-true process for specificity:

\begin{longtable}[]{@{}lrr@{}}
\toprule\noalign{}
Decision rule & mean power & specificity \\
\midrule\noalign{}
\endhead
\bottomrule\noalign{}
\endlastfoot
Retention rule as adopted (H1--H3, 5\% band) & 0.376 & 0.978 \\
Without the comparator gate (H2, H3 only) & 0.476 & 0.972 \\
Persistence only, \(h=1\), 5\% band & 0.562 & 0.955 \\
Persistence only, any margin & 0.542 & 0.765 \\
Full rule, no tie band & 0.436 & 0.772 \\
Full rule, 10\% band & 0.337 & 0.998 \\
MASE \(<1\) against the naive benchmark & 0.651 & 0.675 \\
Information criterion, \(n\log\mathrm{MSE}+2k\) & 0.509 & 0.992 \\
\end{longtable}

The last two rows place the rule against selection criteria in common
use. Scoring by scaled error alone raises power to \(0.651\) but drops
specificity to \(0.675\); an information criterion penalising free
parameters does better than the rule adopted here on both axes,
reaching \(0.509\) power at \(0.992\) specificity. A criterion penalising parameter count therefore dominates the rule
applied here on both axes, and the tie band is a coarser device for the
same purpose. Within the rule as specified, the band
carries most of the weight --- removing it costs \(0.206\) of
specificity to buy \(0.060\) of power --- while the comparator gate
costs \(0.100\) of power for \(0.006\) of specificity. A diagnostic that
identified in advance where the rule has power would be more useful
than any variant, and Section SI-5 reports two candidates that fail:
the better-correlated one is computed from the same scores the rule
consumes and misclassifies both stock-flow cells. Identifying when the
rule has power remains open.

A module may improve the score and still not constitute added
structure. In the companion, the
groundwater stock-flow variant fitted with training-mean fluxes beats
persistence at one year yet was declined on the ground that it collapses
to an autoregression when those fluxes are held constant. M1b raises the
same question in a different form. Its fitted threshold descends towards
an unattained infimum, at which point the production term is the
zero-threshold cubic branch of Definition 2.1 rather than a
positive-threshold Allee model, so its lower error on Specification A is
not evidence for depensation. The scored rule does not encode this
distinction --- it compares errors --- and both cases are resolved by
inspecting what the fitted module has become, not by the score.

\subsubsection{3.8 The outcome years (2016--2021) on the Specification-A horizon}\label{outcome-years}

Specification A is scored on 1983--2015. The assessment of record for the
same model family was subsequently updated, and the per-year continuation
of the series through 2021 is published in the RAM Legacy Stock Assessment
Database (v4.66, stock COD2J3KL; the 2021-vintage reading of the same
NCAM-family assessment). Table~\ref{tab:outcome} reports the outcome years
under both published vintages, with the persistence error from the 2015
origin.

\begin{table}[h]
\centering
\small
\begin{tabular}{@{}lccc@{}}
\toprule
Year & RAM v4.66 SSB (kt) & xteNCAM SSB (kt) & Persistence error (kt) \\
\midrule
2016 & 340 & 339 & 63 \\
2017 & 433 & 451 & 156 \\
2018 & 394 & 369 & 117 \\
2019 & 419 & 400 & 142 \\
2020 & 440 & 414 & 163 \\
2021 & 411 & 423 & 134 \\
\bottomrule
\end{tabular}
\caption{The outcome years on the Specification-A horizon. Persistence
error is $|S_t - S_{2015}|$ from the RAM-vintage origin of 277 kt; the
xteNCAM column is the independent model's reading of the same calendar
years.}
\label{tab:outcome}
\end{table}

Three facts are relevant to the frozen test. First, the outcome window is
an increase regime: every outcome year lies above its 2015 origin in both
vintages, so the window contains no collapse event of the kind the fixed
pre-collapse origin test was designed to probe, and the frozen verdict is
neither confirmed nor reversed by it. Second, the two vintages read the
same years within 1--26 kt (mean absolute difference 16.8 kt), so the
outcome is vintage-robust even where the historical part of the series is
not (Section 3.9). Third, persistence errs low throughout: from the 2015
origin the rule that nothing changes under-predicts every outcome year,
with mean absolute error 129.2 kt and RMSE 133.3 kt on the RAM vintage
(107.5 kt mean absolute error from the Table A2 origin of 298.65 kt) ---
the same order as the scored one-step persistence RMSE of 98 kt on
1983--2015. Re-scoring the module ladder on the extension is archived
pipeline work under the frozen specification and runners of the
reproducibility statement; this section claims the outcome facts and the
baseline errors, and nothing about module rankings on the extension.

\subsubsection{3.9 Threshold definition-dependence on the published series}\label{definition-dependence}

The two specifications use different thresholds on different series. The
published files allow the stricter experiment: both thresholds on one
series. On the xteNCAM file (1954--2024; its ratio column fixes
\(B_{\mathrm{lim}}\) at 276 kt, with per-year implied values 250--289 and
mean 275.5), classifying each year under \(B_{\mathrm{lim}} = 276\) kt
and under the 2010/2016 LRP of 884.6 kt gives 45 years at or above
\(B_{\mathrm{lim}}\), 16 years at or above the LRP (1954--1969 only),
and 29 disagreement years in two blocks: 1970--1992 without
1977--1979, and every outcome year 2016--2024. Three readings follow.
(i) Under the 1980s-mean threshold the stock stands below the reference
point for 55 consecutive years (1970--2024); under the
\(B_{\mathrm{lim}}\) threshold it is below for 26 years (1977--1979 and
1993--2015) and above from 2016 onward. The late-1970s episode (ratios
0.69--0.95) is a critical-zone reading, on the current model, of years
the contemporary fishery regarded as healthy. (ii) The threshold itself
is vintage-dependent: the 1983--1989 mean that defines the LRP equals
884.6 kt on the Table A2 vintage and 790.0 kt on the RAM vintage --- the
reference point moves by 94.6 kt across vintages of the same assessment
family. (iii) The verdict on identical data is therefore
threshold-relative in both senses, of definition and of vintage; the
October 2023 framework revision, which moved the stock's reported zone
from critical to cautious without a change in the stock, is this
experiment executed by the assessment authority itself. The scored RMSE
comparison is untouched --- it never uses the thresholds --- but every
normative verdict layered on the forecasts, the Brier indicator score
included, inherits the dependence.

\subsubsection{3.10 The 2022--2023 observation suspension as a review-timing event}\label{suspension}

The frozen test conditions on a complete annual series. The stock's real
2022--2023 history violates that condition. Partial coverage by the fall
research-vessel survey in 2021 and 2022 left the leading indicator of
stock size uncomputable; no typical assessment update was possible in
2022 or 2023; the stewardship maximum was rolled over unchanged at
12{,}999 t; and the survey vessel that carries the series was condemned
in February 2023. Table~\ref{tab:suspension} assembles the episode from
the assessment of record and the public record.

\begin{table}[h]
\centering
\small
\begin{tabular}{@{}lp{0.72\linewidth}@{}}
\toprule
Period & Event \\
\midrule
2021--22 & Partial fall-survey coverage; leading indicator not computable \\
2022--23 & No typical assessment update; stewardship maximum held at 12{,}999 t \\
Feb 2023 & Survey vessel condemned \\
Mar 2024 & Revised-framework assessment: SSB at 1.2 \(B_{\mathrm{lim}}\) (95\% CI 0.7--2.1; 22\% critical-zone probability) \\
Jun 2024 & Commercial fishery reopened at 18{,}000 t after 32 years \\
2025 & TAC 38{,}000 t; reference point revised after being deemed overestimated \\
Apr 2026 & Healthy-zone report: SSB $\approx$ 540 kt (CI 420--700); 70\% probability above the upper stock reference \\
\bottomrule
\end{tabular}
\caption{The 2022--2023 observation suspension and its resolution
(sources: Regular et al. 2025, Table 17 file; DFO reopening release
2024; the 2025 and 2026 assessment coverage).}
\label{tab:suspension}
\end{table}

Read against the review-timing rule --- a review must land no later than
the worst-case loss of the information it consumes --- the episode is
that rule's boundary case in the wild: the information loss ran past the
review cadence, and the removals rule froze for two cycles while the
observation map was dark. The companion monitoring paper states the formal identity (viability at the deadline requires the review no
later than worst-case exit, boundary inclusive); the 2022--2023
suspension is its empirical instance, and it marks the practical limit
of any scoring exercise that presumes the series exists.

\subsection{4. Discussion}\label{discussion}

Section 3.7 separates two readings of the empirical result. Where the
rule has power, non-retention is evidence about this series; where it
does not, non-retention constrains little. The discussion below is
scoped accordingly: mechanisms are proposed for the autonomous module,
whose non-retention the simulation shows to be informative, and withheld
for the stock-flow and depensation alternatives, whose are not.

The observed path falls far below the fitted lower equilibrium and then
recovers. No exact fixed autonomous noise-free scalar trajectory can
reproduce that. The obstruction is structural, and is stated below for
the map held at constant catch; the scored collapse failure has a
simpler and more direct cause, given after the proposition.

\textbf{Proposition 4.1 (No recovery from below the lower equilibrium,
for autonomous noise-free scalar surplus maps).} \emph{Let
\(S_{t+1}=F(S_t)\) be an autonomous one-dimensional surplus-production
map of Definition 2.1 evaluated with \(\varepsilon_t \equiv 0\) and a
constant catch \(C_t \equiv C\), with at most two positive equilibria --- a lower
repelling point \(S_-\) and an upper attracting point --- and a one-step
map monotone below its maximum; the single-equilibrium monotone regime
is the special case, and the fitted collapse-window map (\(r = 1.935\),
\(K = 1032.7\) kt, constant catch \(240\) kt) has exactly this form
(repeller \(144\) kt, attractor \(889\) kt, monotone below \(783\) kt).
If the observed path \((S_t)\) falls strictly below \(S_-\) at some
time and is strictly larger at some later time, then no trajectory of
\(F\) reproduces \((S_t)\) exactly. The Specification A collapse path
satisfies the hypothesis: NCAM SSB falls to \(9.7\) kt by 1995, far
below the fitted repeller of \(144\) kt, and rises thereafter. }

The argument is immediate: below the lower equilibrium \(F(S) < S\), and
\([0, S_-)\) is forward-invariant because \(F\) is monotone there and
\(F(S_-) = S_-\), the interval being closed at the left because the
implementation clips the state at a non-negative floor. Every trajectory
entering \([0, S_-)\) therefore decreases toward the floor and never
returns above \(S_-\), so a path that goes below \(S_-\) and is later
higher is not a trajectory of \(F\).

Two limits on the scope of this proposition should be stated
explicitly, because both are easy to overstate. First, it does not say
that autonomous maps cannot produce non-monotone paths: the fitted map
has \(F'(S^\star) \approx -0.39\) at its attractor, so approach to the
attractor is by damped oscillation, and declines followed by recoveries
are the generic local behaviour there (from \(950\) kt the fitted map
gives \(950 \to 857 \to 899\) kt). What is excluded is recovery after a
fall below the lower repeller, which is the feature the observed
collapse path actually has. Second, the result is a statement about
autonomous noise-free scalar maps only. It does not apply to the
stochastic map of Definition 2.1 with \(\varepsilon_t \neq 0\), nor to
M3 and M4, which carry a residual state and a stale initial state
respectively and are therefore not autonomous scalar maps. The
proposition applies directly to M1 under the deterministic forecast
convention used here, and to M1b whenever its own fitted map has the
stated equilibrium structure. It does not formally cover M2, whose
prescribed year-by-year catch makes the update a time-indexed family
\(F_t\) rather than a single autonomous map; for M2 the argument is
motivating rather than binding, and the collapse-window result of
Section 3.1 is the relevant evidence. It is not an impossibility
theorem for surplus-production modelling in general.

A third limit is the most consequential for ecological reading, and it
follows from the same algebra. The lower equilibrium is generated by the
removals, not by any depensatory term in the production function. For
the Schaefer map the positive equilibria of \(rS(1-S/K)=C\) are
\(S_{\pm}(C) = (K/2)\left[1 \pm \sqrt{1-4C/(rK)}\right]\), so the
repeller is a function of catch: at the fitted \(r = 1.935\),
\(K = 1032.7\) kt it sits at \(144\) kt under \(C = 240\) kt, at
\(66\) kt under \(C = 120\) kt, and at \(2.6\) kt under the
moratorium-level \(C = 5\) kt, approaching zero as \(C\) does. The
1995 SSB of \(9.7\) kt is far below the first of these and above the
last. The obstruction is therefore a statement about a stock held under
sustained removals of the pre-collapse magnitude; it is not evidence of
biological depensation, and it does not explain non-recovery after
catches fell. Nothing in Definition 2.1's Schaefer branch contains an
Allee effect --- that is the separate M1b branch --- so a threshold
recovered from an M1 fit should not be read as a biological one.

The collapse forecasts fail for a more direct reason than the
equilibrium structure. Every module carries stationary production, so
each predicts positive surplus at the biomass levels of the late 1980s
while the reconstruction falls. For the modules supplied with the catch
path --- M2, M3 and M4 --- the 1992 decline in prescribed catch raises
the predicted increment directly, so they turn upward exactly where the
reconstruction turns down; M1 and M1b use a training-mean constant and
turn upward because that constant is smaller than the losses the
reconstruction records. That mechanism, not the crossing of a
catch-dependent repeller computed at \(C = 240\) kt, is what the
fixed-origin collapse scores record.

The fixed-window sign-hit rates of \(\Delta S\) for the structural
modules (0.00--0.50 on collapse; NA for the baselines, which supply no
directional prediction; Table 3) record the directional character of that failure: the
structural models miss not only the magnitude of the crash but the sign
of the trajectory exactly when the official catch series drops.

Under both the coarse catch regime and official landings, lowering \(C\)
in 1992 cannot generate the crash. Neither prescribed catch treatment reproduced the reconstructed decline
in these implementations; which process accounts for the shortfall ---
productivity, unallocated mortality, or the observation and
reconstruction of the series --- is not identified by this comparison.
Nor does this test whether fishing caused the historical collapse: in
this accounting \(S_t\) is spawning-stock biomass while \(C_t\) is total
landings, which is not an age-structured removal from the spawning
stock, so the exercise bounds a scalar catch-driven surplus story and
not the causal role of the fishery. That is
compatible with DFO's caution that NCAM \(M\) can absorb unreported
deaths: the crash remains such an event under the best available public
catch reconstruction, which tightens that limitation rather than
relaxing it.

Unidentified extra structure does not close the gap to persistence, and
in two of the five structural modules (M2 and M4) it widens it. The direction is not uniform
and should not be stated as though it were: relative to its own declared
comparator M3 reads below M2 in all six rolling comparisons (by 6 to 129
kt), and M1b reads below M1 at \(h=1\) on Specification A and on the
fixed recovery window, though not on Specification B, where M1b is the
higher of the two (152 against 120 kt at \(h=1\)). What fails is H2, not always H1. Where the added
structure does hurt, the mechanism is visible: autoregressive residuals
fitted on short, regime-changing windows persist the wrong sign, an
Allee parameter goes to the boundary, and a one-year delay removes
information. These modules stay descriptive unless the conditions under
which a module's contribution could be certified rather than merely
scored hold. The delay mechanism itself is formalized exactly: the
certification companion's delayed-information theorem bounds guaranteed
recognition time by the pre-observation state (Abaee, 2026, An obstruction calculus), the
theorem-level form of the one-year lead documented here.

The 2016 LRP is the 1980s mean SSB, so over its own reference years the
mean deviation of SSB from that bound is zero by construction, though
individual annual deviations are not. Defining the threshold from the
same years that a leading-indicator claim would be tested on does not
establish prospective early-warning performance for 1983--1990.

The moratorium is a change in implementable catch and is already in
\(C_t\). A separate delay switch does not add a degree of freedom beyond
stock-flow. Whether the 1992 management architecture admits any feasible
catch path that keeps the stock above its reference threshold is a
separate question from one-year forecast error, treated for this stock
in Abaee (2026, Robust viability of the 2J3KL limit reference point), and is not addressed
here.

One-dimensional surplus production is not NCAM: age structure,
migration, and survey catchability are omitted. The test is whether this
ladder is retained under the stated score, not whether the assessment is
a good filter. M4 is a stale-start experiment rather than a fully delayed information
set: its trajectory begins at \(S_{t-1}\), but its parameters are
M3's, fitted on training transitions that end at \(S_t\), so the
later state still reaches the forecast through the fitted coefficients
and the carried residual. Because the path starts one year earlier,
reaching a target \(h\) years after the origin takes \(h+1\) updates
rather than \(h\). The catch path is indexed from the earlier start, so
the first update applies \(C_{t-1}\), and the residual recursion is
applied before the first step, so the first projected residual is
\(\hat\phi\,e_{\mathrm{last}}\) where \(e_{\mathrm{last}}\) is the
residual of the transition \(S_{t-1}\to S_t\); the first update
therefore re-forecasts the transition on which that residual was fitted.
It is therefore a perturbed
initialisation evaluated at calendar time \(t\), not propagation from a
genuinely year-old information set. A genuinely delayed variant would truncate
training as well, and is not scored here. Its one-year gap against
persistence therefore mixes a stale initial state with the model's own
cost.

The separating control --- persistence issued from the last available
assessment, \(S_{t-1}\) --- decomposes the gap. On Specification A under
the coarse catch regime (the primary pass) the control reads 184.4 kt at
\(h = 1\) and 329.8 kt at \(h = 5\), against M4's 195.6 kt and 488.3 kt.
The information delay, the year-old start under the same persistence
rule, therefore accounts for 86.4 of the 97.5-kt one-year gap and the
surplus model's own cost for the remaining 11.1 kt; at \(h = 5\) the
split is 65.1 against 158.4 kt.

On Specification B every term is computed on M4's own origin set, so
that the control, the baseline and the module share origins (\(n=59\)
at \(h=1\), \(n=55\) at \(h=5\)). Matched persistence reads 84.4 and
300.0 kt and the stale-start control 158.0 and 337.4 kt, against M4's
205.7 and 1030.7 kt. The delay accounts for 73.6 of the 121.3-kt
one-year gap and the model's own cost for the remaining 47.7 kt; at
\(h = 5\) the model's own cost dominates (693.3 of 730.8 kt). These are descriptive comparisons between two
forecast rules, not a causal decomposition of the value of assessment
timeliness: the delayed control still reads off the same smoothed
reconstruction, and access to a later reconstruction is not the same
thing as access to a timely real-world assessment. With that reading,
M4's error is numerically dominated by the initialisation component at the
one-year horizon on both specifications, and mostly the surplus model's own structure at the
five-year horizon on Specification B: at \(h = 1\) the surplus model's
own penalty is only about 11 kt --- relative to the stale-start control the gap decomposes numerically into
an initialisation component and a residual-surplus component, an
arithmetic split between forecast rules rather than a causal
attribution. For reference, the one-year lag in the start state accounts for \(86.4\) kt
at \(h = 1\) on Specification A, against structural costs over timely
persistence of \(23\) kt for M1, \(17\) for M1b, \(46\) for M2 and
\(37\) for M3, and M4's own structure-given-delay cost of \(11.1\) kt.
These are differences between forecast rules on a smoothed
reconstruction and are not an estimate of the operational value of
assessment timeliness.

The predictand is an assessment smoother (NCAM/xteNCAM SSB), not a raw
observation: forecasting it is forecasting a filtered state, and
persistence inherits the smoother's autocorrelation --- the comparison
is valid for ``does this ladder track this assessment,'' not for a
vintage or early-warning reading, which would require assessment
vintages at each origin, and none is available. Persistence is
correspondingly not an empty null on this predictand. Spawning-stock
biomass aggregates surviving mature cohorts of a long-lived species, and
the assessment adds further temporal smoothing, so strong one-year
autocorrelation is the expected baseline on biological and observational
grounds alike. Beating it requires useful prediction of subsequent change, whether a
sustained trend or a transition such as the collapse, the onset of
rebuilding, or the post-2015 stall. Persistence is retained as the scoring benchmark the retention rule is
defined against, not as a biological projection, and no demographic
balance is asserted on its behalf at either horizon.
Recreational catch remains incompletely measured.

The log-RMSE scores of Table 4 carry a floor caveat: structural
trajectories absorb at the numerical floor
\(\varepsilon_{\mathrm{log}} = 10^{-3}\) kt rather than at zero (the
trajectory code clips the state to
\([\varepsilon_{\mathrm{log}}, 10^{6}]\) kt), and the reported values
use \(\log\max(\hat S, \varepsilon_{\mathrm{log}})\);
\(\varepsilon_{\mathrm{log}}\) is part of the registered scoring
configuration and is distinct from the process noise \(\varepsilon_t\)
of Definition 2.1. The floor binds often --- for M1 and M1b it binds on a majority of
five-year origins on both specifications --- and the per-origin counts
are given in Section SI-4. The raw-RMSE column is the retention score, so no retention verdict
depends on the floor; it affects the log column only, which is reported
as a secondary diagnostic. Recomputing that column with the floor
lowered to \(10^{-6}\) and \(10^{-9}\) leaves the log-RMSE ordering of
the five modules unchanged at both horizons on both specifications, so
no secondary ranking depends on it either. A constant-productivity map is
misspecified for a series whose assessment attributes the crash to
\(M \approx 2.5\); the test is whether that misspecification still
forecasts SSB better than persistence, and the scored comparison answers
that for this ladder on these series. Sample sizes are small
(\(n=33\) years on Specification A; rolling \(n=25\) at \(h=1\) and
\(n=21\) at \(h=5\)). They suffice to rank models and do not suffice to
certify a small skill difference. The 2023 40\% \(B_{\mathrm{MSY}}\) LRP
is the Specification B setting, and Section 3.3 is that review: the
ladder was run on the xteNCAM series under a fresh
specification-matching review, with the same non-retention verdict.

\textbf{Reconstruction-level corroboration.} An independent,
assessment-level record sharpens the same boundary. Rose (2026) compares
two surplus-production reconstructions of the stock spanning 1983--2023
--- the Rose and Walters (2019) model and the DFO (2024a) assessment
model --- and documents that surplus production and stock growth stalled
after 2015, with some years negative, the stall-point biomass being
controversial between the two reconstructions and well below historical
norms. The recovery-stall window of Specification B (train 1995--2012,
test 2013--2024) overlaps that period (its 2016--2024 portion). A stock
whose surplus production and stock growth stall together, with some
years negative and the stall-point biomass below historical norms ---
the configuration Rose (2026) reports --- is precisely the configuration
a constant-productivity surplus law cannot reproduce, and the ladder's
single-origin comparison on that window is consistent with it on the
rolling score, which is what decides retention. On that single fixed
origin M1 (254 kt) and especially M1b (178 kt) do read below frozen
persistence (262 kt); the retention verdict is nonetheless negative,
because it is decided on rolling-origin RMSE, where persistence is not
beaten at either horizon. The fixed-window result is reported because
selecting favourable episodes would tell a different story from rolling
evaluation, and that contrast is itself part of the finding. The two reconstructions' disagreement
at the stall point is moreover the assessment-level twin of the ladder's
own non-uniqueness: M1b's apparent recovery-window skill is carried by a
carrying capacity extrapolated far beyond the training range, exactly
the extrapolation on which the two reconstructions diverge. Rose (2026)
also documents the limit-reference-point trajectory --- successive
downward revisions to just over 0.3 Mt, then to about 0.25 Mt in 2025
--- which is the norm-field drift that separates the two specifications
evaluated here (884.6 kt versus 276 kt) and the reason their columns are
never pooled.

\textbf{Attribution and the prey modules.} Rose (2026) infers from
structural-equation models that the post-2015 production deficit is
limited by capelin and amplified by harp seal predation, with fishing
almost certainly not the sole cause. That attribution is consistent with
the ladder's two negative findings, and it sharpens both. First, it
matches the catch-treatment null: neither prescribed catch treatment
repairs the collapse, which is what one would expect if the dominant
missing term is not catch. Both treatments could omit relevant
removals, and other structural errors could prevent either from helping,
so the scored result is consistent with that attribution rather than
establishing it. Second, it draws the
identification boundary inside the prey results: capelin limitation can
be a real production driver --- the 1991 acoustic collapse is large
(1985--90 median 3704 kt versus 1991--2022 median 174 kt) --- while the
corresponding forecast modules still fail the stated score. A driver's
reality at the reconstruction level does not make it a usable forecast
module at the surplus-production level: the two-regime and
acoustic-index variants are weakly constrained by their short training
windows and missing index years, and they fail the retention score. No
parameter profile was computed for them, so their non-retention rests
on prediction error rather than on a demonstrated identification
failure.

\textbf{Protocol status.} The evaluation windows and the scoring core
--- the primary RMSE comparison and the persistence bar --- were coded
before the first scoring pass and applied unchanged. Two components of
the retention rule were not: the 5\% tie band and the comparator
declarations for the non-nested rungs were formalised after the scores
of Section 3 had been computed. Neither affects any outcome reported
here; the tie band never decides a retention result, the smallest deficit
being \(17.1\%\). The design is therefore a fixed computational
protocol, not a prospective registration, and the additional passes
(annual landings, survey start, capelin, Specification B) were declared
in the text as they were added. No dated pre-scoring protocol file exists for this study.
Section~SI-1 records the order in which each pass was run.

The results are specific to the objects tested. They do not extend to
other estimators of the same model family, to closed-loop filtering
schemes, or across the two assessment specifications, which are scored
separately throughout. The delay evaluated here is a one-year
information delay in the forecast start state, not a continuous-time
retarded system.

On this evidence, within the scope of this estimator and ladder, modules
that are not identified on the training data did not reduce forecast
error below matched persistence at both horizons here. Weak
identification is documented for the recovery-window Allee fits and is
not established as the sole cause of non-retention. The retained
forecast is persistence.

Sections 3.8--3.10 locate the practical margin. The outcome years contain
no event the ladder could win, the verdict layer moved under the stock
twice in a decade (threshold revision, vintage revision), and the one
genuine decision failure of the period was an observation failure. On
this stock, the binding constraint is the monitoring and reference-point
layer, not forecast structure; the retention rule's empty set is a
statement about structure, and the monitoring companion's design rules
are the instrument matched to where the loss of value actually occurred.

\subsection{5. Conclusions}\label{conclusions}

On locked NCAM M-shift SSB for 1983--2015, last-value persistence is
more accurate, on the pooled rolling score, than surplus production,
catch-driven stock-flow, autoregressive residuals, lagged
initialisation, and prey-informed productivity. Neither prescribed catch
treatment reproduced the reconstructed decline in this accounting,
though which process accounts for the shortfall is not identified here. Extra structure did
not recover the gap to persistence at both horizons, and the
lagged-initialisation module enlarged it; the autoregressive module
improved on its own declared comparator without closing the gap
to persistence.

The strength of that finding is not uniform across modules, and a
pre-registered simulation is what separates the cases. Under known
ground truth the rule recovers a true autonomous module in \(97\%\) of
replicates at collapse-window parameters, so the empirical
non-retention of that module is evidence about the series rather than
about the instrument; for the stock-flow and depensation alternatives at
recovery-window parameters power falls below \(15\%\), so their
non-retention is weak evidence against those structures. The rule is
specific --- it declines to retain structure under a persistence-true
process in \(97\)--\(99\%\) of replicates --- but an information
criterion penalising free parameters would have been the stronger
instrument on both axes. Whether a diagnostic can tell an analyst in
advance where the rule has power remains unresolved. Scored on a
groundwater system in Abaee (2026, Does a one-pool water-balance model improve forecasts of Edwards Aquifer head) and restated under the rule used
here, the same criterion withholds retention from a module that beats
persistence on the point ranking, so the tie band and the two-horizon
requirement are load-bearing rather than nominal. The retained set is
empty on both systems, while the series are not pooled and no verdict is
transferred between them. The same retention outcome holds on the
unpooled xteNCAM series, whose post-2015 stall period fails in the same
way the reconstructions stall.

The paper reports a forecast comparison. It does not conclude that the
stock is unsustainable, and it does not conclude that the evaluation
framework is empirically confirmed on this stock.

The outcome years, the threshold experiment, and the 2022--2023
suspension sharpen the scope statement. The forecast comparison stands as
scored; the stock's decision-relevant uncertainty over 2016--2024
accumulated in the observation and reference-point layer --- which
vintage reads the state, and which threshold reads the verdict --- and in
a two-cycle suspension of the observing system itself.

\subsection{References}\label{references}

Abaee, A., 2026. Does a one-pool water-balance model improve forecasts
of Edwards Aquifer head? A scored test at J-17. Zenodo.
https://doi.org/10.5281/zenodo.22552680.

Abaee, A., 2026. Periodic review as sampled governance: sample-and-hold
dynamics of assessment-driven effort control, a selected 42-stock
spectral screen, and the Northern Cod case. Zenodo.
https://doi.org/10.5281/zenodo.22554297.

Abaee, A., 2026. Robust viability of the 2J3KL limit reference point
under a surplus-production map: policy scoring, expansion, and when
catch cannot help. Zenodo. https://doi.org/10.5281/zenodo.22552060.

Abaee, A., 2026. An obstruction calculus for viability under incomplete
observation. Manuscript submitted for publication.

Cadigan, N.G., 2016. A state-space stock assessment model for northern
cod, including under-reported catches and variable natural mortality
rates. Can. J. Fish. Aquat. Sci. 73, 296--308.

Diebold, F.X., Mariano, R.S., 1995. Comparing predictive accuracy. J.
Bus. Econ. Stat. 13, 253--263.
https://doi.org/10.1080/07350015.1995.10524599

Carvalho, F., Winker, H., Courtney, D., Kapur, M., Kell, L., Cardinale,
M., Schirripa, M., Kitakado, T., Yemane, D., Piner, K.R., Maunder, M.N.,
Taylor, I., Wetzel, C.R., Doering, K., Johnson, K.F., Methot, R.D.,
2021. A cookbook for using model diagnostics in integrated stock
assessments. Fish. Res. 240, 105959.

DFO, 2009. A fishery decision-making framework incorporating the
Precautionary Approach. Fisheries and Oceans Canada, Ottawa.

DFO, 2016. Stock Assessment of Northern cod (NAFO Divs. 2J3KL) in 2016.
DFO Can. Sci. Advis. Sec. Sci. Advis. Rep.~2016/026.

DFO, 2024a. NAFO Divisions 2J3KL Northern Cod stock assessment to 2024.
DFO Can. Sci. Advis. Sec. Sci. Advis. Rep.~2024/049.

DFO, 2024b. Assessment of capelin in NAFO Divisions 2J3KL. DFO Can. Sci.
Advis. Sec. Sci. Advis. Rep.~2024/050.

Hutchings, J.A., Myers, R.A., 1994. What can be learned from the
collapse of a renewable resource? Atlantic cod, Gadus morhua, of
Newfoundland and Labrador. Can. J. Fish. Aquat. Sci. 51, 2126--2146.

Hyndman, R.J., Koehler, A.B., 2006. Another look at measures of forecast
accuracy. Int. J. Forecast. 22, 679--688.

Kell, L.T., Kimoto, A., Kitakado, T., 2016. Evaluation of the prediction
skill of stock assessment using hindcasting. Fish. Res. 183, 119--127.

Kell, L.T., Sharma, R., Kitakado, T., Winker, H., Mosqueira, I.,
Cardinale, M., Fu, D., 2021. Validation of stock assessment methods: is
it me or my model talking? ICES J. Mar. Sci. 78, 2244--2255.

Kokkalis, A., Berg, C.W., Kapur, M.S., Winker, H., Jacobsen, N.S.,
Taylor, M.H., Ichinokawa, M., Miyagawa, M., Medeiros-Leal, W., Nielsen,
J.R., Mildenberger, T.K., 2024. Good practices for surplus production
models. Fish. Res. 275, 107010.

K\"unsch, H.R., 1989. The jackknife and the bootstrap for general
stationary observations. Ann. Stat. 17, 1217--1241.
https://doi.org/10.1214/aos/1176347265

Murphy, H.M., Adamack, A.T., Lewis, R.S., Bourne, C.M., 2025. Assessment
of capelin in NAFO Divisions 2J+3KL to 2023. DFO Can. Sci. Advis. Sec.
Res. Doc. 2025/022.

Northwest Atlantic Fisheries Centre, 2025. 2J3KL cod and capelin biomass
indices. Zenodo. https://doi.org/10.5281/zenodo.17515115

Pella, J.J., Tomlinson, P.K., 1969. A generalized stock production
model. Bull. Inter-Am. Trop. Tuna Comm. 13, 419--496.

Regular, P.M., et al., 2025. Assessment of the Northern Cod stock in
NAFO Divisions 2J3KL in 2024. DFO Can. Sci. Advis. Sec. Res. Doc.
2025/048.

Rose, G.A., 2026. Northern cod comeback: 10 years after. Can. J. Fish.
Aquat. Sci. 83, 1--14. https://doi.org/10.1139/cjfas-2025-0141

Rose, G.A., Rowe, S., 2015. Northern cod comeback. Can. J. Fish. Aquat.
Sci. 72, 1789--1798.

Rose, G.A., Walters, C.J., 2019. The state of Canada's iconic Northern
cod: a second opinion. Fish. Res. 219, 105314.

Schaefer, M.B., 1954. Some aspects of the dynamics of populations
important to the management of the commercial marine fisheries. Bull.
Inter-Am. Trop. Tuna Comm. 1(2), 27--56.

Schijns, R., Froese, R., Hutchings, J.A., Pauly, D., 2021. Five
centuries of cod catches in Eastern Canada. ICES J. Mar.~Sci. 78,
2675--2683.

Shelton, P.A., Healey, B.P., 1999. Should depensation be dismissed as a
possible explanation for the lack of recovery of the northern cod (Gadus
morhua) stock? Can. J. Fish. Aquat. Sci. 56, 1521--1524.

Walters, C., Maguire, J.-J., 1996. Lessons for stock assessment from the
northern cod collapse. Rev.~Fish Biol. Fish. 6, 125--137.

\section*{Declarations}

\subsection*{Data availability}

All input data, analysis scripts, result files, and the frozen model
specification are archived at
\url{https://github.com/MIKEAA2020/general-sustainability} and
\url{https://zenodo.org/records/22553609}. The outcome-year, vintage,
and threshold analyses of Sections 3.8--3.10 use the RAM Legacy Stock
Assessment Database v4.66 (10.5281/zenodo.14043038; stock COD2J3KL), the
digitized NCAM Table A2 and xteNCAM Table 17 files, and the capelin
acoustic index; these are archived in-repo under
\texttt{paperE1\_calibration\_data\_v1\_wave\_e\_cod/} with locked
provenance, together with the calibration record
\texttt{paperE1\_calibration\_data\_v1.md}.

\subsection*{Reproducibility}

All computations are deterministic: within a fixed interpreter and
library stack, repeated executions reproduce every result file byte for
byte. An independent re-execution verified all 29 recorded output
checksums and regenerated all 30 result files byte-identically; one
registered file carries no pinned checksum and is covered by the
regeneration comparison. Across environments, the intervention runners
and the OLS-based prediction runners regenerate their outputs byte for
byte, while the four optimizer-based forecast runners (L-BFGS-B fits)
reproduce every scored row at printed precision with one exception: the
M1b rolling row on Specification B, whose Allee optimum is
environment-sensitive at the \ensuremath{\pm}17 kt level (\(h=1\): 151.6
versus 153.2 kt; \(h=5\): 445.5 versus 462.5 kt, on Python
3.12/numpy 2.1.3/scipy 1.14.1 against the recorded original
environment). That sensitivity is consistent with the identification
fragility reported for M1b in Section 3.1 (\(\mathfrak s\) descending
towards zero, \(K\) resting at the 500-kt multi-start initialiser), and
it is immaterial to the retention verdict, which persistence wins at
both horizons by margins far larger than the instability. The archive
records the checksums of the original-environment outputs; the
deterministic-in-environment claim and this sensitivity note together
are the reproducibility statement. Sections 3.8--3.10 are produced by
\texttt{paperE1\_cod\_forecast\_ladder\_v59\_verification.py}
(standard library only), which regenerates every number in those
sections from the archived data files and asserts the cross-checks
recorded there.

The uncertainty layer of Section 3.5 is produced by a seeded,
deterministic script applying Diebold--Mariano (HAC) statistics and a
K\"unsch moving-block bootstrap to the archived per-origin forecast
files; the persistence baseline is recomputed there on the identical
origin sets from the registered series and asserted against the recorded
origin-matched values. The script and its output are archived with the
repository.

The section scripts are archived at \texttt{wave\_e\_cod/src/} in the
repository:
\begin{verbatim}
python3 wave_e_cod/src/run_ladder.py
python3 wave_e_cod/src/run_xte.py
python3 wave_e_cod/src/run_capelin_regime.py
python3 wave_e_cod/src/run_capelin_index.py
python3 wave_e_cod/src/compare_catch.py
python3 wave_e_cod/src/make_figures.py
\end{verbatim}

\subsection*{CRediT authorship contribution statement}

Amin Abaee: Conceptualization, Methodology, Software,
Validation, Formal analysis, Investigation, Data curation, Writing --
original draft, Writing -- review \& editing, Visualization.

\subsection*{Funding}

Amin Abaee: Conceptualization, Methodology, Software,
Validation, Formal analysis, Investigation, Data curation, Writing --
original draft, Writing -- review \& editing, Visualization.

\subsection*{Declaration of competing interest}

The authors declare that they have no known competing financial
interests or personal relationships that could have appeared to
influence the work reported in this paper.

\subsection*{Code availability}

The scored ladder and every auxiliary campaign are produced by the scripts
named above, archived with their outputs. The verification script
\texttt{paperE1\_cod\_forecast\_ladder\_v60\_verification.py} re-runs
the record and recompares the printed scores, reference points and
window facts against the locked input files; it is archived alongside
this manuscript.

\subsection*{AI declaration}
GLM (Z.ai), Qwen (Alibaba Cloud) and DeepSeek AI assisted with drafting
and iterative review. The author reviewed and edited outputs and takes
responsibility for the final work.

\end{document}


% ===== PART: E3_v17  (source: paperE3_edwards_forecast_ladder_v17.tex) =====
\clearpage

% LaTeX source generated from paperE3_edwards_forecast_ladder_v16.md by wave13/build_latex_v13.py (batch 7 (audits of agent arena 1 paper rewrites)/wave13).
% Pandoc body identical to the wave-9/11 output (asserted at build time), with the
% declaration subsections relocated to a trailing Declarations section.
% Front matter: Amin Abaee, Independent Researcher, clickable ORCID
%(https://orcid.org/0000-0002-0019-1842) and email (amin_abaee@ut.ac.ir)
% beneath the name; date September 6, 2026.
% Back matter: the paper's own declarations, each a titled subsection, plus
% the AI declaration (GLM (Z.ai), Qwen (Alibaba Cloud) and DeepSeek AI
% assisted with drafting and iterative review).
% References carry the owner's Zenodo DOIs (see wave13/apply_md_doi.py).
% Edit the markdown source and re-run the build script to regenerate.
% Compiles error-free with tectonic (also compatible with pdflatex/xelatex).
\documentclass[11pt]{article}
\usepackage[a4paper,margin=1in]{geometry}
\usepackage{amsmath,amssymb}
\usepackage{graphicx}
\usepackage{booktabs,longtable,array,calc}
\usepackage{xcolor}
\usepackage[colorlinks=true,allcolors=blue!45!black]{hyperref}
\usepackage[font=small,labelfont=bf]{caption}
\usepackage{etoolbox}
\AtBeginEnvironment{longtable}{\small}  % dense data tables
\setlength{\tabcolsep}{4pt}
\providecommand{\tightlist}{\setlength{\itemsep}{0pt}\setlength{\parskip}{0pt}}
\setcounter{secnumdepth}{-1}
\emergencystretch=3em
\graphicspath{{../}}
\begin{document}
\title{Does a one-pool water-balance model improve forecasts of Edwards Aquifer head? A scored test at J-17}
\author{Amin Abaee\\[0.35em]
{\small Independent Researcher}\\[0.55em]
{\small\href{https://orcid.org/0000-0002-0019-1842}{ORCID: 0000-0002-0019-1842}}\\[0.3em]
{\small\href{mailto:amin\_abaee@ut.ac.ir}{amin\_abaee@ut.ac.ir}}}
\date{September 6, 2026}
\maketitle

{\centering\small \textbf{Prepared in the format of Groundwater (Wiley/NGWA)}\par}
{\centering\small \textbf{Article type: Research Paper}\par}

\begin{abstract}


\textbf{Problem.} Index-well head forecasting is a recurring operational need, while deliberately simple process-based water-balance models are rarely tested against naive benchmarks under a locked retention rule.

\textbf{Approach.} On the 1934--2023 annual-mean J-17 Edwards Aquifer head series, a fixed ladder was scored by rolling and fixed-window out-of-sample RMSE: persistence, training-mean climatology, univariate AR(1), a one-pool stock-flow water balance, residual/delay variants, and climate-informed recharge. A causal module was retained only if it beat both persistence and the next-simpler causal model under a protocol frozen before any score.

\textbf{Results.} The stock-flow map that persists last year's recharge loses to persistence at one year (14.70 vs. 13.23 ft), because annual recharge is near-white (corr(R\_t, R\_{t\ensuremath{-}1}) = 0.17 vs. corr(\ensuremath{\Delta}H\_t, R\_t) = 0.74). The univariate AR(1) improves one-year RMSE by 0.39 ft (12.84 ft), but its margin is a coin-flip (mean absolute error a tie, five-year loss, bootstrap interval covering zero). The climatological-flux affine map (M2m) is the best one-step forecaster (12.28 ft) yet is declined by a protocol clause. Given realized fluxes the same map nowcasts (7.55 ft); climate modules gain at most 0.13 ft and none is retained. At five years climatology beats persistence robustly (16.80 vs. 21.11 ft).

\textbf{Implications.} For annual J-17 head forecasts persistence and an AR(1) are hard to beat; the one-pool balance, given the year's fluxes, nowcasts rather than forecasts; multi-year planning should prefer climatology to persisted recharge.

\end{abstract}

\textbf{Keywords:} Edwards Aquifer; J-17 index well; groundwater level
forecasting; water balance; forecast evaluation; persistence benchmark;
prediction skill

\textbf{Article Impact Statement.} Deliberately simple process-based
water-balance models are rarely tested against naive benchmarks under a
locked retention rule. Across groundwater-forecast evaluation, adding
structure should be justified against the forecast that nothing changes.
On the 1934--2023 J-17 record, the one-pool balance that persists last
year's recharge loses to persistence at the annual origin (recharge is
near-white, r = 0.17), while with realized fluxes the same map nowcasts
(7.6 ft) rather than forecasts. The transferable lesson is a
benchmark-discipline one: climate-independent baseline models are hard
to beat at the annual origin in rapidly recharged, institutionally
bounded aquifers, and at multi-year horizons the training mean (climatology)
is the robust default.

\subsection{1. Introduction}\label{introduction}

Forecasting aquifer heads at index wells is a recurring operational
problem of groundwater management. Drought-stage declarations,
springflow protection, and permit adjustments all key on forecasted
water levels. The methodological literature has responded with an
increasingly rich inventory of groundwater-level forecasting models,
from conceptual water balances to data-driven regressions and,
prominently, artificial-neural-network and wavelet--neural-network
conjunction models trained on measured heads (Daliakopoulos, Coulibaly,
and Tsanis 2005; Adamowski and Chan 2011). The field has also begun to
institutionalize benchmark culture. GEMS-GER, the first machine-learning
benchmark dataset for long-term groundwater levels, standardizes 32
years of weekly observations from 3,207 German wells together with three
benchmark models of increasing complexity (Ohmer et al.~2026).
Systematic comparisons of nine machine-learning and deep-learning
architectures on a karst catchment now anchor the karst groundwater
forecasting literature (Zhu et al.~2026).

Three hydrological objects are worth separating, because they are read
off the same number. The head record is observed at an access point, the
index well. That head indexes a store of water in the aquifer, which is
the resource. The store is replenished and drawn by fluxes, which are
the flow. The one-pool water balance (a lumped stock-flow model with
head change equal to weighted fluxes plus a linear drain) closes this
loop approximately. The term ``water balance'' is used throughout in its
increment-structure sense: head change equals weighted fluxes plus a
linear drain, with the spring-discharge series stored but deliberately
excluded from the forecasting equation, and a {[}610, 710{]} clip
standing in for a physical storage floor. The map is not a closed mass
balance; this qualification is stated here once and travels with the
term.

Whether added structure improves out-of-sample forecasts --- as opposed
to in-sample fit --- is a separate, testable question. The general
forecasting literature has made the benchmark discipline explicit:
across the M4 competition's 100,000 series, sophisticated methods did
not uniformly beat simple statistical baselines (Makridakis, Spiliotis,
and Assimakopoulos 2020). The minimal bar for any proposed module is
therefore the forecast that nothing changes --- last-value persistence
--- together with the training mean. Yet the groundwater benchmark
studies above compare model families against one another. They do not
subject each module to a retention gate against the naive baselines, and
none subjects a deliberately simple process-based water balance to a
scored ablation against those baselines. This paper supplies that
missing test. It applies a scored model-ablation design in which a
scored ladder (a forward-ordered set of models evaluated by a fixed
retention rule) of incrementally structured one-pool models is evaluated
against the two naive baselines. Complexity is kept only if it improves
the stated score, decided on out-of-sample error on the predictand
itself.

The San Antonio Pool of the Edwards (Balcones Fault Zone) Aquifer,
Texas, indexed by the J-17 well, is the natural test bed. Its 1934--2023
daily head record is among the longest managed groundwater series in
North America. Its institutional thresholds (the 660-ft Stage I line of
the Edwards Aquifer Authority) and its physical threshold (the \ensuremath{\approx}618-ft
level at which Comal Springs approaches cessation) are explicit and
dated. Recharge and pumpage series exist that are constructed
independently of the head series. The aquifer is also a karst system in
which regional flow has long been represented --- and debated ---
through equivalent porous media and lumped approaches (Scanlon et
al.~2003, for the Barton Springs segment; the lumped-versus-EPM question
is inherited by the San Antonio Pool, not settled by that citation). For
this reason, the fate of a deliberately simple one-pool water-balance
map is a live hydrogeologic question rather than a straw man.

The question is whether stock-flow, residual, delay, or climate-informed
recharge modules reduce forecast error of J-17 annual-mean head relative
to last-value persistence and to a univariate AR(1). The climate modules
enter because El Ni\~no/Southern Oscillation (ENSO) phases shift North
American precipitation patterns (Ropelewski and Halpert 1986):
September--November Ni\~no 3.4 and lagged climate-division precipitation
are precisely the signals observable at an annual forecast origin.

A companion study under separate review applies the same scored design
to a marine fishery stock (Northern cod, NAFO 2J3KL); the analogy
between the two systems is that the identified driver is not persistent
at the forecast origin (Author et al., in review). The two systems'
scores are never pooled, and no retention verdict is transferred between
them. This paper does not assess whether the aquifer is sustainable and
does not treat springflow or reconstructed storage as co-primary
predictands; a two-pool exchange module was specified and not fitted.

\subsection{2. Data and Specification}\label{data-and-specification}

The specification separates data fields from empirical constructs,
models, and normative thresholds, so that the scored object is
unambiguous.

\textbf{Table 1.} The San Antonio Pool specification: J-17 annual-mean
elevation, 1934--2023.

\begin{longtable}[]{@{}
  >{\raggedright\arraybackslash}p{(\columnwidth - 4\tabcolsep) * \real{0.3333}}
  >{\raggedright\arraybackslash}p{(\columnwidth - 4\tabcolsep) * \real{0.3333}}
  >{\raggedright\arraybackslash}p{(\columnwidth - 4\tabcolsep) * \real{0.3333}}@{}}
\toprule\noalign{}
\begin{minipage}[b]{\linewidth}\raggedright
Field
\end{minipage} & \begin{minipage}[b]{\linewidth}\raggedright
Contents
\end{minipage} & \begin{minipage}[b]{\linewidth}\raggedright
Type
\end{minipage} \\
\midrule\noalign{}
\endhead
\bottomrule\noalign{}
\endlastfoot
System & Edwards Aquifer, San Antonio Pool, indexed by J-17 & D \\
\(H_t\) & Calendar-year mean of daily-high J-17 elevation (ft AMSL) &
D \\
Well & TWDB 6837203 / EAA AY-68-37-203 & D \\
Domain & San Antonio Pool; calendar years 1934--2023 & D \\
Physical threshold & Head high enough that Comal Springs does not cease
(\ensuremath{\approx}618 ft; 1956 daily minimum 612.51 ft) & N / E \\
Institutional threshold & EAA Stage I, 10-day mean J-17 \textless{} 660
ft (in force after 2007) & N \\
Disturbance & Recharge pulses; unmodeled karst; Uvalde--San Antonio
exchange & M \\
Implementable use & Pre-EAA pumping; post-1996/2007 permits and
critical-period management & E \\
Service series & USGS 08168710 Comal Springs annual mean discharge (cfs)
& D \\
\end{longtable}

Field types: D = data, E = empirical construct, M = model, N = normative
threshold; TWDB = Texas Water Development Board; NSE = Nash--Sutcliffe
efficiency. The predictand is a measured well series. It is not a GRACE
or G3P storage reconstruction, a MODFLOW or GWSIM inversion, EAA
reconstructed storage, the J-27 Uvalde index, San Marcos Springs, or
total spring discharge. The predictand is the calendar-year mean of
daily \emph{highs}; a daily-mean convention would sit slightly lower and
matters near the 660/618-ft lines, and the Authority's Stage I rule is a
10-day mean --- the annual-mean proxy is flagged where used.

Stage I at 660 ft is a 2007 institutional rule and is not applied as if
it existed in 1956. Cessation of Comal Springs near 618 ft is a physical
service bound. The indicator 1\{H \textless{} 660\} on the annual mean
is a coarse proxy for the Authority's 10-day declaration, not the
declaration itself.

Annual means use all available daily highs. Years with fewer than 240
observations are dropped. No year falls below the floor (minimum n =
242, 1939), so the rule never binds on this panel; 1935 (n = 258) and
1939 (n = 242) are retained as incomplete-coverage means; missing days
are not interpolated. The published pre-continuous composite (Beverly
Lodges extension) is used as issued. The complete estimation panel ends
in 2023, whose values carry provisional TWDB status.

Recharge R is USGS total San Antonio-area recharge (10\textsuperscript{3} acre-ft yr\textsuperscript{-1}),
estimated by the Puente (1978) stream-loss method and constructed
independently of J-17 (Umphres and Choi 2025). Pumpage P is well
discharge from Edwards Aquifer Authority Table 1 (Edwards Aquifer
Authority 2024/25, covering 1934 onward). R and P are constructed
fluxes, not observations of J-17. Total spring discharge is stored and
is not used as a driver. Climate predictors, used only in Section 5.4,
are September--November Ni\~no 3.4 (HadISST, anomaly relative to
1991--2020) and calendar-year precipitation in Texas climate divisions
06 and 07 (NCEI nClimDiv).

\subsection{3. Forecast Models}\label{forecast-models}

The forecast models form a fixed ladder. Each rung adds structure to the
previous rung, and each causal rung is scored against the next-simpler
rung under the retention rule of Section 4.

\textbf{Definition 3.1 (One-pool water-balance map).} With head clipped
to {[}610, 710{]} ft, the one-pool map is

\[H_{t+1}=\bigl[H_t+\alpha+\beta\tilde R_{t+1}+\gamma\tilde P_{t+1}+\delta H_t\bigr]_{\mathrm{clip}},\]

where \(\tilde R\) and \(\tilde P\) are the flux values a forecast
substitutes for the (unknown) year \(t+1\) fluxes. Forecasts are issued
at the end of year \(t\) from (H, R, P) through \(t\).

\textbf{Table 2.} Model ladder.

\begin{longtable}[]{@{}
  >{\raggedright\arraybackslash}p{(\columnwidth - 6\tabcolsep) * \real{0.2500}}
  >{\raggedright\arraybackslash}p{(\columnwidth - 6\tabcolsep) * \real{0.2500}}
  >{\raggedright\arraybackslash}p{(\columnwidth - 6\tabcolsep) * \real{0.2500}}
  >{\raggedright\arraybackslash}p{(\columnwidth - 6\tabcolsep) * \real{0.2500}}@{}}
\toprule\noalign{}
\begin{minipage}[b]{\linewidth}\raggedright
ID
\end{minipage} & \begin{minipage}[b]{\linewidth}\raggedright
Class
\end{minipage} & \begin{minipage}[b]{\linewidth}\raggedright
Fluxes at \(t+h\)
\end{minipage} & \begin{minipage}[b]{\linewidth}\raggedright
Role
\end{minipage} \\
\midrule\noalign{}
\endhead
\bottomrule\noalign{}
\endlastfoot
persist & baseline & --- & \(\hat H_{t+h}=H_t\) \\
mean & baseline & --- & training-window mean \\
M1 & autonomous & --- & \(H_{t+1}=a+\varphi H_t\) \\
M2 & causal stock-flow & last \((R_t,P_t)\) persisted & candidate for
retention \\
M2m & climatological stock-flow & training-mean \((R,P)\) & affine
AR(1) \\
M3 & residual & as M2 & AR(1) residual on \(\Delta H\) \\
M4 & delay & as M2 & symmetry control with the companion fisheries
evaluation; not a monitoring constraint of this system \\
M2\_oracle & diagnostic & realized future \(R,P\) & excluded from
retention \\
\end{longtable}

\textbf{Remark 3.1 (Class reduction of M2m).} Under constant fluxes
\(\tilde R_{t+1}=\bar R\), \(\tilde P_{t+1}=\bar P\), M2m reduces to
\(H_{t+1}=(1+\delta)H_t+\mathrm{const}\) and therefore shares M1's
forecast function class --- but not its estimator. M2m pins its
intercept from the in-sample mean fluxes, an additional identifying use
of the recharge and pumpage records in training; the persistence
coefficient \((1+\delta)\) is not mean-flux dependent. The decline of
M2m is a protocol choice recorded in the frozen documents, not a
theorem, and it is not part of the retention rule itself (Section 4.1).

Climate-informed variants (Section 5.4) replace persisted R by a
one-step forecast of R from information known at t: AR(1) on R, lagged
Ni\~no 3.4, lagged climate-division precipitation, or all three. A
precipitation-oracle variant uses year \(t+h\) precipitation and is
excluded from retention. For horizons \(h>1\), the one-step recharge
forecast is held constant.

\subsection{4. Evaluation Design}\label{evaluation-design}

The evaluation design fixes the scoring rule before any score is
computed. The retention rule (a module is kept only if it beats both
persistence and the next-simpler causal model) is the gate against which
every causal rung of the ladder is judged.

\textbf{Definition 4.1 (Primary and secondary scores).} The primary
score is RMSE of annual-mean J-17, in feet. Secondary scores are mean
absolute error and the Brier score for \(\mathbf{1}\{\hat H < 660\}\),
interpreted only for origins at or after 2007; for deterministic 0/1
forecasts this score is a misclassification rate. A sign-hit rate of
\(\Delta H\) was listed in an earlier design iteration and is struck: no
value is reported for it, and it is not part of any retention decision.

\textbf{Definition 4.2 (Retention rule, frozen verbatim).} Let
\(\mathrm{RMSE}(\cdot)\) denote the primary rolling-origin RMSE at a
fixed horizon \(h\). A causal module \(M\) is retained only if:

(H1) \(\mathrm{RMSE}(M) < \mathrm{RMSE}(\text{persist})\), i.e.~it beats
last-value persistence, and

(H2)
\(\mathrm{RMSE}(M) < \mathrm{RMSE}(M_{\text{next-simpler causal}})\),
i.e.~it beats the next-simpler causal model.

Retention is decided at \(h=1\) on the point-RMSE comparison; no tie
tolerance and no minimum margin is imposed. Diagnostic oracles are
excluded from retention. The Comal series is excluded from retention.
Retention is decided on the head series only; the service series is
scored after that decision is frozen.

\subsubsection{4.1 Protocol record and
deviations}\label{protocol-record-and-deviations}

The scoring protocols for the primary pass and the climate pass were
frozen and dated (2026-08-25) before the corresponding RMSE tables were
computed, and are archived with the analysis code in the public
repository. The design is a fixed computational protocol rather than a
prospective clinical-style registration; the phrase ``pre-registered''
is avoided for that reason. The frozen retention rule is Definition 4.2,
verbatim above. Three protocol elements sit outside that rule and are
therefore recorded here as deviations, in one place:

\begin{enumerate}
\def\labelenumi{\arabic{enumi}.}
\tightlist
\item
  \textbf{The M2m class clause.} M2m is listed and then declined on
  class grounds (Remark 3.1). The clause is in the frozen protocol
  document, but it is not in Definition 4.2; under the rule as written,
  M2m satisfies (H1) and (H2) at h = 1.
\item
  \textbf{The climate-rung comparator.} The frozen Pass-2 protocol
  document states the climate question as whether a causal recharge
  forecast reduces primary RMSE on J-17 ``relative to persistence and
  relative to M1'' --- the retained M1. The comparator is read
  accordingly: persistence and M1. Declaring the declined M2m as the
  rung's (H2) comparator instead would be inconsistent with the frozen
  document and circular --- the gate would be a model the protocol had
  itself declined. Section 5.4 reports the climate verdict on the
  declared gate, and the M2m margins are retained as a nested-baseline
  reading rather than as the gate. The choice changes no frozen verdict
  --- no climate module is retained under either statement --- but it
  fixes the stated mechanism.
\item
  \textbf{The struck sign-hit score} (Definition 4.1).
\end{enumerate}

Post-freeze objects, labelled as such: the climate-pass fixed-window
scores, the pumpage counterfactuals of Section 5.6, the Comal
service-series scoring, and the uncertainty layer of Section 5.3.1
together with its independent replication. None replaces or alters a
frozen verdict.

Fixed windows: (1) drought-of-record drawdown, train 1934--1950, test
1951--1956; (2) drought-of-record recovery, train 1934--1956, test
1957--1961; (3) pre-permit wet interval, train 1980--1990, test
1991--1995; (4) critical-period era, train 1997--2014, test 2015--2023.

Rolling origin: minimum 15 training years; horizons h = 1 and h = 5; n =
75 and n = 71 origins respectively (every model in Tables 4 and 7 is
scored on the identical origin sets; M2m uses the same n = 75 / n = 71).
The fixed windows train on 16, 22, 10, and 17 transitions respectively
and fit four parameters; early rolling origins with the 15-year floor
give 14 transitions, so the causal family's rolling scores carry
unstable early-window coefficients (Section 5.2).

\subsection{5. Results}\label{results}

\subsubsection{5.1 The series}\label{the-series}

\begin{figure}[htbp]
\centering
\includegraphics[width=\linewidth]{figs_e3/fig1_series.png}
\caption{J-17 annual mean and daily-high range. 1956 mean 623.15 ft, daily minimum 612.51 ft. 1992 mean 691.96 ft, daily maximum 703.31 ft. 2023 mean 635.68 ft. The annual mean is below 660 ft in 31 of 90 years; the daily minimum is below 618 ft in one year (1956).}
\end{figure}

\subsubsection{5.2 Fixed windows}\label{fixed-windows}

\begin{figure}[htbp]
\centering
\includegraphics[width=\linewidth]{figs_e3/fig2_windows.png}
\caption{Multi-step forecasts from the end of each training window. The oracle (dashed) is diagnostic and uses realized future R, P.}
\end{figure}

\textbf{Table 3.} Fixed-window RMSE (ft). Bold marks the lowest RMSE of
the window including the diagnostic oracle.

\begin{longtable}[]{@{}
  >{\raggedright\arraybackslash}p{(\columnwidth - 16\tabcolsep) * \real{0.0857}}
  >{\raggedleft\arraybackslash}p{(\columnwidth - 16\tabcolsep) * \real{0.1143}}
  >{\raggedleft\arraybackslash}p{(\columnwidth - 16\tabcolsep) * \real{0.1143}}
  >{\raggedleft\arraybackslash}p{(\columnwidth - 16\tabcolsep) * \real{0.1143}}
  >{\raggedleft\arraybackslash}p{(\columnwidth - 16\tabcolsep) * \real{0.1143}}
  >{\raggedleft\arraybackslash}p{(\columnwidth - 16\tabcolsep) * \real{0.1143}}
  >{\raggedleft\arraybackslash}p{(\columnwidth - 16\tabcolsep) * \real{0.1143}}
  >{\raggedleft\arraybackslash}p{(\columnwidth - 16\tabcolsep) * \real{0.1143}}
  >{\raggedleft\arraybackslash}p{(\columnwidth - 16\tabcolsep) * \real{0.1143}}@{}}
\toprule\noalign{}
\begin{minipage}[b]{\linewidth}\raggedright
Window
\end{minipage} & \begin{minipage}[b]{\linewidth}\raggedleft
persist
\end{minipage} & \begin{minipage}[b]{\linewidth}\raggedleft
mean
\end{minipage} & \begin{minipage}[b]{\linewidth}\raggedleft
M1
\end{minipage} & \begin{minipage}[b]{\linewidth}\raggedleft
M2
\end{minipage} & \begin{minipage}[b]{\linewidth}\raggedleft
M2m
\end{minipage} & \begin{minipage}[b]{\linewidth}\raggedleft
M3
\end{minipage} & \begin{minipage}[b]{\linewidth}\raggedleft
M4
\end{minipage} & \begin{minipage}[b]{\linewidth}\raggedleft
oracle
\end{minipage} \\
\midrule\noalign{}
\endhead
\bottomrule\noalign{}
\endlastfoot
Drawdown 1951--56 & 23.75 & 35.19 & 30.94 & \textbf{18.11} & 27.44 &
18.12 & 18.23 & 19.69 \\
Recovery 1957--61 & 43.62 & 14.07 & 56.24 & 55.32 & 37.74 & 55.28 &
55.12 & \textbf{12.26} \\
Pre-permit wet 1991--95 & 30.13 & 18.24 & 20.02 & 16.67 & 23.47 & 16.41
& 15.26 & \textbf{7.18} \\
Critical-period era 2015--23 & 27.41 & 14.77 & 15.62 & 23.37 & 14.79 &
22.84 & 22.17 & \textbf{8.70} \\
\end{longtable}

Fixed windows are diagnostic, not retention. On the drawdown window the
bolded M2 is a continuing-drought artefact, not a forecasting merit: the
last observed R is already low, and the map trained on 1934--1950 has
the wrong sign on pumpage (\ensuremath{\gamma} = +0.021) --- pumping rose as the drought
deepened, so least squares aliases the human response into the state
transition (simultaneity bias; the short window compounds the
identification failure). That identification failure is why M2 ``beats''
the oracle there. The train-mean baseline is the best non-oracle
forecast on the recovery window (14.07 ft) and the critical-period era
(14.77 ft), consistent with its rolling five-year win; the
residual-persistence rungs M3/M4 track M2 on the drawdown (18.12/18.23
versus 18.11 ft); M4 is the best causal model on the pre-permit wet
window (15.26 ft); and both fail with the causal family on the recovery
and critical-period windows. The recovery window is also where the
{[}610, 710{]} clip binds on forecast paths: the M2 trajectory with
persisted (R, P) = (43.7, 321.1) runs 617.1 \ensuremath{\rightarrow} 611.6 \ensuremath{\rightarrow} 606.8 \ensuremath{\rightarrow} 605.3 \ensuremath{\rightarrow}
605.3 ft, reaching the 610-ft floor from 1959 onward. The clip never
binds on the observed record (annual means 623--692 ft); it binds only
on these persisted-drought forecast paths, and it is doing work exactly
there.

The 1957--61 recovery is a recharge pulse (R\textsubscript{1956} = 43.7, R\textsubscript{1957} = 1142.6
\ensuremath{\times} 10\textsuperscript{3} acre-ft). Persistence stays at the 1956 floor (RMSE 44 ft). Causal
M2 persists drought recharge and falls further (RMSE 55 ft, into the
clip). The oracle, given the 1957 flood year, tracks the rise (RMSE 12
ft). Recoveries on this specification are recharge events, not
autonomous mean reversion --- although the training mean (mean
reversion) scores 14.07 ft, second only to the oracle, in this
institutionally bounded, rapidly recharged system. The 1992 peak and the
2015--23 window repeat the same split: the oracle follows the recharge
year; persistence and persisted recharge do not.

\subsubsection{5.3 Rolling origin}\label{rolling-origin}

\begin{figure}[htbp]
\centering
\includegraphics[width=\linewidth]{figs_e3/fig3_rmse.png}
\caption{Rolling-origin RMSE. The oracle is excluded from retention.}
\end{figure}

\textbf{Table 4.} Rolling-origin summary (ft).

\begin{longtable}[]{@{}lrrr@{}}
\toprule\noalign{}
Model & h=1 RMSE & h=1 MAE & h=5 RMSE \\
\midrule\noalign{}
\endhead
\bottomrule\noalign{}
\endlastfoot
persist & 13.23 & 10.73 & 21.11 \\
mean & 16.17 & 13.17 & \textbf{16.80} \\
M1 & 12.84 & 10.72 & 21.25 \\
M2 & 14.70 & 11.45 & 33.49 \\
M2m & 12.28 & 10.22 & 17.44 \\
M3 & 14.46 & 11.12 & 33.46 \\
M4 & 14.30 & 11.17 & 33.39 \\
M2\_oracle & 7.55 & 5.79 & 10.86 \\
\end{longtable}

\textbf{Table 5.} Retention on rolling h = 1 RMSE (point rule; margins
vs persist and vs M1 both shown).

\begin{longtable}[]{@{}
  >{\raggedright\arraybackslash}p{(\columnwidth - 8\tabcolsep) * \real{0.1765}}
  >{\raggedleft\arraybackslash}p{(\columnwidth - 8\tabcolsep) * \real{0.2353}}
  >{\raggedleft\arraybackslash}p{(\columnwidth - 8\tabcolsep) * \real{0.2353}}
  >{\raggedright\arraybackslash}p{(\columnwidth - 8\tabcolsep) * \real{0.1765}}
  >{\raggedright\arraybackslash}p{(\columnwidth - 8\tabcolsep) * \real{0.1765}}@{}}
\toprule\noalign{}
\begin{minipage}[b]{\linewidth}\raggedright
Model
\end{minipage} & \begin{minipage}[b]{\linewidth}\raggedleft
vs persist
\end{minipage} & \begin{minipage}[b]{\linewidth}\raggedleft
vs M1
\end{minipage} & \begin{minipage}[b]{\linewidth}\raggedright
Distinct structure
\end{minipage} & \begin{minipage}[b]{\linewidth}\raggedright
Decision
\end{minipage} \\
\midrule\noalign{}
\endhead
\bottomrule\noalign{}
\endlastfoot
M1 & 12.84 \textless{} 13.23 (\ensuremath{-}0.39) & --- & output only & retained
(point rule; margin within noise) \\
M2 & 14.70 \textgreater{} 13.23 (+1.47) & +1.86 & causal fluxes &
reject \\
M2m & 12.28 \textless{} 13.23 (\ensuremath{-}0.95) & \ensuremath{-}0.56 & no (affine AR(1)
function class) & listed; declined by protocol class clause \\
M3, M4 & worse than persist & worse & yes & reject \\
M2\_oracle & 7.55 & --- & uses future R, P & excluded \\
\end{longtable}

\textbf{Retention verdict (rolling h = 1).} Applying Definition 4.2 to
the h = 1 column of Table 4: M1 satisfies (H1) (12.84 \textless{}
13.23); as the simplest autonomous model it has no next-simpler causal
comparator, so (H2) is vacuous and M1 is retained by the point rule. M2
(14.70), M3 (14.46), and M4 (14.30) each fail (H1). M2m satisfies (H1)
(12.28 \textless{} 13.23) and, against its class-relative comparator M1,
also satisfies the (H2) comparison (12.28 \textless{} 12.84); it is
listed and then declined by the class clause of Section 4.1 --- a
protocol choice, not an outcome of the retention rule. M2\_oracle is
excluded as a diagnostic oracle. No stock-flow, residual, or delay
module is retained; on the rule as written, the verdict ``no stock-flow
module is retained'' is an artefact of the M2m class clause, since the
water-balance-identified affine map with climatological fluxes is the
best one-step forecaster tested. What the ladder actually shows is that
persisting last year's recharge is the failure (Section 5.4: persisted-R
RMSE 702 versus climatological 556 \ensuremath{\times} 10\textsuperscript{3} acre-ft on the recharge
target).

M1 is retained by the point-RMSE rule at the one-year horizon. The
margin is 0.39 ft on n = 75 and is not a significance claim; its
bootstrap interval covers zero (Section 5.3.1), MAE is a tie (10.72
versus 10.73 ft), and at h = 5 M1 (21.25 ft) does not beat persistence
(21.11 ft) while the training mean (16.80 ft) beats both. The retention
is therefore explicitly a one-year, RMSE-level statement ---
provisional, a coin-flip recorded by a point rule --- and it records a
slightly mean-reverting head series, not a confirmation of stock-flow
structure. At the decision scale of annual drought-stage declarations
the difference is operationally nil. The M2m decline and its 0.56-ft
advantage over M1 are likewise within noise (Section 5.3.1): the
estimator-level advantage of water-balance identification is visible in
the point scores but not separated from zero at this sample size.

\textbf{The post-2007 secondary-score record.} Definition 4.1's second
secondary score is scoped to the post-2007 origins, where the 660-ft
Stage I rule is in force. On those origins (\(n = 16\) at \(h = 1\)) the
committed records give 660-ft Brier scores of \(0.31\) (persistence),
\(0.25\) (M1), and \(0.19\) (the oracle), with the one-year RMSE ranking
unchanged (persistence \(13.09\), M1 \(12.16\), M2 \(13.31\), and the
oracle \(8.03\) ft); at \(h = 5\) (\(n = 12\)) the M1--persistence
ordering reverses on the subsample (\(17.16\) versus \(25.10\) ft),
reported without changing the one-year retention statement. The Brier
values are misclassification rates of a coarse annual-mean proxy for the
Authority's 10-day rule on a small subsample, reported from the archived
records and not over-read; the full post-2007 sub-sample record
(\(h = 1\) and \(h = 5\)) is archived with the analysis code
(\texttt{wave\_e\_edwards/results/rolling\_modern\_2007.csv}).

\subsubsection{5.3.1 Uncertainty on the retention margins (post-freeze
layer)}\label{uncertainty-on-the-retention-margins-post-freeze-layer}

A post-freeze uncertainty layer attaches Diebold--Mariano tests (Diebold
and Mariano 1995; Newey--West HAC, lag h \ensuremath{-} 1) and moving-block bootstrap
intervals (K\"unsch 1989; block length 8, 10,000 replications, seeded) to
every load-bearing margin, computed from the archived per-origin
forecast files. It is labelled post-freeze and changes no frozen
verdict.

\textbf{Table 6.} Uncertainty on the load-bearing margins (post-freeze
layer; Diebold--Mariano with Newey--West HAC, lag h \ensuremath{-} 1; moving-block
bootstrap, block length 8, 10,000 replications, seeded).

\begin{longtable}[]{@{}
  >{\raggedright\arraybackslash}p{(\columnwidth - 10\tabcolsep) * \real{0.1429}}
  >{\raggedleft\arraybackslash}p{(\columnwidth - 10\tabcolsep) * \real{0.1905}}
  >{\raggedleft\arraybackslash}p{(\columnwidth - 10\tabcolsep) * \real{0.1905}}
  >{\raggedleft\arraybackslash}p{(\columnwidth - 10\tabcolsep) * \real{0.1905}}
  >{\raggedright\arraybackslash}p{(\columnwidth - 10\tabcolsep) * \real{0.1429}}
  >{\raggedright\arraybackslash}p{(\columnwidth - 10\tabcolsep) * \real{0.1429}}@{}}
\toprule\noalign{}
\begin{minipage}[b]{\linewidth}\raggedright
Comparison (rolling, h = 1 unless noted)
\end{minipage} & \begin{minipage}[b]{\linewidth}\raggedleft
RMSE margin (ft)
\end{minipage} & \begin{minipage}[b]{\linewidth}\raggedleft
DM statistic
\end{minipage} & \begin{minipage}[b]{\linewidth}\raggedleft
DM p
\end{minipage} & \begin{minipage}[b]{\linewidth}\raggedright
95\% block-bootstrap CI (ft)
\end{minipage} & \begin{minipage}[b]{\linewidth}\raggedright
Separates from zero
\end{minipage} \\
\midrule\noalign{}
\endhead
\bottomrule\noalign{}
\endlastfoot
M1 \ensuremath{-} persist & \ensuremath{-}0.39 & \ensuremath{-}0.85 & 0.40 & {[}\ensuremath{-}1.51, +0.71{]} & no \\
M2m \ensuremath{-} persist & \ensuremath{-}0.95 & \ensuremath{-}3.07 & 0.003 & {[}\ensuremath{-}1.45, \ensuremath{-}0.68{]} &
\textbf{yes} \\
M2m \ensuremath{-} M1 & \ensuremath{-}0.56 & \ensuremath{-}1.61 & 0.11 & {[}\ensuremath{-}1.53, +0.13{]} & no \\
M2 \ensuremath{-} persist & +1.47 & +1.27 & 0.21 & {[}\ensuremath{-}0.02, +3.40{]} & no
(borderline) \\
mean \ensuremath{-} persist (h = 5) & \ensuremath{-}4.30 & \ensuremath{-}1.65 & 0.10 & {[}\ensuremath{-}9.72, \ensuremath{-}2.09{]} &
\textbf{yes} \\
M1 \ensuremath{-} persist (h = 5) & +0.15 & +0.06 & 0.95 & {[}\ensuremath{-}5.57, +4.94{]} & no \\
M2\_combo \ensuremath{-} M1 (climate) & \ensuremath{-}0.13 & \ensuremath{-}0.21 & 0.83 & {[}\ensuremath{-}1.37, +1.31{]} &
no \\
M2\_combo \ensuremath{-} M2m (nested baseline) & +0.43 & +0.74 & 0.46 & {[}\ensuremath{-}0.56,
+1.57{]} & no \\
\end{longtable}

Three readings. First, the retention margin itself (M1 \ensuremath{-} persist, 0.39
ft) is within noise, as is the M2m-over-M1 estimator margin and the M2
causal loss (the last only borderline, its interval touching zero).
Under the frozen point-RMSE rule the verdicts stand as frozen; under
uncertainty, the only one-year margin that separates from noise is M2m
over persistence (\ensuremath{-}0.95 ft, p = 0.003). Second, the five-year
climatology win (\ensuremath{-}4.30 ft) is robust: its bootstrap interval excludes
zero. The h = 5 DM statistic (p = 0.10 with HAC lag 4) is known to be
undersized under overlapping origins, and the bootstrap interval is the
primary evidence. Third, the nested-baseline margin (combo \ensuremath{-} M2m, +0.43
ft) is itself within noise: the climate rejection is a point-RMSE rule
outcome, not a significance finding.

An independent replication of this layer is registered with the
repository: the archived script
\texttt{campaign\_e3\_dm\_uncertainty.py} (in the batch-7 audit
directory), a second Diebold--Mariano + moving-block-bootstrap
implementation --- HAC lag h \ensuremath{-} 1 with unweighted truncation and
population-variance scaling, block length h, 20,000 replications, seed
0, deterministic (byte-identical on re-execution) --- reading the same
registered per-origin forecast file. Its verified output (archived
alongside it) reproduces every load-bearing reading: the M1 retention
margin DM z = \ensuremath{-}0.86 with block CI {[}\ensuremath{-}1.22, +0.56{]} ft (bootstrap p =
0.38); the M2m edge over persistence (p = 0.001); the h = 5 climatology
win (p = 0.035). It also registers two rows this table does not carry
--- the h = 1 climatology loss (+2.94 ft against persistence, p = 0.046)
and the M2 h = 5 loss (+12.4 ft, p \textless{} 0.001) --- both
consistent with the readings above. The two implementations use
different seeds, block lengths, HAC weighting, and variance conventions,
and agree on every conclusion: the 0.39-ft AR(1) retention margin is
statistically indistinguishable from zero, the M2m edge is the only
one-year margin that separates from noise, and the five-year climatology
win survives.

\subsubsection{5.4 Climate-informed
recharge}\label{climate-informed-recharge}

Information known at 31 December of year t comprises R\_t, the
September--November Ni\~no 3.4 anomaly of year t, and calendar-year
precipitation in Texas climate divisions 06 and 07. December--February
precipitation that includes January of t+1 is not used.

\begin{figure}[htbp]
\centering
\includegraphics[width=\linewidth]{figs_e3/fig4_pass2.png}
\caption{Rolling RMSE on J-17 for climate-informed recharge. The precipitation oracle uses year t+h precipitation and cannot be retained.}
\end{figure}

\textbf{Table 7.} Rolling RMSE, climate-informed recharge (same origin
sets as Table 4; margins vs M1 --- the frozen Pass-2 gate --- and vs the
nested M2m baseline both shown for the h = 1 column).

\begin{longtable}[]{@{}
  >{\raggedright\arraybackslash}p{(\columnwidth - 8\tabcolsep) * \real{0.1579}}
  >{\raggedleft\arraybackslash}p{(\columnwidth - 8\tabcolsep) * \real{0.2105}}
  >{\raggedleft\arraybackslash}p{(\columnwidth - 8\tabcolsep) * \real{0.2105}}
  >{\raggedleft\arraybackslash}p{(\columnwidth - 8\tabcolsep) * \real{0.2105}}
  >{\raggedleft\arraybackslash}p{(\columnwidth - 8\tabcolsep) * \real{0.2105}}@{}}
\toprule\noalign{}
\begin{minipage}[b]{\linewidth}\raggedright
Model
\end{minipage} & \begin{minipage}[b]{\linewidth}\raggedleft
H, h=1 (ft)
\end{minipage} & \begin{minipage}[b]{\linewidth}\raggedleft
margin vs M1
\end{minipage} & \begin{minipage}[b]{\linewidth}\raggedleft
H, h=5 (ft)
\end{minipage} & \begin{minipage}[b]{\linewidth}\raggedleft
R, h=1 (10\textsuperscript{3} acre-ft)
\end{minipage} \\
\midrule\noalign{}
\endhead
\bottomrule\noalign{}
\endlastfoot
persist H / persist R & 13.23 & +0.39 & \textbf{21.11} & 702 \\
M1 & 12.84 & --- & 21.25 & --- \\
M2\_Rar & 13.25 & +0.41 & 25.38 & 561 \\
M2\_Renso & 12.82 & \ensuremath{-}0.02 & 24.42 & 528 \\
M2\_Rprecip & 12.80 & \ensuremath{-}0.04 & 25.38 & 545 \\
M2\_combo & 12.71 & \ensuremath{-}0.13 & 26.88 & 538 \\
M2m (the declined nested baseline) & 12.28 & \ensuremath{-}0.56 & 17.44 & --- \\
rain climatology & --- & --- & --- & 556 \\
rain oracle & 10.56 & --- & 16.91 & \textbf{354} \\
\end{longtable}

Same-year corr(R, precipitation) = 0.78, which is why the precipitation
oracle reduces head RMSE to 10.56 ft; it remains worse than the full (R,
P) oracle (7.55 ft) because the linear rain map misses 1957-scale
extremes. That last clause is a result: the oracle gap is not closed by
knowing precipitation alone.

Lagged precipitation and September--November Ni\~no 3.4 have modest skill
on R relative to climatology (528--545 versus 556 \ensuremath{\times} 10\textsuperscript{3} acre-ft), and
they do not constitute forecast structure on head. The frozen Pass-2
gate compares the climate modules against persistence and the retained
M1 (Section 4.1): on that gate the three modules are listed by the point
rule --- each beats persist and M1 by margins of 0.02, 0.04, and 0.13
ft, all within noise (Section 5.3.1, Table 6) --- so the listing is not
a skill claim. None is retained, on three stated grounds: the margins
against the gate are within noise; at h = 5 all three have RMSE 3--6 ft
higher than persistence (the h \textgreater{} 1 climate scores reuse the
one-step recharge forecast, held constant over the horizon, so this is a
design consequence, not a multi-year climate test); and structurally
they are the declined M2m with a weakly adjusted intercept --- no
additional forecast structure over the retained AR(1) class --- so
against that nested climatological baseline they lose outright
(12.71--12.82 versus 12.28 ft; the combination's +0.43-ft deficit,
Section 5.3.1, also within noise). The rejection is therefore honest
with both readings visible: the climate modules beat persistence and the
AR(1) by at most 0.13 ft and lose to climatological fluxes; none is
retained, and nothing in that verdict is a significance finding. M2\_Rar
loses at h = 1 (13.25 ft; 0.41 ft worse than M1) --- autoregression on
nearly white recharge is not a recharge forecast.

On the 1957--61 recovery, no causal climate-informed module has lower
RMSE than persistence (persist 43.6 ft; best causal about 48.8 ft;
precipitation oracle 33.7 ft). September--November 1956 is La Ni\~na
(\ensuremath{-}0.92) and does not announce R\textsubscript{1957} = 1143.

The remaining fixed windows complete the same record: drawdown 1951--56,
M2\_Renso 24.30, M2\_Rprecip 28.67, M2\_combo 24.70, precipitation
oracle 18.16 ft; pre-permit wet 1991--95, M2\_Rprecip 22.03, M2\_Renso
23.03, M2\_Rar 23.94, M2\_combo 25.71, oracle 10.98 ft; critical-period
era 2015--23, M2\_Rprecip 14.52, M2\_Rar 14.67, M2\_Renso 16.01,
M2\_combo 16.57, oracle 9.75 ft --- on this one window the two
climate-informed modules M2\_Rprecip and M2\_Rar edge past M1 by about
one foot (14.52 and 14.67 versus 15.62 ft). On the recharge target
itself the fixed-window scores are an order of magnitude coarser on
every window (climate modules 199--937; precipitation oracle 80.7--487,
against the 556 \ensuremath{\times} 10\textsuperscript{3} acre-ft climatology scale of the rolling record),
so the marginal head advantage is a window-specific result, not a
recharge forecast, and it does not enter the abstract.

Closing the gap between persistence and the oracle would require next
year's recharge, which is not available at the annual forecast origin.

\subsubsection{5.5 The service series after the retention
freeze}\label{the-service-series-after-the-retention-freeze}

\begin{figure}[htbp]
\centering
\includegraphics[width=\linewidth]{figs_e3/fig5_fibre.png}
\caption{Comal annual mean versus J-17. Contemporaneous r = 0.986. The service series is a measured rating curve of the same state --- a measured service, not an independent information source.}
\end{figure}

\textbf{The Comal channel (rating-curve reading).} The map Q = c\textsubscript{0} + c\textsubscript{1}H
fitted on 1934--1950 (c\textsubscript{0} = \ensuremath{-}2876, c\textsubscript{1} = 4.77) is a linear channel of the
same state indexed by J-17. The full-sample contemporaneous correlation
is r = 0.986; the 1934--1950 fit is quoted because the channel is used
on that train, and the full-sample correlation is the redundancy
statement.

The channel fails in the drought tail, and the direction matters: at the
1956 annual mean (623.15 ft) it reads \ensuremath{\approx}97 cfs against an observed \ensuremath{\approx}32
cfs annual mean --- it predicts non-cessation --- and its zero-discharge
level (2876/4.77 \ensuremath{\approx} 602.9 ft) lies below every observed head (daily
minimum 612.5 ft), so the map never predicts cessation where cessation
occurred. Near-linear at ordinary heads, the channel fails at cessation
exactly where management cares; that tail failure is the one place Comal
is not redundant with J-17. One-year Comal RMSE (cfs), constructed by
scoring the same ladder directly on the Comal series over the same
rolling origins (n = 75): persist 71.9, M1 69.0, M2m 68.7, M2 74.8, M3
73.8, M4 73.4, train-mean 89.7, oracle 45.3 --- the ranking mirrors the
head ranking, as it must for a linear channel.

The service series does not change retention, and it cannot: it is
nearly a linear transform of head. Comal is an independent measurement
(a USGS spring gauge, not used to construct J-17), but it is not an
independent information source. Gravimetric storage or the J-27 Uvalde
index would be different objects, not a second observation channel of
this specification. Eastern-basin recharge was stored and not used.

\subsubsection{5.6 Pumpage
counterfactuals}\label{pumpage-counterfactuals}

A declared scenario layer passes counterfactual pumpage paths through
the pre-permit affine map (train 1980--1990, the ladder's window 3; the
short-window identification caveat applies to the map itself), from the
observed 1990 head (645.8 ft) over 1991--2023 with the actual recharge
sequence. The scenarios are counterfactual simulations of the fitted map
--- declared scenarios, not forecasts, and no retention implication.

\textbf{Table 8.} Pumpage counterfactuals, 1991--2023 (fitted pre-permit
map; actual recharge).

\begin{longtable}[]{@{}
  >{\raggedright\arraybackslash}p{(\columnwidth - 8\tabcolsep) * \real{0.1579}}
  >{\raggedleft\arraybackslash}p{(\columnwidth - 8\tabcolsep) * \real{0.2105}}
  >{\raggedleft\arraybackslash}p{(\columnwidth - 8\tabcolsep) * \real{0.2105}}
  >{\raggedleft\arraybackslash}p{(\columnwidth - 8\tabcolsep) * \real{0.2105}}
  >{\raggedleft\arraybackslash}p{(\columnwidth - 8\tabcolsep) * \real{0.2105}}@{}}
\toprule\noalign{}
\begin{minipage}[b]{\linewidth}\raggedright
Scenario
\end{minipage} & \begin{minipage}[b]{\linewidth}\raggedleft
Mean pumpage (10\textsuperscript{3} acre-ft yr\textsuperscript{-1})
\end{minipage} & \begin{minipage}[b]{\linewidth}\raggedleft
End head 2023 (ft)
\end{minipage} & \begin{minipage}[b]{\linewidth}\raggedleft
Minimum head (ft)
\end{minipage} & \begin{minipage}[b]{\linewidth}\raggedleft
RMSE vs observed heads (ft)
\end{minipage} \\
\midrule\noalign{}
\endhead
\bottomrule\noalign{}
\endlastfoot
actual pumpage & 382.1 & 641.6 & 640.2 & 8.56 \\
pumpage frozen at 1990 & 489.4 & 630.9 & 630.5 & 14.22 \\
pre-permit mean pumpage & 469.8 & 632.4 & 632.0 & 13.02 \\
20\% pumpage cut & 305.7 & 646.8 & 645.5 & \textbf{7.19} \\
\end{longtable}

The observed 2023 head is 635.7 ft. Four readings follow. First, pumpage
is a secondary lever in this map: the full spread of counterfactual
policies --- from a 20\% cut below actual to freezing at the 1990 peak
--- spans 630.9--646.8 ft at 2023, and the map's own RMSE against the
observed record (7.2--14.2 ft) is of the same order as the policy spread
itself; the 5-ft effects below are within model error and are not policy
estimates. Second, the 20\%-cut path is the closest to the observed
record of the four scenarios (7.19 versus 8.56 ft) --- an artefact of
the map running high on this window (the actual-pumpage path ends 5.9 ft
above the observed 2023 head), not evidence that the Authority's pumpage
was near-optimal; the counterfactuals do not repair the ladder's primary
failure, which Section 5.3 traces to the recharge series' persistence
failure. Third, the direction is right and the magnitude modest: the
20\% cut raises the simulated 2023 head by 5.2 ft (646.8 \ensuremath{-} 641.6)
relative to the actual-pumpage path. Fourth, the recharge coefficient
dominates the map through its range, not its magnitude: a 1000 \ensuremath{\times} 10\textsuperscript{3}
acre-ft recharge difference moves the head by roughly 17.5 ft at the
pre-permit coefficient (\(\hat{\beta}\) \ensuremath{\approx} 0.0175 on the 1980--1990 train; the pumpage
coefficient on that train is \ensuremath{-}0.031 ft per 10\textsuperscript{3} acre-ft, larger per unit
but acting on a range of \textasciitilde300--540 against recharge's
44--2,486), and \ensuremath{\gamma} is an ordinary-least-squares association the paper
itself diagnoses as simultaneity-contaminated (Section 5.2), not a
policy lever. The full-sample coefficients (\(\hat{\beta}\), \(\hat{\gamma}\)) = (0.017, \ensuremath{-}0.026)
differ from the pre-permit window's because the windows differ.

\subsection{6. Discussion}\label{discussion}

The one-pool increment \ensuremath{\Delta}H tracks recharge (r = 0.74), and the oracle's
RMSE is 43\% below persistence at h = 1 and 49\% below at h = 5: the map
accounts for contemporaneous increments when the year's water is known.
As a one-year forecast, causal stock-flow fails because the dominant
increment is not persistent; residual and delayed variants inherit that
timing error. Recoveries in 1957 and 1992 are missed by every causal
model and captured by the oracle. The identified driver has the right
timing and is unforecastable one year ahead --- the same pattern a
companion evaluation finds on Northern cod, where a more accurate catch
series does not rescue constant-productivity surplus production (Author
et al., in review). The two score tables are not pooled.

A second question is why the AR(1)'s edge over persistence is as small
as it is, and here the fitted head autocorrelation is the operative
quantity. An estimated AR(1) with coefficient \ensuremath{\rho} nests persistence at
\ensuremath{\rho} = 1; as the head series approaches that limit the two forecasts
converge, and the AR(1)'s advantage shrinks toward the noise in
estimating \ensuremath{\rho} itself. On this record the head is persistent but not
nearly a unit root (annual AC(1) = 0.64), which leaves the AR(1) a real
but slim 0.39-ft margin --- consistent with the coin-flip diagnostics
reported above. The prediction implied is directional and cheap to
state: on a more strongly autocorrelated head record the AR(1) should
beat persistence by \emph{less}, not more, and at sufficiently high
AC(1) it should lose. An exploratory port of this ladder to a second,
separately regulated pool of the same aquifer (annual head AC(1) = 0.84)
is consistent with that direction --- there the AR(1) loses to
persistence by 0.13 ft --- but a single additional pool is a direction,
not a test, and no claim is made here beyond the mechanism. The point
for the present result is that the AR(1)'s retention is a
weak-autocorrelation phenomenon and should not be read as evidence that
an output-only model is generally preferable.

Two features distinguish this design from the groundwater benchmark
literature it engages. First, the retention gate: GEMS-GER ships three
benchmark models and reports the fraction of wells for which the best
one reaches NSE \textgreater{} 0.5 (Ohmer et al.~2026), and the karst
benchmark of Zhu et al.~(2026) ranks nine architectures by RMSE and R\textsuperscript{2}
--- both compare model families against one another, while here every
module must beat persistence and the next-simpler causal model under a
rule frozen before scoring, and on that rule no causal module is
retained at the one-year horizon (the climate rungs pass the frozen
Pass-2 gate by within-noise margins and are declined on class grounds;
Section 5.4). Second, the horizon contrast: no benchmark study of this
basin reports a training mean beating persistence at a longer horizon.
The five-year climatology result (16.80 versus 21.11 ft) is the
specification's most directly transferable candidate for management
planning that keys on multi-year outlooks, with the scope stated: the
San Antonio Pool is a rapidly recharged, institutionally bounded system,
so the training mean is informative at long horizons here, and a
non-stationary fossil aquifer under sustained depletion would not
sustain a mean-reverting baseline.

Two specification points carry over to the interpretation. First,
next-year pumpage is partially an institutional scenario --- permits and
critical-period rules respond to head --- rather than a purely exogenous
flux like rainfall; persisting last year's pumpage is a declared
simplification, and the oracle pumpage is in no information set. Second,
the karst setting matters. The annual affine one-pool specification is
not a karst model: conduits, the Uvalde--San Antonio divide, unconfined
recharge-zone storage, and the confined-zone pressure response remain in
the residual, and the long-standing question whether lumped or
equivalent-porous-media representations can carry regional flow (Scanlon
et al.~2003) is inherited, not resolved, by this design. What the design
shows is that at the annual origin, the information carried by such a
map is timing-bound: with realized fluxes the same map is a nowcast of
the year's head (7.55 ft RMSE), and the contemporaneous increment \ensuremath{\Delta}H\_t
against R\_t (r = 0.74) is the closure statement; the word
``certificate'' is retired for both. The retained AR(1) admits a
complementary reading: to first order, spring discharge obeys Darcy
proportionality to head above the spring level, and the autonomous
solution of that drainage law is an affine autoregression
\(H_{t+1} \approx (1-k)H_t + k H_s\); the fitted \(\hat\varphi = 0.66\)
(the M1 estimate; the full-sample Pearson correlation corr(H\_t,
H\_\{t\ensuremath{-}1\}) = 0.64 is a different quantity) is consistent with a
drainage-decay coefficient \(\hat k \approx 0.34\) yr\textsuperscript{-1}. The module is
output-only in the protocol's sense --- no flux data enter the forecast
--- but the shape it fits carries the aquifer's own free-drainage
momentum; this is an interpretation in discussion, not a retention
reason.

The wrong leading indicator is worth recording. The 1950s decline is not
a Stage I event: the 660-ft line is a 2007 rule, and slack to 660 ft is
the wrong leading indicator for 1951--56, just as the 2016 Northern cod
limit reference point was the wrong leading indicator for 1983--90 in
the companion study. 1956 is a physical near-cessation at Comal (annual
mean springflow 32 cfs; daily J-17 minimum 612.51 ft).

A two-pool exchange module was specified and not fitted; no barrier or
exchange term was fitted. Post-1997 pumpage no longer spikes with
drought as in 1956 (\(321 \times 10^3\) acre-ft from wells); that change
is already in P\_t, and a separate critical-period switch adds no degree
of freedom beyond the pumpage series.

Limitations follow the data and the map. Total-area R and P mix the San
Antonio and Uvalde pools; recharge is a Puente estimate; pumpage
includes unreported domestic, livestock, and federal use. Neither series
is head: R and P are constructed fluxes, not observations of J-17. M4 is
a one-year information delay and a symmetry control with the companion
fisheries evaluation --- J-17 is a telemetered gauge whose head is
recorded daily, so the module prices a theoretical information lag
rather than a constraint of this system's monitoring. The sample is 90
years with four short test windows, and the 0.39-ft AR(1) margin is not
a significance claim (Section 5.3.1: its interval covers zero).
Additional climate indices (PDO, AMO) were not run; they would enter
through the same one-step recharge forecast and cannot alter the timing
result. A mid-year nowcast would require a new evaluation protocol.

\subsection{7. Conclusions}\label{conclusions}

On locked J-17 annual-mean head, last-value persistence is more accurate
than a causal one-pool balance that persists last year's recharge,
because annual recharge is near-white (r = 0.17) while the
contemporaneous increment tracks it (r = 0.74). The univariate AR(1)
improves one-year RMSE by 0.39 ft and is retained as an output-only
model by the frozen point rule; the margin is a coin-flip (MAE tie,
five-year loss, bootstrap interval covering zero), and the same affine
class fitted through the water balance with climatological fluxes does
better still (12.28 ft, the only one-year margin separated from noise)
but is declined by a protocol class clause outside the retention rule.
The same balance, given realized recharge and pumpage, nowcasts the
current year (7.55 ft, a 43\% reduction); climate variables known at the
annual origin do not recover that gap; and at five years, climatology
wins, robustly (16.80 versus 21.11 ft, interval excluding zero).

\emph{Transferable lesson.} The result is not specific to J-17. The
benchmark-discipline claim --- that a deliberately simple process-based
model must be justified against naive persistence and the training mean,
and that it may fail to earn that justification when the driver it
persists is near-white --- is a statement about the evaluation design, and
it replicates in the companion Northern cod forecast study (Abaee 2026),
where scored surplus-production structure likewise failed to beat
last-value persistence at the annual origin. The underlying mechanism
differs by system (a near-white annual recharge here; a weakly
autocorrelated surplus there), but the test does not: each system is
held to the standard that added structure clears a naive benchmark. That
is the general content of the paper; the J-17 basin is its
demonstration.

\emph{Practical reading.} Two choices follow for a manager. At the
annual origin, do not prefer a one-pool water-balance forecast over
persistence on timing grounds --- the map is only as good as the year's
realized fluxes, which are not known at the forecast origin. For
multi-year planning, prefer climatological baselines to persisted
recharge: the five-year training-mean result (16.80 versus 21.11 ft)
beat persistence with an interval excluding zero, and it is the most
directly transferable candidate for a district that keys on multi-year
outlooks. The pumpage counterfactuals (Section 5.6) add a third,
caveated reading: pumpage is a secondary lever on this map (the 20\% cut
raises the simulated 2023 head by 5.2 ft, but the full policy spread is
within the map's own error band), so pumpage regulation is not a
substitute for forecasting skill. A fourth reading concerns the AR(1)
itself: its retention here is a weak-autocorrelation phenomenon (Section
6), so a district whose head record is more strongly autocorrelated than
J-17's should expect the AR(1)'s edge over persistence to be smaller,
not larger, and should re-score rather than inherit this verdict.

\emph{Generalizability boundary.} The class of systems for which the
finding is expected is stated explicitly: a rapidly recharged,
institutionally bounded aquifer, not a non-stationary fossil aquifer
under sustained depletion. Conducts, the Uvalde--San Antonio divide,
unconfined recharge-zone storage, and the confined-zone pressure
response remain in the residual for this one-pool map (Section 6), and
the single-well test bed is the San Antonio Pool's well J-17 only; no
transfer to the Uvalde Pool (J-27) is claimed, consistent with the
programme's no-pooling discipline. The Uvalde figure quoted in Section 6
is used only to state the direction of a mechanism, not to score a
model: no Uvalde result enters the retention verdicts, the two records
are not pooled, and that exploratory comparison is not offered as a test
of anything claimed here.

The paper reports a forecast comparison. It does not conclude that the
aquifer is unsustainable, and it does not conclude that the evaluation
framework is empirically confirmed on this basin.

\subsection{References}\label{references}

Adamowski, J., and Chan, H.F. 2011. A wavelet neural network conjunction
model for groundwater level forecasting. \emph{Journal of Hydrology}
407: 28--40. https://doi.org/10.1016/j.jhydrol.2011.06.013

Abaee, A. 2026. Does a surplus-production ladder improve forecasts of
Northern cod? A scored test on NAFO 2J3KL. Zenodo.
https://doi.org/10.5281/zenodo.22553609. Companion forecast-evaluation
study (Northern cod, NAFO 2J3KL).

Daliakopoulos, I.N., Coulibaly, P., and Tsanis, I.K. 2005. Groundwater
level forecasting using artificial neural networks. \emph{Journal of
Hydrology} 309: 229--240. https://doi.org/10.1016/j.jhydrol.2004.12.001

Diebold, F.X., and Mariano, R.S. 1995. Comparing predictive accuracy.
\emph{Journal of Business \& Economic Statistics} 13: 253--263.
https://doi.org/10.1080/07350015.1995.10524599

Edwards Aquifer Authority. 2024/25. \emph{2023 Groundwater Discharge and
Usage}, Table 1 (after USGS letter report, 5 April 2024).
https://www.edwardsaquifer.org/wp-content/uploads/2025/05/2023-Groundwater-Discharge-and-Usage.pdf

Edwards Aquifer Authority. Critical Period / Drought Management.
https://www.edwardsaquifer.org/groundwater-users/critical-period-drought-management/

K\"unsch, H.R. 1989. The jackknife and the bootstrap for general
stationary observations. \emph{Annals of Statistics} 17: 1217--1241.
https://doi.org/10.1214/aos/1176347265

Makridakis, S., Spiliotis, E., and Assimakopoulos, V. 2020. The M4
Competition: 100,000 time series and 61 forecasting methods.
\emph{International Journal of Forecasting} 36: 54--74.

NOAA National Centers for Environmental Information. nClimDiv
precipitation (climdiv-pcpndv-v1.0.0).
https://www.ncei.noaa.gov/pub/data/cirs/climdiv/

NOAA Physical Sciences Laboratory. Ni\~no 3.4 monthly SST (HadISST).
https://psl.noaa.gov/data/timeseries/month/data/nino34.long.data

Ohmer, M., Liesch, T., Habbel, B., Heudorfer, B., Gomez, M., Clos, P.,
N\"olscher, M., and Broda, S. 2026. GEMS-GER: A machine learning benchmark
dataset of long-term groundwater levels in Germany with meteorological
forcings and site-specific environmental features. \emph{Earth System
Science Data} 18, 77.

Puente, C. 1978. \emph{Method of estimating natural recharge to the
Edwards Aquifer in the San Antonio area, Texas.} U.S. Geological Survey,
Austin, Texas.

Ropelewski, C.F., and Halpert, M.S. 1986. North American precipitation
and temperature patterns associated with the El Ni\~no/Southern
Oscillation (ENSO). \emph{Monthly Weather Review} 114: 2352--2362.

Scanlon, B.R., Mace, R.E., Barrett, M.E., and Smith, B. 2003. Can we
simulate regional groundwater flow in a karst system using equivalent
porous media models? Case study, Barton Springs Edwards aquifer, USA.
\emph{Journal of Hydrology} 276: 137--158.

Texas Water Development Board. Water Data for Texas, well 6837203
(J-17). https://waterdatafortexas.org/groundwater/well/6837203

Umphres, G.D., and Choi, N.J. 2025. Estimated Annual Recharge to the
Edwards Aquifer in the San Antonio Area, by Stream Basin or Ungaged
Area, 1934--2024. U.S. Geological Survey data release.
https://doi.org/10.5066/P1BI62NY

U.S. Geological Survey. National Water Information System, site
08168710, Comal Springs at New Braunfels, Texas.

Zhu, Q., Zhu, Y., Niu, J., Huang, J., Huang, F., Zhou, X., Liu, D., and
Hu, B.X. 2026. Benchmarking machine learning and deep learning models
for groundwater level prediction in karst aquifers: The dominant role of
hydrogeological complexity. \emph{Water} 18, no. 8: 939.

\section*{Declarations}

\subsection*{Data Availability Statement}

All input data, analysis scripts, and result files are archived in the
public repository at
https://github.com/MIKEAA2020/general-sustainability, together with the
frozen scoring protocols (dated 2026-08-25, locked before any score was
generated). J-17 daily highs: Texas Water Development Board well
6837203. Recharge: Umphres and Choi (2025), USGS data release,
https://doi.org/10.5066/P1BI62NY. Pumpage: Edwards Aquifer Authority
Table 1. Comal Springs: USGS 08168710. Ni\~no 3.4: NOAA PSL HadISST (raw
file registered with the repository). Precipitation: NCEI nClimDiv ---
the raw file is not distributed with the repository (provenance URL
archived in the sources index), so the three precipitation columns of
the fixed panel are not reproducible from the registered code alone,
while the two Ni\~no columns rebuild from the registered file to machine
precision; scoring from the registered analysis panel does not require
the nClimDiv file. The registered twenty-column analysis panel is the
dataset of record for all scored analyses. All computations are
deterministic: re-executing the registered scripts in a fresh
environment regenerated every archived result file byte for byte, and
all scored rows recompute from the per-observation forecast files and
the registered series. The pumpage counterfactual layer (Section 5.6) is
produced by
\texttt{rerun\_campaigns/campaign\_e3\_pumpage\_scenarios.py}, archived
alongside its outputs, and regenerates them exactly. The post-freeze
uncertainty layer (Section 5.3.1) is produced by
\texttt{wave\_e\_edwards/src/e3\_audit\_uncertainty.py} (seeded,
deterministic; Diebold--Mariano with Newey--West HAC and moving-block
bootstrap on the archived per-origin forecast files), with its outputs
archived as
\texttt{wave\_e\_edwards/results/e3\_audit\_uncertainty.json}; the
clip-binding statement of Section 5.2 is computed there from the
registered panel and reproduces both fixed-window M2 RMSEs (18.11 and
55.32 ft) exactly. The independent replication of the post-freeze
uncertainty layer is registered as
\texttt{batch\ 7\ (audits\ of\ agent\ arena\ 1\ paper\ rewrites)/campaign\_e3\_dm\_uncertainty.py}
(archived), with its deterministically reproduced output archived
alongside it at
\texttt{batch\ 7\ (audits\ of\ agent\ arena\ 1\ paper\ rewrites)/results/e3\_dm\_uncertainty.csv};
it recomputes the Diebold--Mariano and moving-block-bootstrap table from
the same registered per-origin forecast file under different seeds,
block lengths, and HAC conventions, and every load-bearing conclusion of
Section 5.3.1 is unchanged under either implementation.

\subsection*{Funding}
None.

\subsection*{Competing interests}
None.

\subsection*{Code availability}
The scored ladder, the post-freeze uncertainty layer and the pumpage
counterfactuals are produced by the scripts named above, archived with
their outputs. The verification script
\texttt{paperE3\_edwards\_forecast\_ladder\_v17\_verification.py}
re-runs the ladder from the registered panel and recompares every cell
of Tables~3, 4 and 7 against the recomputed scores; it is archived
alongside this manuscript.

\subsection*{AI declaration}
GLM (Z.ai), Qwen (Alibaba Cloud) and DeepSeek AI assisted with drafting
and iterative review. The author reviewed and edited outputs and takes
responsibility for the final work.

\end{document}


% ===== PART: ws_v17  (source: ws_v17.tex) =====
\clearpage

\documentclass[10pt,twocolumn]{article}
\usepackage[a4paper,margin=15mm]{geometry}
\usepackage{amsmath,amssymb}
\usepackage{amsthm}
\usepackage{graphicx}
\theoremstyle{definition}
\newtheorem{proposition}{Proposition}
\newtheorem{remark}{Remark}
\usepackage{booktabs,array,calc}
\usepackage[font=small,labelfont=bf]{caption}
\usepackage[colorlinks=true]{hyperref}
\providecommand{\tightlist}{\setlength{\itemsep}{0pt}\setlength{\parskip}{0pt}}
\renewcommand{\thesubsection}{\arabic{subsection}}
\makeatletter
\renewcommand{\@seccntformat}[1]{\csname the#1\endcsname.\quad}
\makeatother
\emergencystretch=3em
\allowdisplaybreaks
\setlength{\tabcolsep}{4pt}

\begin{document}
\title{Exact Audits of Worked Systems for the\\
Obstruction Calculus:\\
Counts, Witnesses, and Monitoring Design Rules}
\author{Amin Abaee\\[0.35em]
{\small Independent Researcher}\\[0.55em]
{\small\href{https://orcid.org/0000-0002-0019-1842}{ORCID: 0000-0002-0019-1842}}\\[0.3em]
{\small\href{mailto:amin\_abaee@ut.ac.ir}{amin\_abaee@ut.ac.ir}}}
\date{September 25, 2026}
\maketitle

\begin{abstract}
Viability certificates are useful in applications only insofar as their behaviour on concrete systems is understood precisely. This paper develops one exact audit framework --- integer and rational arithmetic throughout, every
count regenerated by script --- and applies it to a family of worked
systems that instantiate the mechanisms of the obstruction calculus
: a shared finite two-coordinate system with observation
structures (aggregation, policy-class restrictions, decentralized
observation, regime uncertainty), a review-timing grid for a
hidden-regime resource model, and a continuous two-floor production
benchmark. The audits yield the following. The review-timing identity on
the 93-cell grid is non-strict: viability at the deadline holds if and
only if \(T_{\mathrm{obs}} \le \tau^{*}(z_0) = z_0 - 1\), boundary included; the strict form against the real exit bound fails exactly on the boundary cells (against the floored count \(\sigma^{*}\), on one cell per row from \(z_0 = 2\)). On the four-state
monitoring target, 7 of the 15 partitions are adequate with two incomparable coarsest adequate two-cell partitions (minimal in the refinement order --- hence no single coarsest adequate design exists), and the two maximal common-action sets overlap. Four policy classes have kernel
sizes 24, 26, 25, and 28 of 36 belief pairs --- the classes form a two-by-two product of aggregation level and protocol --- with incomparable middle
cells. Decentralized unanimity over a
\(256\)-law space (\(64\) behaviorally distinct) sustains 12 of 36
pairs; a whole-law authority restriction collapses the kernel to zero,
and either agency's low vote forced against the target kills it. Under
frozen regimes no pair is lost (the corner pins its state); mode-blind
operation loses exactly the pairs containing the corner state
\((1,1)\). The biased certainty-equivalence law has
drift at least \(21/100\) on \([1,2]\); the bias-free law has drift
identically zero. The benchmark's critical aggregate is \(Y^{*} = 27/5\) (feasible,
margin \(0\)); at the audited obstructed aggregate \(Y = 6\) a
rational dual certificate prices the shortfall at margin \(3/50\). Static duality
holds at value \(-1/10\) on the two-floor instance with an \(\ell_1\)-normalized
Farkas pair, and a Helly-tight triple exhibits the sparse \(m{+}1\)-constraint witness. The complete result tables are carried here: the master table
of all 36 belief pairs against the five kernels, the sixteen sustaining
law pairs of the decentralized protocol (an exact \(4 \times 4\)
product), the 93-cell timing grid, the fifteen-partition census with
each fibre's common safe-action set, and the benchmark and drift tables,
with a five-figure set regenerated from the same scripts. Design rules
follow: review deadlines are boundary-inclusive; merged fibres call for
splits, not threshold tuning; aggregated indices support certainly-safe
reporting only; biased observation maps call for structural correction.
\end{abstract}

\textbf{Keywords:} viability kernels; obstruction certificates; exact
arithmetic; monitoring design; \mbox{decentralized observation}; certainty
equivalence; max-plus algebra

\subsection{Introduction}\label{introduction}

The obstruction calculus states sound sufficient conditions for
nonviability --- obstruction certificates --- mechanism by mechanism:
common-action failure (Sections \ref{master}--\ref{adequacy}),
review-timing failure (Section \ref{timing}), finite-time exit under
regime alternation (Section \ref{regimes}), and failure induced by
restricting the policy class or by a biased observation map
(Sections \ref{classes}, \ref{drift}) (Abaee, 2026, An obstruction calculus). Each mechanism is
accompanied in application by quantitative questions. How many belief
states does an observation structure destroy, and which ones? At which
review deadlines does the timing certificate fire, and does it fire at
the boundary? What exactly does a vote protocol cost, and which law in
the protocol's space sustains a given target? This paper answers such
questions exactly for a family of worked systems, under one audit frame
and one notation. The sharp thesis, made explicit by the master
theorem of Section~\ref{mtheorem}: obstructions arise from five
mechanisms --- common-action failure, memory/register coupling,
finite-time boundary exit, decentralized information loss, and joint
convex infeasibility --- and each admits a certificate with an explicit
semantic scope, a monotonicity status, and a robustness margin.

Three properties define the frame. First, all computations are carried
out in integer and rational arithmetic: no solver, no tolerance, no
simulation. Second, every discrete census in the family is finite --- the continuous objects reduce to exact rational calculations on rationally specified domains --- and every count is regenerated by a deterministic standard-library script, so any two executions agree. Third, each audited quantity is stated together with a
witness or a certificate --- an action sequence, a partition, a law
pair, a dual measure --- so the counts carry their own evidence. This paper carries the audit's complete result tables and its figure set:
the per-pair membership table behind the kernel counts, the sustaining
law pairs behind the decentralized count, the full timing grid, the full
partition census, and the benchmark and drift tables, every entry of
which is regenerated from the same scripts that verify the text.

The systems are drawn from the applied model of Abaee (2026, Robust viability of the 2J3KL limit reference point), in which
a two-coordinate resource stock is managed by acting one of two
instruments under incomplete observation. The family comprises: the
shared audit system with its observation structures (Sections
\ref{master}--\ref{decentralized} and \ref{adequacy}--\ref{regimes});
the hidden-regime review grid (Section
\ref{timing}); a continuous two-floor production benchmark (Section
\ref{benchmark}); and a multiple-floor line extending the Farkas
certificate to jointly binding constraints (Section \ref{floors}). The
shared audit system is the running instance on which the paper's counts
are developed; each further member of the family is introduced against
it. The contributions are: the master monotonicity theorem --- one
audited-tuple object whose kernel comparisons unify the policy,
observation, and memory enlargements (Section \ref{mtheorem}); the
complete exact audit of the worked family --- the master table, the
decentralized count, the full timing grid, the partition census, and the
benchmark and drift tables, every entry regenerated by the verification
scripts (Sections \ref{master}--\ref{regimes}); and the design rules
the audits license --- boundary-inclusive review deadlines,
merged-fibre caution, the monitoring-adequacy criterion, and the
five-mechanism obstruction taxonomy with its certificate semantics
(Sections \ref{timing}--\ref{adequacy}).

\subsection{Systems, conventions, and methods}\label{frame}

\emph{The shared audit system.} States are pairs
\(x = (z_1, z_2) \in \{0,1,2,3\}^{2}\); the safe set is
\(\mathcal{V} = \{z_1 \ge 1,\ z_2 \ge 1\}\) with the two coordinates as
floors. Instruments are \(u \in \{0, 1, 2\}\): \(u = j \in \{1,2\}\)
acts coordinate \(j\), which regenerates one unit up to the cap \(3\),
while the other coordinate loses one unit; \(u = 0\) executes no
instrument and both coordinates lose one. The transition is defined index-wise: the acted coordinate regenerates, \(F_j(z)_j = \min\{3,\ z_j + 1\}\), and the other declines, \(F_j(z)_{3-j} = \max\{0,\ z_{3-j} - 1\}\), for \(j \in \{1,2\}\); \(F_0(z)_i = \max\{0,\ z_i - 1\}\) for both coordinates \(i\) --- coordinates clip at zero, the cap is saturating --- and safety is evaluated after each transition. A coordinate-swapped misreading (the regenerated coordinate written first for both \(j\)) reproduces no audited count and is excluded by the verification script. Instrument \(0\) is never uniquely viable on the audited family: admitting \(u = 0\) as a third action leaves all four kernel counts unchanged (\(24, 26, 25, 28\); verified by the verification script, which also checks the structural premises --- \(F_0 \le F_1\) and \(F_0 \le F_2\) pointwise on the grid and the safe set an up-set, the one-step reason no fibre needs \(0\)). The corner \((1,1)\) is the
single state from which no instrument maintains both floors, so its
singleton is nonviable; the nine safe states generate
\(C(9,2) = 36\) two-element belief pairs, the audited population (28 of the pairs full-information-viable; column F of Table~\ref{tab:master}), at review horizon \(H = 12\) under the audited two-horizon verdict (the winning recursion at \(H\) and at \(H - 2\), which agree setwise on the audited universe and on the whole \(2^9\)-belief universe (stabilized from horizon \(2\) there, and from horizon \(1\) for the two non-alternation classes) --- \(W_{12} = W_{10}\) for each policy class, verified by the verification script --- so the verdict certifies a stabilized recursion; Abaee, 2026, An obstruction calculus, Theorem 1).

\emph{Observation structures.} An observation map
\(\mathrm{info}: x \mapsto c\) partitions a belief into fibres; a policy
acts per fibre on the read cell. The audited policy classes combine two
aggregation levels (full information; aggregation
\(c(x) = z_1 + z_2\)) with two protocols (free per-step choice; the
no-repeat protocol that forbids repeating the previous action. The
protocol's register is formalized as follows: one previous-action register
serves each review, shared by all fibres of the belief and not one per
branch; the policy fixes one action per fibre, the register then holds
some fibre's executed action, and the audited rule is pessimistic --- the
belief survives only if every admissible thread (each fibre's action
entering the register) leaves a winning continuation. Register values are
the two instruments, the initial register is unconstrained, and the survival condition quantifies over every
admissible thread, so the semantics is thread-independent by
construction (Proposition~\ref{prop:register} computes what
single-thread rules and per-branch registers give instead).

\emph{Exactness.} Counts are computed by backward recursion over fibre
partitions with exact state arithmetic; witnesses are extracted as step-one per-fibre action vectors along a winning branch; the certificate itself is the recursion invariant --- winning-set membership at every horizon level, quantified universally over adversarial successor states and register threads --- of which the archived step-one vector is the first level. The census of
Section \ref{adequacy} enumerates all \(15\) partitions of a four-state
target; partitions are ordered by refinement --- \(P\) is coarser than
\(P'\) when every fibre of \(P'\) lies in a fibre of \(P\) --- and
``coarsest adequate'' is relative to this order. The continuous objects of Section \ref{benchmark} are rationals.

\emph{Conventions.} On the hold class of the hidden-regime model, the
worst-case exit time is \(\tau^{*}(z_0) = z_0 - 1\) and the discrete
survival supremum (the largest admissible integer review count) is
\(\sigma^{*}(z_0) = \lfloor z_0 - 1 \rfloor = \lfloor \tau^{*}(z_0)
\rfloor\), with level sets \(\{z_0 \ge 1 + k\}\); for an integer deadline
the inclusive criterion \(T_{\mathrm{obs}} \le \tau^{*}(z_0)\) coincides
with \(T_{\mathrm{obs}} \le \sigma^{*}(z_0)\), while the strict form
against the floored count is stricter than against the real exit time off
integer \(\tau^{*}\). Two
certainty-equivalence laws are audited: an uncorrected law whose index
leads the state by one step, \(s \mapsto s + 0.1\), and a corrected law
with the lead removed; both are read through the index \(g(s) = s^{2}\).

\emph{Figures and tables.} Every figure and every tabulated entry in
this paper is produced by a generation script that first recomputes the
quantity from the primitives above, in exact rational arithmetic, and
asserts it before drawing or printing; the master verification script of
Section \ref{methods} chains both scripts and re-derives the tabulated
entries independently.

\medskip
\noindent\textbf{Shared symbols across the companion papers.} The
companion papers share this vocabulary with paper-local meanings; the
letters below carry meanings here that differ there:
\begin{center}\small
\begin{tabular}{@{}p{0.40\linewidth}p{0.52\linewidth}@{}}
\toprule
\textbf{Here} & \textbf{In the companions} \\
\midrule
\(\lambda\), \(\lambda^{*}\): the Farkas and dual certificate weights &
the Farkas multiplier and observer decay rate (obstruction calculus);
certificate weight vectors (certification; minimax duals) \\
\(\tau^{*}(z_0)\): the worst-case exit time &
observation and revelation times \(\tau\) (obstruction calculus;
certification; minimax duals) \\
\(\sigma^{*}(z_0)\): the floored exit count &
the kernel time variable \(\sigma\) (certification; minimax duals);
a noise scale (forecast ladder) \\
\bottomrule
\end{tabular}
\end{center}

\subsection{A master monotonicity theorem}\label{mtheorem}

The audited family is an instance of one abstract object, and its
kernel comparisons are instances of one theorem. An \emph{audited
tuple} is
\(\Theta = (X, \mathcal{V}, A, F, \mathrm{info}, \Pi, m)\): a finite
state set \(X\) with safe set \(\mathcal{V}\), action set \(A\),
transition \(F\), observation map, admissible policy class \(\Pi\)
(per-cell action rules), and register convention \(m\) (shared, per-branch,
or none); viability is the horizon-stabilized recursion of
Section~\ref{frame}.

\begin{proposition}[master monotonicity]\label{prop:master}
Fix the audited tuple. Kernel membership is monotone in each of the
three enlargements, setwise: (i) \emph{policies}: if
\(\Pi \subseteq \Pi'\) then
\(\mathcal{W}_{\Pi} \subseteq \mathcal{W}_{\Pi'}\) --- restriction is
antitone (Proposition~\ref{prop:lattice}); (ii) \emph{observations}:
if \(\mathrm{info}'\) is finer than \(\mathrm{info}\) (every fibre of
\(\mathrm{info}\) is a union of fibres of \(\mathrm{info}'\)) then
\(\mathcal{W}_{\mathrm{info}} \subseteq
\mathcal{W}_{\mathrm{info}'}\); (iii) \emph{memory}: a per-branch
register dominates the shared register --- every shared-register policy
is a per-branch policy, so the per-branch kernel contains the shared
one. Finer observation maps and larger memories can only enlarge the
kernel; all three statements are proved by the same induction on the
horizon: each enlargement enlarges the admissible action sets at every
cell and horizon level, so a winning continuation transfers upward.
\end{proposition}

\begin{proof}
Induction on \(h\). For \(h = 0\) the kernels coincide
(\(\mathcal{V}\) is fixed). Step: a belief winning at \(h\) under the
restricted tuple uses, at each cell, an action admissible and winning
there; the same choice is admissible under the enlarged tuple (the
action set at each cell only grows: more policies choose more actions,
finer cells shrink the quantifier over merged branches, a branch-local
register relaxes the no-repeat constraint on the other branch), and the
continuation is winning at \(h-1\) by the inductive hypothesis. The
finite population makes the induction an exact certificate: checks
P16--P18 verify the containments setwise on all \(36\) pairs.
\end{proof}

The audited counts are the theorem's rows: \(24 \le 26\) and
\(25 \le 28\) (observation axis, free protocol), \(24 \le 25\) and
\(26 \le 28\) (protocol axis at fixed aggregation), and the
per-branch-register kernels \((26, 26, 27)\) dominate the shared ones
\((24, 26, 25)\) (memory axis; Proposition~\ref{prop:register}).
\begin{proposition}[stationary policies]\label{prop:stationary}
Enumerate all deterministic stationary policies --- one action per cell,
fixed in time: \(2^{9} = 512\) at full information (nine singleton
cells), \(2^{5} = 32\) at aggregation (sums two to six). At full
information the union of the stationary policies' safe-pair sets is
exactly the kernel: memoryless sufficiency holds on the instance, and
the \(28\)-pair kernel admits an independent constructive description.
At aggregation the union is \emph{empty} --- no stationary policy
sustains any of the \(36\) pairs: the merged index is read against
states that demand conflicting actions across visits, and every
aggregate-viable pair relies on actions scheduled against time. The
\(26\)-pair aggregate kernel is therefore entirely a time-varying-policy
phenomenon --- the register-coupling mechanism in another guise, with
the review clock in the register's place (checks P17--P18; the fixpoint
computation of each class's horizon family independently reproduces all
four kernels).
\end{proposition}

\begin{proof}
Both counts are exhaustive enumerations with exact forward simulation
under each fixed policy (check P17), cross-validated against the
backward recursion by the horizon fixpoint (check P18).
\end{proof}

Strictness is instance-specific: the strict differences are exactly the
register-blocked pairs \(\{(1,2),(3,1)\}\) and \(\{(1,3),(2,1)\}\) on
the protocol axis, the single pair \(\{(1,3),(3,1)\}\) on the
observation axis, with the alternation-forced pair \(\{(1,2),(2,1)\}\)
outside both middle kernels; the incomparability of the
middle cells is a shared-register phenomenon, and the decentralized
kernel sits outside the chain because unanimity restricts \emph{who}
acts, not what the acting policy may do --- an orthogonal axis the
theorem prices separately (Section~\ref{decentralized}).

\subsection{The master table}\label{master}

Table \ref{tab:master} is the audit's master table: all 36 belief pairs
against the five kernels --- the four policy classes of Section
\ref{classes} (institutional: one shared compliance register per
review; aggregate: one pooled index per branch; full-information codex:
pooled registers across branches; full-information: unrestricted state
observation) and the decentralized kernel of Section
\ref{decentralized}. Rows are grouped by the pair's lexicographically
smaller state, written \(z_1 z_2\).

\begin{table*}[t]
\centering
\footnotesize
\caption{The master table: membership of all 36 belief pairs in the five
audited kernels. I = institutional (aggregation, no-repeat); A =
aggregate; C = full-information codex; F = full information; D =
decentralized unanimity. A dot marks a non-member; every entry is
regenerated by script.}
\label{tab:master}
\begin{tabular}{@{}l cccc c@{\hspace{12pt}}l cccc c@{}}
\toprule
Pair & I & A & C & F & D & Pair & I & A & C & F & D\\
\midrule
$12|21$ & $\cdot$ & $\cdot$ & $\cdot$ & F & D &
$22|31$ & I & A & C & F & D\\
$12|31$ & $\cdot$ & A & $\cdot$ & F & $\cdot$ &
$22|23$ & I & A & C & F & $\cdot$\\
$12|23$ & I & A & C & F & D &
$22|33$ & I & A & C & F & D\\
$12|33$ & I & A & C & F & $\cdot$ &
$22|32$ & I & A & C & F & $\cdot$\\
$12|22$ & I & A & C & F & $\cdot$ &
$23|31$ & I & A & C & F & $\cdot$\\
$12|32$ & I & A & C & F & D &
$23|33$ & I & A & C & F & $\cdot$\\
$12|13$ & I & A & C & F & $\cdot$ &
$23|32$ & I & A & C & F & D\\
$13|21$ & $\cdot$ & A & $\cdot$ & F & $\cdot$ &
$31|33$ & I & A & C & F & D\\
$13|31$ & $\cdot$ & $\cdot$ & C & F & D &
$31|32$ & I & A & C & F & $\cdot$\\
$13|23$ & I & A & C & F & $\cdot$ &
$32|33$ & I & A & C & F & $\cdot$\\
$13|33$ & I & A & C & F & D &
$11|12$ & $\cdot$ & $\cdot$ & $\cdot$ & $\cdot$ & $\cdot$\\
$13|22$ & I & A & C & F & D &
$11|21$ & $\cdot$ & $\cdot$ & $\cdot$ & $\cdot$ & $\cdot$\\
$13|32$ & I & A & C & F & $\cdot$ &
$11|31$ & $\cdot$ & $\cdot$ & $\cdot$ & $\cdot$ & $\cdot$\\
$21|31$ & I & A & C & F & $\cdot$ &
$11|23$ & $\cdot$ & $\cdot$ & $\cdot$ & $\cdot$ & $\cdot$\\
$21|23$ & I & A & C & F & D &
$11|33$ & $\cdot$ & $\cdot$ & $\cdot$ & $\cdot$ & $\cdot$\\
$21|33$ & I & A & C & F & $\cdot$ &
$11|22$ & $\cdot$ & $\cdot$ & $\cdot$ & $\cdot$ & $\cdot$\\
$21|22$ & I & A & C & F & $\cdot$ &
$11|32$ & $\cdot$ & $\cdot$ & $\cdot$ & $\cdot$ & $\cdot$\\
$21|32$ & I & A & C & F & D &
$11|13$ & $\cdot$ & $\cdot$ & $\cdot$ & $\cdot$ & $\cdot$\\
\bottomrule
\end{tabular}
\end{table*}

The column counts are the audited kernel sizes
\(|\mathcal{W}_{\mathrm{inst}}| = 24\),
\(|\mathcal{W}_{\mathrm{agg}}| = 26\),
\(|\mathcal{W}_{\mathrm{full\text{-}codex}}| = 25\),
\(|\mathcal{W}_{\mathrm{full}}| = 28\), and
\(|\mathcal{W}_{\mathrm{dec}}| = 12\). Figure \ref{fig:kernels} displays
the four policy-class kernels, whose middle cells are incomparable. Two set-level facts are recorded exactly. The meet of the
two middle kernels equals the institutional kernel setwise. At class
level the join is the full-information class; at kernel level the union
of the two middle kernels falls one pair short of the full kernel:
exactly the pair \(\{(1,2),(2,1)\}\) lies in neither middle kernel. The
pair is the alternation-forced pair: under full information each branch
survives by its own strict alternation --- \((1,2)\) under \(1, 2, 1,
\dots\) and \((2,1)\) under \(2, 1, 2, \dots\) --- which the
idealized per-state policy can carry on both branches; aggregation cannot
see either requirement (one fibre, no common action), and the no-repeat
protocol, whose single previous-action register must serve both branches,
forbids the second branch's required next action once the first branch's
choice occupies the register. The full kernel is exactly the
\(C(8,2) = 28\) pairs avoiding the corner \((1,1)\): the corner's
singleton nonviability accounts for every remaining pair-level loss at
full information, as the eight \(11|\cdot\) rows of Table
\ref{tab:master} show.

\begin{figure}[t]
\centering
\includegraphics[width=0.94\linewidth]{figs_ws4/fig_kernels.pdf}
\caption{The four policy classes form a two-by-two product of aggregation
level and protocol; the audited kernel sizes are \(24, 26, 25, 28\) of
36. The meet of the middle kernels is the institutional kernel setwise;
the class join is the full-information class, whose kernel is the 28
corner-free pairs, one pair beyond the middle-kernel union.}
\label{fig:kernels}
\end{figure}

\subsection{Policy classes}\label{classes}

The four audited classes combine full information and aggregation with
free and no-repeat protocols. Their kernels among the 36 pairs have
sizes:
\begin{gather*}
|\mathcal{W}_{\mathrm{inst}}| = 24,\quad
|\mathcal{W}_{\mathrm{agg}}| = 26,\\
|\mathcal{W}_{\mathrm{full\text{-}codex}}| = 25,\quad
|\mathcal{W}_{\mathrm{full}}| = 28 .
\end{gather*}
The four classes form a two-by-two product of aggregation level and
protocol; their kernels have incomparable middle cells:
the aggregate kernel contains pairs outside the full-information codex
kernel (for instance \(\{(1,2),(3,1)\}\) and \(\{(1,3),(2,1)\}\)), and
the codex kernel contains pairs outside the aggregate (for instance
\(\{(1,3),(3,1)\}\)). Within each observation column the free protocol's
kernel contains the no-repeat protocol's (both inclusions strict); across
columns neither middle kernel contains the other. The precise
meet-and-join reading is the one recorded in Section \ref{master}. Two
per-pair readings illustrate the columns of Table \ref{tab:master}. The
pair \(\{(1,2),(3,1)\}\) is aggregate-viable but not codex-viable: the
free protocol admits it under both observation structures, while the
no-repeat protocol excludes it under both --- at the second review the
thread that records the \((3,1)\)-branch's instrument \(2\) forbids the
forced instrument \(2\) of the successor branch \((2,1)\), the same
shared-register collision that excludes the alternation-forced pair. The pair \(\{(1,3),(3,1)\}\) is
codex-viable but not aggregate-viable: the two branches share the fibre
\(c = 4\), and no single per-fibre instrument serves both, while
full information separates them. Three structural facts behind the table
admit short proofs.

\begin{proposition}[full information decomposes branchwise]\label{prop:decomp}
At full information under the free protocol, a two-state belief is
viable if and only if each of its singletons is viable; the kernel is
exactly the \(C(8,2) = 28\) corner-free pairs.
\end{proposition}

\begin{proof}
Full-information fibres are singletons, so a belief's step-one action
sets are the product of its branches': a choice is admissible exactly
when each branch's part is, and the recursion decomposes --- a belief is
viable exactly when both branches are. For the singletons: the corner
\((1,1)\) is nonviable because each instrument drops a coordinate to
zero in one step (\(u = 1\) forces \(z_2 = 0\), \(u = 2\) forces
\(z_1 = 0\), \(u = 0\) forces both); every other singleton is viable
--- the memoryless policy that plays instrument \(1\) exactly at
\(z_1 = 1\) and instrument \(2\) otherwise is safe forever: at
\(z_1 = 1\) the exclusion of the corner gives \(z_2 \ge 2\), so
\(u = 1\) restores \(z_1 = 2\) with \(z_2 \ge 1\); at
\(z_1 \ge 2\), \(u = 2\) keeps \(z_1 \ge 1\) and regenerates
\(z_2\).
\end{proof}

\begin{proposition}[class lattice]\label{prop:lattice}
Kernel membership is antitone in the policy class: restricting the class
can only lose viable beliefs. Hence
\(\mathcal{W}_{\mathrm{inst}} \subseteq \mathcal{W}_{\mathrm{agg}},
\mathcal{W}_{\mathrm{full\text{-}codex}} \subseteq
\mathcal{W}_{\mathrm{full}}\); on the audited 36-pair population the
two middle inclusions are strict, the meet
\(\mathcal{W}_{\mathrm{agg}} \cap
\mathcal{W}_{\mathrm{full\text{-}codex}}\) equals the institutional
kernel setwise, and exactly one pair --- the alternation-forced pair
\(\{(1,2),(2,1)\}\) --- lies in the full kernel outside the middle
union.
\end{proposition}

\begin{proof}
Antitonicity --- order-reversal between class inclusion and kernel inclusion --- is immediate: a policy admissible for a restricted class is
admissible for any enlargements, so viability transfers upward. The
setwise facts are finite statements about a 36-element population; they
are proved by exhaustive certification --- the master table's
\(36 \times 5\) membership marks are recomputed row-by-row by the paper's verification script (checks M2 and M5), which is a complete
proof for a finite family.
\end{proof}

\begin{proposition}[register-semantics invariance]\label{prop:register}
Endow the no-repeat protocol with the pessimistic shared-register
semantics: the policy fixes one action per fibre, the register then holds
some fibre's executed action, and the belief survives only if every
admissible thread leaves a winning continuation. Because the survival condition
quantifies universally over threads, the
pessimistic kernels depend on no thread choice: the audited counts
\(24, 26, 25, 28\) are thread-independent by construction. Two further
readings are computed for the record. A single-thread rule (one fixed
fibre's action entering the register) is order-sensitive: with the
first, or the last, fibre in sorted order the counts are
\(25, 26, 26, 28\) --- neither fixed rule yields the audited
institutional and codex counts --- so the well-defined semantics is the
pessimistic one. The optimistic semantics (some thread) is strictly
larger (codex \(27\)) and is not the audited reading. If instead each
branch carries its own register --- one compliance record per unit,
threaded forward along the branch --- the protocol cost collapses: the
institutional kernel equals the aggregate kernel setwise (\(26 = 26\)),
and the codex kernel becomes \(27\), regaining the alternation-forced
pair and losing exactly \(\{(1,3),(3,1)\}\) relative to full
information, with the aggregate kernel contained in it. The
incomparability of the middle cells is therefore a shared-register
phenomenon; the shared register is the audited institutional object ---
one compliance register per review across the units managed under the
protocol.
\end{proposition}

\begin{proof}
Every clause is a finite statement, computed from the shared primitives
of Section \ref{frame} and asserted by the verification script: the
pessimistic counts (M2b), the optimistic containment (M2c), both
single-thread orders, and the per-branch kernels with the
\(\{(1,3),(3,1)\}\) loss (checks P9 and P10).
\end{proof}

\subsection{Decentralized observation}\label{decentralized}

Two agencies each observe one coordinate and vote; an instrument
executes only on agreement. Each agency's vote is a map from its
coordinate's four values to the two instruments, so the law space has
\(2^{4} \times 2^{4} = 256\) law pairs. Unanimity sustains 12 of the
36 pairs (column D of Table \ref{tab:master}), and exactly 16
coordinated-viable pairs are lost. The loss instance
\(\{(1,2),(3,1)\}\) is representative: each branch is individually
sustainable, and no law pair sustains the belief. For the audited belief
\(\{(1,2),(2,1)\}\), exactly 16 of the 256 law pairs sustain it; the
count coincides numerically with the 16 lost pairs, and no identity is
implied. The authority restriction is exact:

\begin{proposition}[authority restriction]\label{prop:authority}
Fixing agency 2's vote at \(z_2 = 1\) to instrument \(1\) sustains no
law pair for the belief \(\{(1,2),(2,1)\}\): branch \((2,1)\) exits
at its first step under every such law pair. Fixing agency 1's vote at
\(z_1 = 1\) to instrument \(1\) leaves exactly the \(16\) sustaining
law pairs of Table \ref{tab:laws}, and either whole-law
restriction (one agency's vote fixed to a single instrument at every
level) collapses the decentralized kernel to zero.
\end{proposition}

\begin{proof}
On branch \((2,1)\) at the first step: \(z_2 = 1\) must not
decline, so agreement on instrument \(2\) is required; agency 2's
fixed vote at \(z_2 = 1\) is \(1\), so agreement can hold only on
instrument \(1\), which drives \(z_2\) to zero, and disagreement
executes \(u = 0\), which does the same --- the branch exits at step
one under every admissible law pair. With agency 1's vote at \(z_1 = 1\) fixed to \(1\), the
sustaining pairs are exactly those with agency 1 forced to \(1\) at
\(z_1 = 1\) and \(2\) at \(z_1 = 2\) and agency 2 forced to \(2\)
at \(z_2 = 1\) and \(1\) at \(z_2 = 2\) (each forcing is a
one-step safety requirement: the low coordinate must regenerate), and
the free votes at \(\{0, 3\}\) never bind --- the \(4 \times 4\)
product of Table \ref{tab:laws}. A whole-law restriction
fixes one agency's vote everywhere; the other agency's coordinate then
cannot regenerate at the low states where the fixed vote disagrees, and
the kernel collapses (verified exhaustively over all \(256\) law pairs, the whole-law
kernels \(0\) in all four cases; check P12 of the verification script).
\end{proof}

\begin{table}[t]
\centering
\footnotesize
\caption{The 16 law pairs sustaining the audited belief
\(\{(1,2),(2,1)\}\) factor exactly: agency 1's vote is forced to
\(1\) at \(z_1 = 1\) and \(2\) at \(z_1 = 2\), free at
\(z_1 \in \{0, 3\}\); agency 2's vote is forced to \(2\) at
\(z_2 = 1\) and \(1\) at \(z_2 = 2\), free at
\(z_2 \in \{0, 3\}\). Every combination sustains --- a complete
\(4 \times 4\) product (votes listed as
\((l(z{=}0), l(z{=}1), l(z{=}2), l(z{=}3))\));
agency 1's votes are the row labels, agency 2's the column labels.}
\label{tab:laws}
\begin{tabular}{@{}l cccc@{}}
\toprule
 & \multicolumn{4}{c}{agency 2}\\
\cmidrule(l){2-5}
agency 1 & 1211 & 1212 & 2211 & 2212\\
\midrule
1121 & \checkmark & \checkmark & \checkmark & \checkmark\\
1122 & \checkmark & \checkmark & \checkmark & \checkmark\\
2121 & \checkmark & \checkmark & \checkmark & \checkmark\\
2122 & \checkmark & \checkmark & \checkmark & \checkmark\\
\bottomrule
\end{tabular}
\end{table}

The law space itself is counted before quotienting: a \(z = 0\) vote is
consultable only on a trajectory that has already left the safe set, so
the audited kernel is invariant under the \(z{=}0\)-vote quotient ---
the effective law space is \(2^{3} \times 2^{3} = 64\) behaviorally
distinct law pairs, and the \(16\) sustaining law pairs of Table
\ref{tab:laws} comprise \(4\) distinct behavioral laws times the \(4\)
never-consulted \(z{=}0\) votes. The sustaining set has an
exact product structure (Table
\ref{tab:laws}): on the trajectory alternating between the branches,
agency 1 is read only at \(z_1 \in \{1, 2\}\) and must vote \(1\) then
\(2\); agency 2 is read only at \(z_2 \in \{2, 1\}\) and must vote
\(1\) then \(2\); the votes at the unreachable coordinates
\(z \in \{0, 3\}\) are free, giving \(2^{2} \times 2^{2} = 16\). The
restricted kernels are exact. Fixing either agency's low vote to the
regenerating instrument (agency 1 to \(1\) at \(z_1 = 1\), agency 2 to
\(2\) at \(z_2 = 1\)) leaves the decentralized kernel untouched at
\(12\) --- the constraint does not bind. Fixing either agency's low
vote against it (agency 2 to \(1\) at \(z_2 = 1\), agency 1 to \(2\)
at \(z_1 = 1\)) shrinks the kernel to four beliefs ---
\(\{(2,3),(3,2)\}\), \(\{(2,2),(3,3)\}\), \(\{(1,3),(3,3)\}\),
\(\{(1,3),(2,2)\}\), up to the symmetry below --- and kills the
audited target; a whole-law restriction collapses the kernel to \(0\)
in all four cases. The model's swap symmetry (coordinates and instruments
together) interchanges the agencies and maps the audited target to
itself, so no agency is privileged: what carries the target is the
direction of each forced vote --- each branch's low coordinate must be
allowed to regenerate --- not the agency holding it. Two modelling choices are declared. The deadlock transition
\(u = 0\) executes no instrument and both coordinates lose one; and
each agency votes on its own coordinate's current value: the audited
protocol is memoryless and decentralized --- no agency observes the
belief, the other agency's observation, or anything like common
knowledge. The count has a theorem behind it:

\begin{proposition}[the decentralized kernel factors statewise]
\label{prop:req}
Unanimity executes statewise: the executed instrument at a state depends
only on the two votes at that state's coordinates, and safety is
evaluated per state. Hence a belief is sustained if and only if its
branches share a common sustaining law pair,
\(\bigcap_{x \in B} S(x) \neq \emptyset\), where \(S(x)\) is the set
of law pairs sustaining the singleton \(\{x\}\); on the audited
population the factored criterion reproduces the joint kernel exactly
(the \(12\) sustained beliefs of the master table), and the \(16\)
sustaining law pairs of the audited belief
(Table~\ref{tab:laws}) are invariant under the model's coordinate-swap
symmetry. The obstruction is legible branchwise: a belief is lost
exactly when its branches' sustaining sets do not intersect --- as for
the representative loss \(\{(1,2),(3,1)\}\), each branch individually
sustainable, no law pair sustaining both.
\end{proposition}

\begin{proof}
The belief trajectory under a law pair is the union of the branch
trajectories (statewise execution, statewise safety), so it stays in
\(\mathcal{V}\) if and only if every branch trajectory does --- which is
the displayed intersection. Both computations (the factored singleton
intersection and the joint \(256\)-pair enumeration) are run on all
\(36\) beliefs and agree setwise (check P20). The symmetry is the swap
map of this section's closing remark applied to the product structure.
\end{proof}

\subsection{Review timing on the hidden-regime grid}\label{timing}

The hidden-regime resource model: a single stock held at level
\(z_0 \ge 1\) above the floor; each review period the unobserved
adversarial regime realizes one unit of decline and there is no
regeneration (the hold class), so the stock tracks \(z(t) = z_0 - t\)
and remains safe exactly while \(z(t) \ge 1\). The model is audited on
the grid \(z_0 \in \{1.0, 1.1, \dots, 4.0\}\) and
\(T_{\mathrm{obs}} \in \{1, 2, 3\}\), 93 cells. Viability at the
deadline holds if and only if
\(T_{\mathrm{obs}} \le \tau^{*}(z_0) = z_0 - 1\), with zero
mismatches across the grid. The identity is boundary-inclusive: the
viable cells number \(33\), and the strict form \(T_{\mathrm{obs}} <
\tau^{*}(z_0)\) fails exactly on the boundary cells \((2,1)\),
\((3,2)\), and \((4,3)\) --- those with \(z_0 = 1 +
T_{\mathrm{obs}}\); against the floored count \(\sigma^{*}\) the
strict form fails on one cell per row \(z_0 \ge 2\) --- \(21\) cells,
the boundary cells included --- which is why the
inclusive predicate is the audited identity.
Table \ref{tab:timing} carries the full grid and Figure
\ref{fig:timing} displays it. The practical form of the rule: a review
must land no later than the worst-case exit time, and a deadline exactly at the worst-case exit is admissible under the audited endpoint convention (closed safe set; safety evaluated at the deadline review); no open-set or alternative-exit-convention claim is made. The audited identity is in fact the
restriction of a general law:

\begin{proposition}[review timing, general form]\label{prop:timing}
On the hold class of the hidden-regime model, for every initial level
\(z_0 \ge 1\) and every integer deadline \(T \ge 0\), the cell
\((z_0, T)\) is deadline-viable if and only if \(T \le z_0 - 1\).
Accordingly \(\sigma^{*}(z_0) = \lfloor z_0 - 1 \rfloor\) with
superlevel sets \(\{z_0 \ge 1 + k\}\), \(k \in
\mathbb{Z}_{\ge 0}\), and the audited 93-cell grid is the restriction
to \(z_0 \in \{1.0, \dots, 4.0\}\), \(T \in \{1, 2, 3\}\).
\end{proposition}

\begin{proof}
Under every admissible hold-class control the adversarial regime realizes
the declining branch at each step, so the worst stock declines at unit
rate and remains safe exactly through time \(z_0 - 1\) --- the last
safe review, the first violation occurring at \(z_0\):
deadline-viability forces \(T \le z_0 - 1\). Conversely
the hold policy tracks the decline path \(z(t) = z_0 - t\) branchwise,
which stays admissible exactly on \([0,\, z_0 - 1]\) and terminates at
the floor, which is safe because the safe set is closed; hence
\(T \le z_0 - 1\) is viable, boundary included. The set identity
\(\{T : T \le z_0 - 1\} = \{T : T \le \sigma^{*}(z_0)\}\)
(integer \(T\), all \(z_0\)) --- equivalently \(\{T < z_0 - 1\} =
\{T \le \sigma^{*}(z_0)\}\) off the integers, with the boundary
failure at integer \(z_0 - 1\) --- gives the displayed supremum.
\end{proof}

\begin{proposition}[hitting times on regime graphs]\label{prop:regime}
Let the adversary choose a regime path \(r_1 r_2 \dots\) through a
finite regime graph with per-step damages \(d_r \ge 0\), subject to
the graph's adjacency constraint. The worst cumulative damage by time
\(t\) is the max-plus recursion
\(D_t = \max_{\text{paths}} \sum_{s \le t} d_{r_s}\) --- the standard
timed-event-system algebra (Baccelli et al., 1992) --- and the cell
\((z_0, T)\) is deadline-viable if and only if \(D_T \le z_0 - 1\)
(the audited endpoint convention). Two exact instances: with free
switching the adversary plays the worst regime every step, \(D_t = t
\cdot \max_r d_r\), recovering Proposition~\ref{prop:timing} on the
audited grid; with two regimes of damages \(1\) and \(1/2\) forced to
alternate, the worst damage is \(D_{2k} = 3k/2\) and \(D_{2k+1} = 3k/2
+ 1\), and the re-audited \(93\)-cell grid has exactly \(43\) viable
cells (\(21, 16, 6\) at \(T = 1, 2, 3\)), zero mismatches between the
closed form and path enumeration. The max-plus dynamic program equals
path enumeration exactly on a three-regime graph with a forbidden
self-loop (all horizons to six).
\end{proposition}

\begin{proof}
Deadline-viability is safety of the worst adversarial path to the
deadline, which is the max-plus evaluation of the damage sum; free
switching decouples the steps, giving \(t \cdot \max_r d_r = t\) on
the audited damages and hence the audited identity. Under alternation
the two interleavings give \(D_{2k} = k(1 + \tfrac12) = 3k/2\) and
\(D_{2k+1} = (k+1) + \tfrac{k}{2}\); the grid recomputation enumerates
both interleavings per cell (check P21). The dynamic-programming (DP) evaluation over
(last regime, step) is exact rational arithmetic and is asserted equal
to explicit path enumeration.
\end{proof}

\begin{table}[t]
\centering
\footnotesize
\caption{The timing grid, printed in full (\(93\) cells): viability at the deadline against \(\sigma^{*}(z_0) = \lfloor
\tau^{*}(z_0) \rfloor\) (the deadline count; the identity columns use
the inclusive predicate, on which the floored and real-valued forms
coincide). The checkmark of column
\(T_{\mathrm{obs}}\) means the cell is viable; the identity
column ``id.'' is the predicate
\(T_{\mathrm{obs}} \le \sigma^{*}(z_0)\) and column ``cell'' the audited
viability; the two agree on all \(93\) cells (zero mismatches; \(33\)
viable cells). The circled cell of Figure \ref{fig:timing} is
\((2.0, 1)\): viable, on the boundary; the strict form disagrees
exactly on the boundary cells \((2,1)\), \((3,2)\), \((4,3)\).}
\label{tab:timing}
\begin{tabular}{@{}l c c cc c cc c@{}}
\toprule
 & & \multicolumn{2}{c}{$T_{\mathrm{obs}} = 1$} &
\multicolumn{2}{c}{$T_{\mathrm{obs}} = 2$} &
\multicolumn{2}{c}{$T_{\mathrm{obs}} = 3$}\\
\cmidrule(lr){3-4}\cmidrule(lr){5-6}\cmidrule(lr){7-8}
$z_0$ & $\sigma^{*}$ & id. & cell & id. & cell & id. & cell\\
\midrule
1.0 & 0 & $\cdot$ & $\cdot$ & $\cdot$ & $\cdot$ & $\cdot$ & $\cdot$\\
1.1 & 0 & $\cdot$ & $\cdot$ & $\cdot$ & $\cdot$ & $\cdot$ & $\cdot$\\
1.2 & 0 & $\cdot$ & $\cdot$ & $\cdot$ & $\cdot$ & $\cdot$ & $\cdot$\\
1.3 & 0 & $\cdot$ & $\cdot$ & $\cdot$ & $\cdot$ & $\cdot$ & $\cdot$\\
1.4 & 0 & $\cdot$ & $\cdot$ & $\cdot$ & $\cdot$ & $\cdot$ & $\cdot$\\
1.5 & 0 & $\cdot$ & $\cdot$ & $\cdot$ & $\cdot$ & $\cdot$ & $\cdot$\\
1.6 & 0 & $\cdot$ & $\cdot$ & $\cdot$ & $\cdot$ & $\cdot$ & $\cdot$\\
1.7 & 0 & $\cdot$ & $\cdot$ & $\cdot$ & $\cdot$ & $\cdot$ & $\cdot$\\
1.8 & 0 & $\cdot$ & $\cdot$ & $\cdot$ & $\cdot$ & $\cdot$ & $\cdot$\\
1.9 & 0 & $\cdot$ & $\cdot$ & $\cdot$ & $\cdot$ & $\cdot$ & $\cdot$\\
2.0 & 1 & \checkmark & \checkmark & $\cdot$ & $\cdot$ & $\cdot$ & $\cdot$\\
2.1 & 1 & \checkmark & \checkmark & $\cdot$ & $\cdot$ & $\cdot$ & $\cdot$\\
2.2 & 1 & \checkmark & \checkmark & $\cdot$ & $\cdot$ & $\cdot$ & $\cdot$\\
2.3 & 1 & \checkmark & \checkmark & $\cdot$ & $\cdot$ & $\cdot$ & $\cdot$\\
2.4 & 1 & \checkmark & \checkmark & $\cdot$ & $\cdot$ & $\cdot$ & $\cdot$\\
2.5 & 1 & \checkmark & \checkmark & $\cdot$ & $\cdot$ & $\cdot$ & $\cdot$\\
\midrule
2.6 & 1 & \checkmark & $\cdot$ & $\cdot$\\
2.7 & 1 & \checkmark & $\cdot$ & $\cdot$\\
2.8 & 1 & \checkmark & $\cdot$ & $\cdot$\\
2.9 & 1 & \checkmark & $\cdot$ & $\cdot$\\
3.0 & 2 & \checkmark & \checkmark & $\cdot$\\
3.1 & 2 & \checkmark & \checkmark & $\cdot$\\
3.2 & 2 & \checkmark & \checkmark & $\cdot$\\
3.3 & 2 & \checkmark & \checkmark & $\cdot$\\
3.4 & 2 & \checkmark & \checkmark & $\cdot$\\
3.5 & 2 & \checkmark & \checkmark & $\cdot$\\
3.6 & 2 & \checkmark & \checkmark & $\cdot$\\
3.7 & 2 & \checkmark & \checkmark & $\cdot$\\
3.8 & 2 & \checkmark & \checkmark & $\cdot$\\
3.9 & 2 & \checkmark & \checkmark & $\cdot$\\
4.0 & 3 & \checkmark & \checkmark & \checkmark\\
\bottomrule
\end{tabular}
\end{table}

\begin{figure}[t]
\centering
\includegraphics[width=0.96\linewidth]{figs_ws4/fig_timing.pdf}
\caption{The timing grid. Cells viable at the deadline (green) satisfy \(T_{\mathrm{obs}} \le \tau^{*}(z_0)\); the identity is
boundary-inclusive, and the strict form \(T_{\mathrm{obs}} <
\tau^{*}(z_0)\) fails at exactly one cell,
\((z_0, T_{\mathrm{obs}}) = (2.0, 1)\) (circled); against the floored count \(\sigma^{*}(z_0) = \lfloor z_0 - 1 \rfloor\) the strict form fails on all six viable cells, which is why the inclusive predicate is the audited identity.}
\label{fig:timing}
\end{figure}

\subsection{Monitoring adequacy}\label{adequacy}

Let \(R(x)\) denote the safe actions of state \(x\): those keeping
\(x\) in \(\mathcal{V}\) for one step. On the audited target,
\(R((1,2)) = \{1\}\), \(R((2,1)) = \{2\}\), \(R((2,2)) = \{1,2\}\), and
\(R((3,1)) = \{2\}\). A partition of a target belief is \emph{adequate}
when every fibre has a nonempty common safe-action set. On the
four-state target \(\{(1,2), (2,1), (2,2), (3,1)\}\), 7 of the 15
partitions are adequate, and two minimal adequate partitions with two
cells exist. Table \ref{tab:census} carries the complete census with
each fibre's common safe-action set; Figure \ref{fig:census} displays
it.

\begin{table}[t]
\centering
\footnotesize
\caption{The fifteen-partition census on the target
\(\{(1,2), (2,1), (2,2), (3,1)\}\) (states written \(z_1 z_2\); fibres
separated by ``;''). A partition is adequate when every fibre's common
safe-action set (brackets; \(\cdot\) = empty) is nonempty. Seven
partitions are adequate; the minimal adequate two-cell partitions are
rows 6 and 7.}
\label{tab:census}
\begin{tabular}{@{}l l l@{}}
\toprule
Partition & Adequate & Common safe actions\\
\midrule
$12|21|22|31$ & no & $[\cdot]$\\
$12|21|31$ ; $22$ & no & $[\cdot\ 12]$\\
$12|22|31$ ; $21$ & no & $[\cdot\ 2]$\\
$12|31$ ; $21|22$ & no & $[\cdot\ 2]$\\
$12|31$ ; $22$ ; $21$ & no & $[\cdot\ 12\ 2]$\\
$21|22|31$ ; $12$ & \textbf{yes} & $[2\ 1]$\\
$21|31$ ; $12|22$ & \textbf{yes} & $[2\ 1]$\\
$21|31$ ; $22$ ; $12$ & yes & $[2\ 12\ 1]$\\
$22|31$ ; $12|21$ & no & $[2\ \cdot]$\\
$22|31$ ; $21$ ; $12$ & yes & $[2\ 2\ 1]$\\
$31$ ; $12|21|22$ & no & $[2\ \cdot]$\\
$31$ ; $12|22$ ; $21$ & yes & $[2\ 1\ 2]$\\
$31$ ; $21|22$ ; $12$ & yes & $[2\ 2\ 1]$\\
$31$ ; $22$ ; $12|21$ & no & $[2\ 12\ \cdot]$\\
$31$ ; $22$ ; $21$ ; $12$ & yes & $[2\ 12\ 2\ 1]$\\
\bottomrule
\end{tabular}
\end{table}

\begin{figure}[t]
\centering
\includegraphics[width=0.98\linewidth]{figs_ws4/fig_census.pdf}
\caption{The census. Each panel draws one partition: a bar over a group
of states is a fibre (green when its common safe-action set is
nonempty, red when empty), and the margin marks adequacy
(\(+\)/\(\times\)). The two adequate two-cell partitions are panels 6
and 7 of the middle row.}
\label{fig:census}
\end{figure}

The maximal subsets with a nonempty common safe-action set are
\(\{(1,2), (2,2)\}\) and \(\{(2,1), (2,2), (3,1)\}\). They overlap in
\((2,2)\), so no partition into maximal common-action sets exists:
coarsening to maximal fibres is not an available design, and the
adequate partitions found by the census are the design options. The criterion is one-step
monitoring adequacy; dynamic viability is not claimed from it (the
kernels are computed by the full recursion). The census excludes
instrument \(0\) without loss: \(F_0 \le F_1\) and \(F_0 \le F_2\)
pointwise with the safe set an up-set (verified in the frame section)
give every one-step-safe fibre a safe instrument in \(\{1,2\}\). Two
consequences follow. Pairwise intersection of safe-action sets is necessary for fibre
adequacy; in the two-instrument audited universe it is also sufficient
(Helly number two: a pairwise-intersecting family of subsets of a
two-point action set has a common action), so every inadequate fibre in
the census fails with a literal empty intersection, while in action sets
of three or more elements pairwise intersection ceases to suffice
(\(\{1,2\}, \{2,3\}, \{1,3\}\)); and an
observation map merging states with disjoint safe-action sets obstructs
at that information state while every branch remains individually
viable --- the remedy is a fibre split, not a finer threshold on the
same index. The census makes both visible at the level of individual
fibres: in each inadequate partition of Table \ref{tab:census} the
failing fibre is exactly one merging \((1,2)\) --- whose safe-action
set is the singleton \(\{1\}\) --- with a state whose safe actions
exclude \(1\), i.e. \((2,1)\) or \((3,1)\), and each adequate
partition merges \((1,2)\) with neither. The census is itself an
instance of a structure:

\begin{proposition}[adequacy as an incompatibility-graph criterion]
\label{prop:adequacy}
Let \(G\) be the incompatibility graph of the target: vertices are
states, with an edge \((x, x')\) when \(R(x) \cap R(x') =
\emptyset\). A partition is adequate if and only if every fibre is an
independent set of \(G\); on the audited target the edges are exactly
\(\{(1,2)\text{--}(2,1),\ (1,2)\text{--}(3,1)\}\), so the adequate
partitions are the partitions into independent sets, counted by
inclusion--exclusion as \(15 - (5 + 5 - 2) = 7\) (the two edge events
overlap in exactly the \(2\) partitions whose fibre contains
\((1,2), (2,1), (3,1)\) jointly), and the adequate two-cell partitions
are exactly the graph's two proper \(2\)-colorings --- rows 6 and 7.
In an \(a\)-action universe, a fibre fails adequacy only through a
subfamily of at most \(a\) states with empty common intersection: the
Helly number of the family of subsets of an \(a\)-element action set
is \(a\) (verified exhaustively for \(a = 2\) and \(a = 3\), where the
minimal failing family \(\{\{1,2\}, \{2,3\}, \{1,3\}\}\) attains it);
pairwise intersection suffices only for \(a = 2\), and for convex
\(m\)-dimensional action domains the bound is \(m + 1\)
(Proposition~\ref{prop:helly}).
\end{proposition}

\begin{proof}
A fibre is adequate exactly when its states' safe-action sets intersect,
which is the complement of an edge exactly on two-state fibres and,
since the audited universe has Helly number two, for all fibres
(Proposition~\ref{prop:helly} is the \(m\)-dimensional analogue). The
edge count is the two disjointness facts above. The inclusion--exclusion
count: partitions whose fibres include the edge \((1,2)\text{--}(2,1)\)
are partitions of the three objects \(\{(1,2)(2,1)\}, (2,2), (3,1)\) ---
five; symmetrically five for the other edge; both edges force the
fibre \(\{(1,2), (2,1), (3,1)\}\) --- two partitions (with and without
\((2,2)\) merged). A two-cell adequate partition is a bipartition into
independent sets, i.e. a proper \(2\)-coloring of the star; there are
exactly two. The Helly statements are finite computations over all
subfamilies of the \(2^{a}\) subsets of an \(a\)-element set (exact for
\(a = 2, 3\); check P19).
\end{proof}

\subsection{Regime uncertainty}\label{regimes}

Two regimes alternate: an adversary chooses at the end of step
\(t - 1\) the regime realized at step \(t\). The model, explicit: the
controller acts one instrument \(u \in \{1,2\}\) per step (coordinate
\(u\) regenerates by one, cap \(3\)), and the acting regime deals one
unit of damage to its own coordinate (floor \(0\)); there is no other
change. With the regime frozen and announced, the controller drives the state
to a fixed point --- acting the damaged coordinate holds it wherever
its level is below the cap, and one step of the other instrument frees
a capped coordinate --- so the corner \((1,1)\) is
viable under each regime separately, and each frozen kernel comprises
all \(36\) belief pairs. Mode-blind operation, with the regime read
only through the state, yields a kernel of exactly 28 --- every pair
containing the corner is lost (from \((1,1)\) either instrument
regenerates one coordinate while the adversarial damage zeroes the
other), no merged-pair loss exists, and the loss is generated entirely
at the singleton level. The mode-blind kernel
agrees setwise with the full-information kernel of Section
\ref{classes}; the agreement is coincidental, arising under different
damage routings, and no identity is claimed. All three counts (frozen
\(36\) per regime, corner viable under each, mode-blind \(28\)) are
machine-verified from these primitives (check P13 of the verification script).

\subsection{The continuous benchmark}\label{benchmark}

A state is a pair of stocks \((x_1, x_2) \ge 0\) with a demand floor
on each; the aggregate is \(Y = x_1 + x_2\), and the \(Y\)-fibre is
the set of states with that aggregate. One control set
\(U = \{u \ge 0 : u_1 + u_2 \ge 2\}\) faces the floors through
Y-dependent caps \(\mathrm{cap}_1(Y) = 3/2 - (Y - 2)/10\) and
\(\mathrm{cap}_2(Y) = 59/50 - (Y - 2)/10\); an allocation \(u \in U\)
is feasible on the fibre when \(u_i \le \mathrm{cap}_i(Y)\) --- fibre
feasibility, not transition dynamics, is this section's audited object.
The critical aggregate is \(Y^{*} = 27/5\): the cap sum equals \(2\)
exactly there. Table
\ref{tab:bench} evaluates the caps on the audited aggregates and Figure
\ref{fig:bench} displays them. The \(Y = 5\) fibre is viable with
witness \((6/5, 4/5)\); the \(Y = 6\) fibre is obstructed with cap sum
\(47/25 < 2\), certified by the dual measure \((1/2, 1/2)\) on the two
active-floor atoms with margin \(3/50\). The crossover is again a general law:

\begin{proposition}[exact benchmark crossover]\label{prop:bench}
For every aggregate \(Y \ge 2\) with \(\mathrm{cap}_i(Y) \ge 0\), the
\(Y\)-fibre admits a feasible input (some \(u \in U\) with
\(u_i \le \mathrm{cap}_i(Y)\)) if and only if
\(\mathrm{cap}_1(Y) + \mathrm{cap}_2(Y) \ge 2\), which holds if and
only if \(Y \le 27/5\). For \(Y \ge 27/5\) the
\(\ell_1\)-normalized dual measure \((1/2, 1/2)\) on the two cap rows
certifies the obstruction with the exact margin
\(\tfrac12\bigl(2 - \mathrm{cap}_1(Y) - \mathrm{cap}_2(Y)\bigr) =
(Y - 27/5)/10\).
\end{proposition}

\begin{proof}
If the cap sum is at least \(2\), the input
\(u = (\mathrm{cap}_1(Y),\, 2 - \mathrm{cap}_1(Y))\) is nonnegative
(\(\mathrm{cap}_1(Y) \le 3/2 < 2\) on \(Y \ge 2\)), lies in
\(U\) (its coordinate sum is \(2\)), and respects the caps
(\(2 - \mathrm{cap}_1(Y) \le \mathrm{cap}_2(Y)\) is the cap-sum
inequality). Conversely any \(u \in U\) under the caps has
\(2 \le u_1 + u_2 \le \mathrm{cap}_1(Y) + \mathrm{cap}_2(Y)\), so a
cap sum below \(2\) is infeasible. The cap sum is the affine function
\(67/25 - (Y - 2)/5\), equal to \(2\) exactly at \(Y = 27/5\) and
declining beyond; the dual measure gives
\(\lambda^\top(\mathrm{cap} - u) = (\mathrm{cap\ sum})/2 -
(u_1 + u_2)/2 \le (\mathrm{cap\ sum})/2 - 1\), the stated margin.
\end{proof}

The derivation of the caps from a two-stock resource model, and the
fibre-width formula \(\Delta x_1(Y) = Y - 4\), are carried in the
dual-certificate companion (Abaee, 2026, Minimax dual certificates); this section owns the
per-aggregate audit and the radius reading.

\begin{table}[t]
\centering
\footnotesize
\caption{The benchmark's cap arithmetic on the audited aggregates
(exact rationals). The cap sum crosses the demand level \(2\) exactly
at \(Y^{*} = 27/5\); at \(Y = 6\) the cap sum falls \(3/25\) short of
\(2\), and the \(\ell_1\)-normalized dual measure \((1/2, 1/2)\) certifies the
obstruction with margin \(3/50\).}
\label{tab:bench}
\begin{tabular}{@{}l c c c l@{}}
\toprule
\(Y\) & \(\mathrm{cap}_1\) & \(\mathrm{cap}_2\) & sum & reading\\
\midrule
4 & $13/10$ & $49/50$ & $57/25$ & above\\
$9/2$ & $5/4$ & $93/100$ & $109/50$ & above\\
5 & $6/5$ & $22/25$ & $52/25$ & viable (witness $(6/5, 4/5)$)\\
$27/5$ & $29/25$ & $21/25$ & $2$ & critical\\
$11/2$ & $23/20$ & $83/100$ & $99/50$ & below\\
6 & $11/10$ & $39/50$ & $47/25$ & obstructed, margin $3/50$\\
\bottomrule
\end{tabular}
\end{table}

\begin{figure}[t]
\centering
\includegraphics[width=0.96\linewidth]{figs_ws4/fig_benchmark.pdf}
\caption{The benchmark. Both caps decrease linearly in the aggregate
\(Y\); their sum crosses the demand level \(2\) exactly at
\(Y^{*} = 27/5\). The \(Y = 5\) fibre is viable (green), the \(Y = 6\)
fibre obstructed with margin \(3/50\) (red).}
\label{fig:bench}
\end{figure}

\subsection{Static duality}\label{duality}

Static duality for the benchmark's two floor rows, abstracted to drift
rows on \([0,1]\): the max--min equals the
min-\(\lambda\) max at value \(-1/10\), with multiplier
\(\lambda = (1/2, 1/2)\) the \(\ell_1\)-normalized Farkas pair of the calculus's
worked certificate (Abaee, 2026, An obstruction calculus). Figure \ref{fig:duality} displays
the instance: under the \(\ell_1\)-normalized multiplier the expected drift is
constant \(-1/10\) on all of \([0,1]\). Convexity is load-bearing: with
\(U = \{0, 1\}\) and rows \((-1, +1)\), \((+1, -1)\), the common set is
empty, the max--min is \(-1\), and the min-multiplier maximum is
\(0\). Under convexity the witness is sparse:

\begin{proposition}[static minimax duality on the two-floor instance]
\label{prop:saddle}
On the obstructed instance with drift rows \(\psi_1(u) = 2/5 - u\),
\(\psi_2(u) = u - 3/5\) on \(U = [0,1]\),
\[
\max_{u \in [0,1]} \min_i \psi_i(u) \;=\;
\min_{\lambda \in \Delta_2} \max_{u \in [0,1]}
\bigl(\lambda_1 \psi_1 + \lambda_2 \psi_2\bigr) \;=\; -\,\frac{1}{10},
\]
with the unique saddle point \((u^{*}, \lambda^{*}) = \bigl(1/2,
(1/2, 1/2)\bigr)\); under \(\lambda^{*}\) the expected drift is
constant \(-1/10\) on all of \([0,1]\).
\end{proposition}

\begin{proof}
The min envelope \(\min_i \psi_i(u)\) is concave and piecewise
linear with its peak at the crossing \(2/5 - u = u - 3/5\), i.e.\
\(u^{*} = 1/2\), value \(-1/10\). For \(\lambda \in \Delta_2\)
the pooled drift \((2\lambda_1 - 1)\cdot(-u) + \text{const}\) is
affine with slope \(\lambda_2 - \lambda_1 = 1 - 2\lambda_1\); it is
constant exactly at \(\lambda^{*} = (1/2, 1/2)\), where the value is
\(\tfrac12(2/5) + \tfrac12(-3/5) = -1/10\) for every \(u\). At any
other \(\lambda\) the maximum over \([0,1]\) is attained at an
endpoint; writing \(A(\lambda)\) for the \(u = 0\) endpoint value, the
two endpoint values are \(A(\lambda)\) and \(A(1 - \lambda)\), and
the averaging bound
\(\max(A(\lambda), A(1-\lambda)) \ge \tfrac12\bigl(A(\lambda) +
A(1-\lambda)\bigr) = -\tfrac1{10}\) --- the average is exactly the
constant value, for every \(\lambda\) --- with equality only at
\(\lambda^{*}\): the convexity of the max, not symmetry alone, locates
the minimum. Weak duality (max-min \(\le\) min-max)
holds unconditionally; the exhibited pair attains both sides at
\(-1/10\), which is therefore the exact common value and the unique
saddle.
\end{proof}

\begin{proposition}[sparse obstruction witnesses]\label{prop:helly}
An infeasible finite family of affine constraints on a convex domain in
\(\mathbb{R}^m\) has an infeasible subfamily of at most \(m + 1\)
members, and \(m + 1\) is tight. The instance realizes the bound: the
three constraints \(\{u_1 \ge 1,\ u_2 \ge 1,\ u_1 + u_2 \le 1\}\)
are pairwise feasible, jointly infeasible, and no two of them certify
the obstruction.
\end{proposition}

\begin{proof}
Intersect each constraint's halfspace with the domain: an infeasible
family of compact convex sets in \(\mathbb{R}^m\) has an infeasible
subfamily of size at most \(m + 1\) (Helly's theorem), and \(m + 1\)
is tight: the reversed halfspaces \(\{u_i \ge 1\}\) and
\(\sum_i u_i \le 1\) in \(\mathbb{R}^m\) --- the facets of an
\(m\)-simplex, translated outward --- meet \(m\) at a time and never
all \(m+1\) (the facets themselves meet in the simplex, which is why
the reversal is needed). For the
instance, each pair is feasible --- at \((1,1)\), \((1,0)\), and
\((0,1)\) respectively --- while all three give the contradiction
\(2 \le 1\); the verification script verifies each pair at its printed
feasible point and the triple by the exact contradiction.
\end{proof}

The Helly-tight family of three constraints
\(\{u_1 \ge 1,\ u_2 \ge 1,\ u_1 + u_2 \le 1\}\) --- no two of them
already infeasible --- attains the obstruction with \(m{+}1 = 3\)
constraints. The feasible contrast at rows
\((2/5 - u,\ u - 1/5)\) carries value \(+1/10\) on both sides of the
saddle: a certificate exists exactly where the value is negative, and at
a zero value the instance is boundary-critical with no strictly positive
margin available.
Weak duality (max--min \(\le\) min--max) holds unconditionally. Strong duality holds where its hypotheses are verified: the control set compact \emph{and convex}, the rows affine (or concave) in \(u\), and the label family finite --- as in the static instance above and the finite-horizon static-observation classes of the calculus (Abaee, 2026, An obstruction calculus). Both hypotheses are load-bearing: the paper's own nonconvex example (\(U = \{0,1\}\), rows \((-1,+1)\), \((+1,-1)\)) is compact, finite, and row-affine, yet gaps (\(-1\) against \(0\)); and continuity of the rows does not repair convexity, as the nonconvex rows example above shows --- concavity of the rows in \(u\), not continuity, is what the minimax argument consumes.

\begin{figure}[t]
\centering
\includegraphics[width=0.96\linewidth]{figs_ws4/fig_duality.pdf}
\caption{Static duality on the two-floor instance. The two drift rows
cross at \(u = 1/2\); the max--min equals the min-multiplier maximum at
\(-1/10\), and the \(\ell_1\)-normalized Farkas pair \((1/2, 1/2)\) makes the
expected drift constant at the value.}
\label{fig:duality}
\end{figure}

\subsection{The certainty-equivalence drift audit}\label{drift}

For the index \(g(s) = s^{2}\) the uncorrected law's one-step drift is
\(g(s + 1/10) - g(s) = s/5 + 1/100\) on \([1,2]\): increasing, with
minimum \(21/100\) at \(s = 1\) and maximum \(41/100\) at \(s = 2\)
(Table \ref{tab:drift}); the corrected law's drift is identically zero --- by
construction, the corrected law being the lead removed. A biased
observation map misreads the state by a fixed one-step lead; on any
decision model conditioned on the map, that misreading is the mechanism
through which a nonempty kernel could empty, and the drift audit bounds
it exactly --- the remedy is structural correction of the map, not a
better sensor.

\begin{table}[t]
\centering
\footnotesize
\caption{The certainty-equivalence drift audit:
\(g(s + 1/10) - g(s) = s/5 + 1/100\) exactly; the corrected law's
drift is \(0\) at every probed state.}
\label{tab:drift}
\begin{tabular}{@{}l c c c@{}}
\toprule
\(s\) & uncorrected & corrected & increment per \(1/4\)\\
\midrule
1 & $21/100$ & 0 & \\
$5/4$ & $26/100$ & 0 & $1/20$\\
$3/2$ & $31/100$ & 0 & $1/20$\\
$7/4$ & $36/100$ & 0 & $1/20$\\
2 & $41/100$ & 0 & $1/20$\\
\bottomrule
\end{tabular}
\end{table}

\subsection{Multiple floors}\label{floors}

With several binding floors the Farkas certificate obstructs jointly: a
single certificate over all floors certifies nonviability where
per-floor certificates do not compose. On the audited benchmark's two
floors (Section \ref{benchmark}), the minimal infeasible subsystem
carries explicit witnesses --- the Helly-tight triple of Section
\ref{duality} --- and the corner \((1,1)\) is dynamically nonviable
as in the shared system. \emph{Quantifier and dual form.} Joint maintenance of several floors is the quantifier \(\exists u\ \forall j \in I(x)\) --- one control satisfying every active floor --- not \(\forall j\ \exists u_j\); per-floor certificates do not compose. The gap is visible in dimension one: on \(U = [-1, 1]\) the floors \(g_1(u) = u - \tfrac12 \ge 0\) and \(g_2(u) = -u - \tfrac12 \ge 0\) are each separately maintainable, yet no common control maintains both, and the \(\ell_1\)-normalized multiplier \(\lambda = (\tfrac12, \tfrac12)\) certifies the joint infeasibility with \(\sup_{u \in U}(\lambda_1 g_1 + \lambda_2 g_2)(u) = -\tfrac12 < 0\) --- a strict-separation dual certificate of exactly the Proposition~\ref{prop:helly} shape, with \(m + 1 = 2\) constraints necessary in dimension \(m = 1\) (each single floor is feasible). Scalar aggregation inherits the gap: a positive weighted average \(\sum_j \lambda_j h_j\) can conceal a violated component, and a multiplier inequality \(\sum_j \lambda_j \mathrm{D}h_j \ge 0\) does not imply that every active drift is nonnegative; the \(\ell_1\)-normalized dual measure of Section \ref{benchmark} certifies the aggregate shortfall exactly because both cap rows are charged jointly. In the audited convex static examples, infeasibility is certified by a multiplier supported on at most \(m + 1\) active constraints; no such bound is claimed for nonconvex controls or dynamic-policy images. Two further objects are recorded for subsequent study: a
first-breach timing bound for multi-floor systems, and a corner-cone
reading of the infeasibility region.

\subsection{Robustness margins}\label{margins}

Each certificate carries a margin, and each margin is a robustness
radius --- the size of the perturbation the certificate's verdict
tolerates.

\begin{proposition}[margins as radii]\label{prop:margins}
(i) \emph{Benchmark}: at \(Y = 6\) the obstruction survives any
per-cap increase \(\delta < 3/50\) of both caps and only there: the
certificate's margin \(3/50\) is exactly the per-cap robustness radius
(feasibility returns at \(\delta = 3/50\)). (ii) \emph{Timing}: the
boundary cells have zero slack (\(z_0 - 1 - T_{\mathrm{obs}} = 0\)):
any adversarial damage above unit rate flips them --- the inclusive
convention is exactly the statement that the audited identity holds on
zero-margin cells. (iii) \emph{Drift}: the corrected law's zero drift
has the sign of the correction error \(\varepsilon\) (\(2\varepsilon +
\varepsilon^2 = \varepsilon(2 + \varepsilon)\) at \(s = 1\)): the
sign-certificate is robust for every \(\varepsilon \in (-2, 0) \cup
(0, \infty)\). (iv) \emph{Decentralization}: each of the four forced
votes of the audited belief is knife-edge --- flipping that single vote
sustains nothing (zero margin, exact); the protocol's value is carried
by four coordinates, each of which is load-bearing.
\end{proposition}

\begin{proof}
(i) feasibility at \(Y = 6\) under uniform increases \(\delta\) is
\(\tfrac{47}{25} + 2\delta \ge 2 \iff \delta \ge
\tfrac{3}{50}\) --- exact arithmetic, checked at \(\delta =
\tfrac{3}{200}\) (infeasible) and \(\delta = \tfrac{3}{50}\)
(feasible). (ii) at \((z_0, T) = (2, 1)\) the worst damage under rate
\(1 + \varepsilon\) is \(1 + \varepsilon > 1 = z_0 - 1\). (iii) the
sign evaluation is the displayed factorization. (iv) for each forced
cell, the count of sustaining law pairs with that vote flipped is zero
(exhaustive over the \(256\) pairs; check P22).
\end{proof}

\subsection{Design rules}\label{rules}

The rules hold on the audited systems --- each an exact census identity or a proved general law, tagged by its evidence --- and transfer as default guidance only at that evidence level. \emph{Timing:} (proved general law, Proposition \ref{prop:timing}) reviews are boundary-inclusive ---
\(T_{\mathrm{obs}} \le \tau^{*}(z_0)\), the deadline at worst-case exit admissible under the audited endpoint convention (Section \ref{timing}). \emph{Coarseness:} (exact census, Section \ref{adequacy}) a fibre
merging states with disjoint safe-action sets obstructs at that
information state; the remedy is a fibre split, not a finer threshold
on the same index (Section \ref{adequacy}). \emph{Aggregation:} (exact census, Sections \ref{classes} and \ref{adequacy}) the aggregate index's fibres cross the safe boundary --- each mixes states of conflicting verdict --- and the aggregate kernel (\(26\) against the full \(28\)) certifies only a certainly-safe inner approximation; the certainly-safe set is the sound report.
\emph{Bias:} (exact audit, Section \ref{drift}) the uncorrected law's drift is at least \(21/100\) on a
101-point grid of \([1,2]\), while the corrected law's drift is
identically zero; a fixed bias is exactly correctable, and the remedy is
structural correction of the map, not a better sensor. The decentralized audit adds a final rule: (proved --- the restricted kernel is zero, Proposition \ref{prop:authority}) a whole-law authority restriction on one agent's vote collapses the sustainable set outright, so protocol
design precedes instrument tuning (Section \ref{decentralized}).

\subsection{Verification methods}\label{methods}

Every count, witness, identity, bound, figure, and tabulated entry in
this paper is regenerated by deterministic scripts in exact integer and
rational arithmetic, standard library only. The master verification
script recomputes the counts of Sections \ref{master}--\ref{regimes}
from the shared system's primitives --- the transition, the fibre
partition function, and the backward recursion with the no-repeat
protocol threaded as a parameter --- re-derives every tabulated entry of
Tables \ref{tab:master}, \ref{tab:laws}, \ref{tab:timing},
\ref{tab:census}, \ref{tab:bench}, and \ref{tab:drift} independently,
enumerates the census of Section \ref{adequacy}, evaluates the
benchmark identities of Section \ref{benchmark} in rational
arithmetic, asserts the drift bounds, and checks that each headline
number appears in the text. The figure-and-table generation script
recomputes and asserts each plotted or printed quantity from the same
primitives before drawing; the two scripts agree on every shared
quantity. The audited identities are additionally covered by the
per-system verification scripts of the companion papers (Abaee, 2026, An obstruction calculus;
see the Code availability statement). Two independent recomputation
paths guard the central recursions (checks P17--P18): an exhaustive
enumeration of all deterministic stationary policies --- \(2^{5} = 32\)
aggregate, \(2^{9} = 512\) full-information --- whose union of
safe-pair sets reproduces the kernels exactly (memoryless determinacy
on the instance, itself an exact verified fact), and a fixpoint
computation of each class's horizon family, whose stabilization equals
the audited kernel setwise. The recursion's cost is measured exactly:
the memoized backward recursion over the \(36\)-pair population uses
the distinct-subproblem counts printed by the script (states \(\times\)
horizon \(\times\) register per class), the policy enumerations use
\(32\) and \(512\) candidates, the decentralized census \(256\) law
pairs (\(64\) behavioral), the adequacy census \(15\) partitions, and
the two timing grids \(93\) cells each.

\subsection{Conclusion}\label{conclusion}

Exact audits bind the obstruction calculus to its worked systems: each
mechanism acquires a count, a witness, and a boundary. The timing
certificate fires everywhere beyond \(z_0 - 1\) and nowhere at the
boundary; the adequate monitoring designs on the audited target are
enumerable, and no coarsest one exists; policy-class value drops
unequally across aggregation and protocol restrictions, with the two
middle cells incomparable, the meet of their kernels institutional
setwise, and exactly one pair --- the alternation-forced pair --- beyond
both; unanimity costs a fixed and identified set of beliefs, sustained
by exactly the \(4 \times 4\) product of free votes, recoverable only
while both votes remain free; regime blindness loses exactly the
corner's pairs, while no frozen regime loses any; and the continuous
benchmark separates the viable from
the obstructed at \(Y^{*} = 27/5\) with rational certificates on both
sides. The incomparability of the middle cells is a shared-register
phenomenon: with per-branch registers the protocol cost collapses
(Proposition~\ref{prop:register}), which is precisely why the shared
register is the audited institutional object. The design rules --- boundary-inclusive reviews, fibre splits
over threshold tuning, certainly-safe reporting under aggregation,
structural repair of bias, and protocol design before instrument tuning
--- are the audit's operational content, and the master table is its
evidence: every count in the paper is one column of one table, every entry of which is regenerated by script. These examples do not constitute a general taxonomy of viability failures; they supply exact counterexamples, finite censuses, and reproducible benchmarks that separate mechanisms often conflated in qualitative discussion. Read through the master theorem, the audit is a small theory: the five mechanisms --- common-action failure, memory/register coupling, finite-time boundary exit, decentralized information loss, and joint convex infeasibility --- each carry a certificate, a monotonicity status under the theorem's three axes, a requirement-level characterization (statewise factoring for unanimity, incompatibility-graph adequacy for fibres, max-plus hitting times for regimes), and a robustness radius (\(3/50\) per cap on the benchmark, zero on the boundary cells and the forced votes, sign-robustness on the drift). Refinement completeness on the continuous side is the matching axis: the certification companion's meshes, stored scenarios, and belief cells tighten its sandwich exactly as the three axes here enlarge the kernels, with positivity non-monotone on both sides (Abaee, 2026, Computational viability certification)

\subsection*{References}

Abaee, A., 2026. Robust viability of the 2J3KL limit reference point
under a surplus-production map: policy scoring, expansion, and when
catch cannot help. Zenodo.
\url{https://doi.org/10.5281/zenodo.22552060}

Abaee, A., 2026. An obstruction calculus for viability under
incomplete observation (submitted).

Abaee, A., 2026. Computational viability certification under incomplete
observation: a continuous-to-finite bridge with exact certificates, a
solver-assisted certified campaign, and a finite-side reference library.
Manuscript submitted for publication.

Abaee, A., 2026. Minimax dual certificates for the common-action
obstruction. Manuscript submitted for publication.

Baccelli, F., Cohen, G., Olsder, G.J., Quadrat, J.-P., 1992.
Synchronization and Linearity: An Algebra for Discrete Event Systems.
Wiley, Chichester.

\subsection*{Declarations}

\subsection*{Funding}
None.

\subsection*{Competing interests}
None.

\subsection*{Data availability}
No external data were used.

\subsection*{Code availability}
The verification script \url{paper2_worked_systems_v17_verification.py}
and the figure-and-table generation script
\url{paper2_worked_systems_figures.py} are publicly available at
\url{https://github.com/MIKEAA2020/general-sustainability} (folder
\texttt{arena agent 1/paper rewrites/latex}); they reproduce all
counts, witnesses, bounds, figures, and tabulated entries verbatim.

\subsection*{AI declaration}
GLM (Z.ai), Qwen (Alibaba Cloud) and DeepSeek AI assisted with drafting
and iterative review. The author reviewed and edited outputs and takes
responsibility for the final work.

\end{document}
